\documentclass[../main.tex]{subfiles}
\begin{document}
%==========================================TIÊU ĐỀ TẬP========================================
\begin{table}[h]
    \begin{adjustwidth}{-1cm}{}
        \begin{tabular}{>{\centering\arraybackslash}p{6cm}|>{\raggedright\arraybackslash}p{12.5cm}}
            \multirow{5}{6cm}
            {\\[-40pt]
                \flaregothic\fontsize{20pt}{20pt}\selectfont \centering
                \color{eptype}ANNÉE\color{black}\\
                \freshbot\fontsize{72pt}{72pt}\selectfont
                \color{epnum}2022\\[6pt]
                \cafeteria\fontsize{20pt}{20pt}\selectfont\color{epnum}(C-B1)\color{black}
                \\[18pt]
                \color{black}
            }
             &
            % Dòng thứ nhất: tên tập tiếng Nhật
            {\fotskip\fontsize{18pt}{18pt}\selectfont ビーチでの\ruby{家}{か}\ruby{族}{ぞく}\ruby{旅}{りょ}\ruby{行}{こう}！！}
            \\[6pt]
            % Dòng thứ hai: tên tập tiếng Anh
             & {\cafeteria\fontsize{18pt}{18pt}\selectfont Beach Party-Rock, It's Time Four Summer Days!}
            \\[4pt]
            % Dòng thứ ba: tên tập tiếng Việt
             & {\vnmsans\fontsize{18pt}{18pt}\selectfont Buổi đi biển đáng nhớ}
            \\[8pt]
            % Dòng thứ tư: Nhân vật xuất hiện
             & {\caviar{\fontsize{11pt}{11pt}\selectfont Nhân vật: \emp{kuon2} \emp{kaiho2} \emp{suzuno2} \emp{ryuuhou2} \emp{naomi2} \emp{chloe} \emp{kson} \emp{kanata} \emp{luna} \emp{towa} \emp{watame}}}
            \\[8pt]
            % Dòng cuối cùng: Thời gian diễn ra của tập
             & {\sen\fontsize{16pt}{16pt}\selectfont Mốc thời gian: 10/07/2022 -- 12/07/2022}
            \\[18pt]
        \end{tabular}
    \end{adjustwidth}
\end{table}
%=========================================TRUYỆN CHÍNH========================================
\color{black}

\txtbx[2]{
    {\color{vol} \uline{Tóm tắt:}} Khép lại chiến dịch giải cứu và thanh tẩy dở khóc dở cười dành cho cô nàng cá voi sát thủ Sakamata Chloe, đại gia đình nhà Hikawa cùng các thành viên Thế hệ thứ Tư của \hll\ chính thức khởi động chuyến hành trình nghỉ mát ba ngày hai đêm tại bờ biển phía nam. Đối với Kuon, đây là chuyến đi biển gia đình thực thụ đầu tiên sau gần ba năm miệt mài chìm đắm trong guồng quay công việc quản lý. Dưới ánh đèn đêm nhộn nhịp của trạm tập kết kho ngoại ô và chuyến xe lăn bánh rạng sáng dưới bầu trời sao, những bước chân chập chững của Suzuno, Naomi và Chloe cùng mở ra những kỷ niệm đầu đời tràn ngập sắc màu hy vọng.
}

\stpt{-Jour 1-}
\stcap{Khởi hành, chuyến xe đêm và đại dương lần đầu hội ngộ}

\pov{Hikawa Kuon}

Mùi dầu máy thoang thoảng hòa cùng hơi gió nồng đượm của một đêm gần giữa tháng Bảy khẽ lướt qua sống mũi tôi.

Đứng trước khoảng sân bê tông rộng thênh thang của khu kho ngoại ô thuộc tổng công ty nơi mẹ tôi làm việc, tôi khẽ hít một hơi thật sâu rồi nhìn lên vòm trời đêm đầy sao lấp lánh. Đã lâu lắm rồi, tôi mới lại có cảm giác nhịp tim mình đập một cách thảnh thơi và bình thản đến nhường này.

Nếu có ai đó hỏi tôi rằng, lần gần nhất tôi cùng gia đình đi biển là khi nào, có lẽ tôi sẽ phải mất vài giây ngập ngừng lục lại trong ký ức.

Đó là một buổi trưa hè oi ả của tháng Tám năm 2019. Khi ấy, tôi vừa mới chân ướt chân ráo nhận chức Tổng quản lý tại văn phòng \hll. Văn phòng lúc đó còn thiếu thốn trăm bề, chiếc điều hòa cũ kỹ đình công khiến cả đám thần tượng quẫn bách đến mức tự bày trò đào cát, dựng bãi biển nhân tạo mini ngay giữa sàn nhà rồi vô tình làm vỡ toang cả ô cửa kính lớn. Đúng lúc tôi đang vò đầu bứt tai lo tiền đền bù thiệt hại thì một cuộc gọi bất ngờ từ mẹ Kaien báo tin công ty mẹ tổ chức xe đi biển và còn dư hẳn mười chỗ trống\footnote{{\flaregothic ÉPISODE 14}, quyển số Một thuộc \textit{holo no graffiti}.}. Chuyến đi biển năm ấy quả thực rất đáng nhớ, nhưng nói cho đúng, với tôi nó lại giống một chuyến công tác áp lực hơn là kỳ nghỉ. Suốt cả mấy ngày ở bờ biển, tôi phải cầm cuốn kịch bản chạy đôn chạy đáo ngoài nắng, vừa chỉ đạo đội hình, vừa canh góc quay cho phân cảnh vũ đạo chung của các cô gái trong đợt ghi hình MV \emph{``Shiny Smily Story''}.

Và rồi, suốt gần ba năm tiếp theo, tôi gần như ``mất tích'' khỏi tất cả các chuyến dã ngoại của gia đình.

Sự bùng nổ vũ bão của COVER với hàng loạt dự án dồn dập, sự ra đời của các nhánh quốc tế cùng khối lượng công việc hậu cần khổng lồ đã cuốn phăng tôi vào một guồng quay không hồi kết. Tôi từng là một kẻ cuồng công việc, luôn ám ảnh bởi tiền bạc và sự tích lũy phòng xa, đến mức tự biến mình thành một cỗ máy thức trắng đêm liên miên để xử lý khủng hoảng. Cho đến khi biến cố đa vũ trụ ập đến vào mùa đông năm ngoái, trải qua những giây phút sinh tử giữa ranh giới sống chết của The Void và thị trấn mịt mù sương ERROR, tôi mới bừng tỉnh nhận ra đâu là giá trị cốt lõi nhất của cuộc đời mình. Tiền bạc hay danh vọng suy cho cùng chẳng thể nào bù đắp được những khoảnh khắc ấm áp bên người thân yêu. Quyết định bàn giao bớt công việc vận hành trực tiếp để chuyển sang vai trò hướng dẫn, điều phối vào cuối năm 2021 không chỉ giúp tôi giải phóng sức lao động, mà quan trọng hơn cả, nó đã mang tôi trở về đúng nghĩa với mái ấm nhà Hikawa.

Và chuyến đi biển lần này chính là minh chứng rõ ràng nhất cho sự trở về ấy.

Mọi chuyện bắt đầu từ một chuỗi sự kiện đầy hỗn loạn cách đây đúng một tuần. Vụ việc cô nàng thực tập sinh cá voi sát thủ Sakamata Chloe bỏ bê việc tắm gội suốt một tháng trời đã biến thành một ``thảm họa sinh học'' cấp độ ba. Mẹ tôi và Marine đã phải đích thân mở một chiến dịch áp giải nghẹt thở, tống cổ con cá voi bướng bỉnh về thẳng ký túc xá Hikawa để thực hiện một cuộc tổng tẩy uế long trời lở đất. Thế nhưng, đằng sau lớp vỏ bọc ngổ ngáo và lười biếng ấy, chúng tôi mới vỡ lẽ ra Chloe thực chất chỉ là một cô gái nhút nhát, sợ bị phán xét và khao khát được sẻ chia. Cuộc hội ngộ bất ngờ với người bạn thanh mai trúc mã thuở ấu thơ ven làng chài mười sáu năm trước đã gắn kết cô nàng với gia đình tôi.

Để rồi trong bữa tiệc sườn nướng tối hôm ấy, kế hoạch về chuyến đi biển ba ngày hai đêm đã chính thức bùng nổ.

Kson đã dùng uy tín của mình để liên lạc và nhận được sự phê duyệt từ công ty cho bốn cô bạn Thế hệ thứ Tư gồm Kanata, Watame, Towa và Luna được nghỉ trọn vẹn ba ngày để xả hơi. Chloe cũng chẳng chịu kém cạnh khi giở tài lẻ ``tình báo'', soạn một bản báo cáo khẩn cấp gửi Tổng soái La+ Darknesss, xin phép được bám sát mục tiêu nguy hiểm mang tên Hikawa Kuon để thu thập cơ mật cho tổ chức holoX - một cái cớ vụng về đến mức chỉ mất ba giây để được sếp của cô nàng phê duyệt với điều kiện phải mua bánh mực nướng mang về làm quà.

Bố tôi cùng Delutaya tình nguyện ở lại ký túc xá để cáng đáng xưởng cơ khí, chăm nom bà nội đã cao tuổi và quan trọng nhất là làm chỗ dựa tinh thần, vỗ về cho Mikeneko ở tầng tám khi tâm trạng cô nàng vẫn còn nhiều nhạy cảm. Vì thế, lực lượng tháp tùng hướng về biển khơi lần này chỉ gồm mẹ làm trưởng đoàn kiêm tài xế, năm anh em chúng tôi, cô nàng cá voi Chloe và hội chị em Gen 4.

\begin{center}
    \uline{Mười một giờ đêm\\Khu kho bãi ngoại ô của Kaien.}
\end{center}

Khuôn viên kho hàng sáng rực dưới ánh đèn pha cao áp vàng vọt. Tiếng cười nói rôm rả, tiếng gọi nhau í ới của các chú bác công nhân hòa cùng tiếng kim loại va chạm leng keng tạo nên một bầu không khí nhộn nhịp hệt như một ngày hội nhỏ. Đoàn xe du lịch gồm ba chiếc xe khách loại vừa và hai chiếc xe gia đình bảy chỗ đang xếp hàng ngay ngắn bên mép sân, nắp cốp sau mở toang để mọi người chất đồ đạc.

Với tôi, Kaiho và Ryuuhou, nơi này chẳng có gì xa lạ. Thỉnh thoảng vào những dịp nghỉ lễ hoặc cuối tuần rảnh rỗi, ba đứa chúng tôi vẫn thường chạy xe xuống xưởng cơ khí của bố Ryujin gần đó, rồi tạt qua kho phụ mẹ bốc dỡ vài kiện hàng nhẹ hoặc tụ tập ăn kem cùng các bác thủ kho.

Vừa bước xuống xe, Ryuuhou đã nhảy chân sáo, giơ hai tay vẫy loạn xạ về phía một bác trung niên vạm vỡ đang vác thùng xốp lớn: ``Bác Toei ơi! Lâu quá không gặp bác! Năm nay bác vẫn mang theo cần câu cá biển tự chế của bác đấy chứ!?''

Người đàn ông quay lại, gương mặt sạm nắng nở một nụ cười sảng khoái: ``A, con bé con nhà Hikawa đấy à! Lớn tướng thế này rồi cơ đấy! Yên tâm đi, đợt này bác chuẩn bị cả chục bộ lưỡi câu xịn, ra ngoài bãi đá bác cháu mình tha hồ giật cá đối nhé!''

Kaiho đứng bên cạnh tôi, hai tay đút hờ vào túi áo khoác bomber tối màu. Dù ít khi mở miệng chuyện trò dài dòng, cô nàng vẫn lễ phép nghiêng đầu khẽ chào những người quen mặt, rồi không một động tác thừa, cô sải bước tiến thẳng đến chiếc xe chở đồ, cúi người nhấc bổng hai thùng đá lạnh nặng trĩu đặt gọn gàng vào góc cốp sau chỉ bằng một tay. Cánh đàn ông trong xưởng đứng gần đó thấy cảnh tượng ấy đều tròn xoe mắt thán phục: ``Ái chà chà! Con bé Kaiho này sức vóc còn khỏe hơn cả mấy tay bốc vác chuyên nghiệp của xưởng nữa! Khá lắm cháu gái!''

Kaiho chỉ nhàn nhạt đáp lại bằng một cái gật đầu cụt lủn: ``Tiện tay thôi ạ.''

Thế nhưng, trái ngược hoàn toàn với vẻ dạn dĩ của ba chúng tôi, hai bóng hình nhỏ bé phía sau lại đang bối rối thu mình lại.

Suzuno và Naomi là lần đầu tiên đặt chân đến một khu kho xưởng công nghiệp rộng lớn đến thế này. Suzuno mặc một chiếc áo len mỏng màu be bên ngoài váy dài trang nhã, mái tóc xõa ngang vai khẽ lay động theo làn gió đêm. Cô nàng mở to đôi mắt dịu dàng, tò mò quan sát những chiếc xe nâng hàng đang chạy chầm chậm và những kiện gỗ xếp cao ngút ngàn, đôi bàn tay đan nhẹ vào nhau trước bụng với vẻ hơi dè dặt.

Còn Naomi thì nép sát hẳn vào bên cạnh Suzuno, hai bàn tay nhỏ xíu níu chặt lấy vạt áo chị gái, chiếc nón cói rộng vành đội hơi sụp xuống che bớt nửa khuôn mặt non nớt. Đôi mắt to tròn trong veo của cô bé chớp chớp liên tục, vừa sợ sệt vừa kinh ngạc nhìn những người đàn ông cao to, bắp tay cuồn cuộn đang khuân vác rầm rập quanh mình.

Mẹ tôi lúc này đang cầm cuốn sổ kiểm tra danh sách hành lý bước đến, vỗ vai hai đứa nhỏ rồi cất giọng sang sảng giới thiệu với các đồng nghiệp trong xưởng: ``Này các ông các bác ơi! Giới thiệu với mọi người, đây là hai đứa em gái mới của thằng Kuon nhà chị đấy! Con bé lớn là Suzuno, còn đứa nhỏ này là Naomi. Lần đầu tiên chúng nó được đi xa cùng với công ty mình, mọi người có trêu chọc thì cũng nhẹ tay thôi kẻo hai đứa nó sợ nhé!''

Một tràng trầm trồ đầy ngạc nhiên lập tức rộ lên từ các công nhân chung quanh: ``Ối chà! Nhà bà Kaien khéo nuôi con thật đấy! Thằng Kuon thì nhanh nhẹn, con bé Kaiho thì mạnh mẽ, Ryuuhou lại dẻo mép, giờ lại có thêm hai cô bé xinh xắn, ngoan ngoãn như búp bê thế này!''

Một bác bảo vệ đứng tuổi, hiền từ bước lại gần, cúi người xuống ngang tầm mắt Naomi rồi rút từ túi áo ra một thanh kẹo sô-cô-la hạt dẻ đưa về phía cô bé: ``Bé con ơi, cầm lấy ăn cho đỡ say xe này. Đi chơi biển với mọi người thì phải cười thật tươi lên chứ, việc gì mà phải sợ bác?''

Naomi tròn xoe mắt, rụt rè nhìn thanh kẹo rồi ngước lên nhìn Suzuno như để xin phép. Cô chị mỉm cười ấm áp, nhẹ nhàng xoa đầu cô bé: ``Nhận đi em, bác cho đấy.''

Được chị cho phép, cô em út mới dám chìa hai bàn tay nhỏ xíu ra đón lấy thanh kẹo, rồi cúi gập người thật sâu, cất giọng thỏ thẻ, lí nhí như tiếng mèo con: ``Cháu\dots\ cháu cảm ơn bác ạ.''

Cử chỉ lễ phép đến mức đáng yêu của cô bé khiến cả đám người lớn xung quanh bật cười ồ lên thích thú, kẻ xoa đầu, người trầm trồ khen ngợi không ngớt.

Cách đó một đoạn, một chú kỹ thuật viên máy móc hay làm việc chung kíp với mẹ tôi bước sang hỏi thăm Suzuno: ``Nghe bà Kaien khoe con bé Suzuno học giỏi lắm, cháu đang theo học ngành gì thế?''

Suzuno hơi đỏ mặt, khẽ vén lọn tóc mai ra sau vành tai rồi từ tốn đáp: ``Dạ\dots\ cháu theo đuổi mảng kỹ thuật hóa học và điều chế dược liệu ạ.''

Người chú nghe vậy liền trố mắt, gật gù liên hồi: ``Kỹ thuật hóa học á!? Giỏi thế! Con gái mà học ngành kỹ thuật khó nhằn đấy là hiếm lắm nha! Đã học giỏi lại còn thùy mị, xinh xắn thế này, sau này khối chàng trai xếp hàng dài từ Kawachi ra tận biển để xin cưới cho mà xem!''

Mặt Suzuno lập tức đỏ ửng như một quả cà chua chín, hai vành tai nóng ran khiến cô nàng chỉ biết bối rối cúi gằm mặt xuống. Đứng ngay bên cạnh, Ryuuhou được đà khoái chí, huých nhẹ khuỷu tay vào hông chị gái trêu chọc: ``Nghe chưa, chị Suzu! Có khi đợt này ra biển lại khối anh dính tiếng sét ái tình với chuyên gia hóa học nhà mình ấy chứ lị!''

``Ryuuhou này\dots!'' Suzuno ngượng ngùng khẽ gắt một tiếng.

Chưa kịp để cô em thứ tư lấn lướt trêu chọc tiếp, một bóng người cao gầy đã bất thình lình bước chéo lên, đứng chắn ngay trước mặt Suzuno. Kaiho khoanh hai tay trước ngực, đôi mắt sắc lạnh quét qua người chú kỹ thuật viên rồi liếc nhìn Ryuuhou, cất giọng trầm thấp, đều đều nhưng mang sức răn đe tuyệt đối: ``Em ấy chưa có ý định yêu đương gì đâu. Mấy trò tán tỉnh vớ vẩn ngoài bãi biển thì dẹp hết đi. Đứa nào dám bén mảng lại gần làm phiền Suzu, tôi bẻ khớp từng đứa một.''

Lời tuyên bố đanh thép, lạnh như băng kèm theo cái bẻ khớp ngón tay \textit{rắc rắc} rợn người của Kaiho làm người chú kia rụt cổ cười trừ: ``Ấy chết, bác chỉ đùa một tí cho vui thôi mà cháu gái! Khiếp, có cô chị gái bảo kê gắt gao thế này thì bố bảo đứa nào dám léng phéng!''

Suzuno khẽ thở phào, đưa tay kéo nhẹ vạt áo sau lưng Kaiho, mỉm cười dịu dàng: ``Chị đừng dọa bác ấy thế chứ, Onee-sama\dots\ Dù sao thì, cảm ơn chị nhé.'' Kaiho chỉ khẽ hừ một tiếng trong cổ họng rồi đứng sát cạnh cô em gái, như một bức tường thép bảo vệ Suzuno khỏi mọi ánh nhìn soi mói.

Ở một góc khuất cạnh bánh sau chiếc xe bảy chỗ, Chloe đang đứng lặng lẽ một mình. Cô mặc một chiếc áo phông thụng dài trùm qua hông và chiếc quần soóc bò đơn giản, tay cầm một que kem vani cắn dở. Đôi mắt đỏ rực của cô khẽ hướng về phía đám đông đang vây quanh Suzuno và Naomi. Nhìn nụ cười hồn hậu của những người lao động chân chất và sự che chở ấm áp mà gia đình Hikawa dành cho nhau, ánh mắt Chloe thoáng chùng xuống, đượm một nỗi niềm xa xăm khó tả. Có lẽ hình ảnh ấy đã vô tình chạm vào những khoảng trống cô độc trong suốt những năm tháng cô phải tự mình bươn chải nơi đô thị sầm uất. Thế nhưng chỉ một thoáng sau, nhận thấy tôi đang liếc nhìn về phía mình, cô nàng liền giật thót, vội vàng quay ngoắt mặt sang hướng khác, giả bộ chăm chú vuốt màn hình điện thoại tra cứu bản đồ tuyến đường.

Cách đó một đoạn, cạnh chiếc xe bảy chỗ thứ hai, bốn bóng hình nhỏ nhắn trùm kín áo khoác chống nắng và mũ lưỡi trai đã đến kho từ trước đang chụm đầu xì xào bàn tán.

Kanata hai tay đang ôm chặt một bọc đồ ăn nhẹ, đôi mắt sáng rực háo hức: ``Oa, đông vui náo nhiệt thật đấy! Cảm giác như cả một binh đoàn chuẩn bị ra trận vậy! Mấy thùng xốp đằng kia chắc chắn là chứa đầy đồ ăn ngon rồi!''

Watame đứng bên cạnh, kéo thấp vành mũ che đi đôi sừng cừu đặc trưng, mũi khẽ hít hà không khí rồi xuýt xoa: ``Mùi gió đêm mát quá đi mất\dots\ Em nghe Kson-chan bảo ra tới biển sẽ có tiệc nướng hải sản tươi sống nữa đó! Tôm nướng, mực nướng, sò điệp nướng mỡ hành\dots\ Nghĩ thôi là bụng em reo ầm ĩ rồi này!''

Towa khoanh hai tay, thở dài bất lực nhìn hai cô bạn háu ăn nhưng khóe môi vẫn khẽ cong lên: ``Hai người bớt mơ mộng chuyện ăn uống lại giùm tôi cái! Chúng ta đang ở chỗ đông người đấy, nhớ che chắn cho kỹ vào. Nhỡ bị người ta nhận ra rồi bu lại xin chữ ký hay chụp hình thì lộ hết cả chuyến đi bí mật của công ty bây giờ!''

Ngồi vắt vẻo trên chiếc vali kéo cạnh đó, Luna đang ngậm kẹo mút, cất giọng nũng nịu ngái ngủ: ``Ưm\dots\ Luna bùn ngủ quá-nanora\dots\ Xe vẫn chưa chại sao-nora? Đến nơi thì nhớ gọi dậy để còn đi bơi đó nha\dots''

Và rồi, sự tĩnh lặng hiếm hoi ấy lập tức bị xé toạc bởi một tràng cười oang oang từ phía cổng kho.

Kson - khoác trên mình chiếc áo sơ mi hoa hòe đậm chất du lịch Hawaii mở phanh cúc ngực khoe lớp áo ba lỗ trắng bên trong - hiên ngang bước vào. Cô nàng vỗ bốp một phát thật mạnh vào vai một bác tài xế lớn tuổi như thể hai người là bạn nối khố từ kiếp trước: ``Chào buổi tối các chiến hữu! Tinh thần anh em thế nào rồi!? Để tôi nói cho các ông nghe, chuyến đi này mà thiếu tôi là bớt vui một nửa đấy nhé, khà khà khà!''

Bác tài xế trố mắt nhìn người phụ nữ tóc vàng cao ráo với phong thái giang hồ bạt mạng: ``Ối mẹ ơi, cô em Tây này ở đâu ra mà bốc lửa thế!? Bạn của thằng Kuon đấy à?''

Kson khoái chí ngoác miệng cười lớn, giơ ngón tay cái chỉ ngược vào tôi: ``Mấy ông mà biết con bé này với đám đàn em đi sau tôi là ai thì có mà ngã ngửa cả lũ ra đất--''

Tôi hoảng hồn lao tới, dùng cả hai tay bịt chặt lấy mồm cô nàng trước khi cái mồm bô bô của nàng cựu Tổng trưởng kịp tiết lộ danh tính của bốn cô nàng Thế hệ thứ Tư: ``Im ngay giùm anh cái, Kson! Em muốn cả khu kho bãi này náo loạn lên vì nhận ra idol hay sao!?''

Kson ú ớ trong tay tôi một hồi rồi mới giãy ra, phủi phủi áo rồi cười sặc sụa: ``Rồi rồi, biết rồi! Bí mật cơ mà, làm gì mà căng thế, paisen! Với cả, ở cái xứ Kawachi này thì người ta chỉ lo làm ăn buôn bán, có mấy ai biết \hll\ là cái đinh gì đâu chứ!''

Tôi trừng mắt, tiến sát mặt cô nàng hạ giọng nhắc nhở: ``Thà cẩn thận còn hơn em ạ. Ai dám chắc trong cái kho hàng cả trăm con người này không có ai lướt mạng xem stream hay biết đến mấy đứa? Nhỡ có một người nhận ra rồi quay clip tung lên mạng thì cả đám khỏi có kỳ nghỉ yên ổn luôn đấy. Đừng có biến mình thành cái loa phường như thế!''

Kson nhún vai, giơ hai tay đầu hàng: ``Tuân lệnh, sếp tổng! Miệng em khóa van một trăm phần trăm rồi đây!''

\ct

Sau gần một tiếng đồng hồ điểm danh, xếp đồ và chuẩn bị chu đáo, đồng hồ trên bảng điều khiển xe chính thức điểm mười hai giờ ba mươi phút sáng.

Đoàn xe khởi hành.

Mẹ tôi đích thân cầm lái chiếc xe bảy chỗ dẫn đầu của gia đình, tôi ngồi ở ghế phụ phía trước để hỗ trợ định vị và quan sát đường sá. Hàng ghế giữa là Suzuno, Naomi và Chloe, trong khi Kaiho và Ryuuhou ngồi ở hàng ghế sau cùng bên cạnh những chiếc ba-lô lớn. Chiếc xe bảy chỗ thứ hai chở nhóm bốn thành viên Gen 4 do Kson cầm lái chạy bám sát ngay phía sau, tiếp theo là đoàn xe khách của các cán bộ công nhân viên trong kho. Cậu tôi - em trai của mẹ kiêm luôn sếp của kho - cùng với bà ngoại Akiko ngồi ở một trong những chiếc xe đó.

Khi chiếc xe lăn bánh qua cánh cổng sắt của khu kho, ánh đèn đường vàng vọt của thành phố cứ thưa thớt dần rồi nhường chỗ cho bóng tối bao la của tuyến đường cao tốc liên tỉnh.

Naomi ngồi ở giữa, một tay níu vạt áo Suzuno, mặt dán sát vào khung cửa kính trong suốt. Đôi mắt to tròn của cô bé sáng long lanh, phản chiếu những vệt sáng lướt qua ngoài ô cửa. Đây là lần đầu tiên trong suốt cuộc đời mình, cô bé được ngồi trên một chiếc xe chạy xuyên qua màn đêm, hướng về một nơi mà trước đây em chỉ thấy qua những trang sách cũ và màn hình vô tuyến.

Hành trình tiến về đại dương bao la đã chính thức bắt đầu.

\ct

Một giờ ba mươi phút sáng.

Không gian trong khoang xe bảy chỗ chìm vào tĩnh mịch, chỉ còn lại tiếng động cơ gầm gừ đều đặn cùng tiếng gió rít nhè nhẹ bên ngoài mặt kính. Chiếc xe lướt đi êm ru trên dải nhựa đường hun hút, bỏ lại phía sau những ngọn đèn cao áp lẻ loi của vùng ngoại ô.

Ở hàng ghế sau cùng, Ryuuhou đang phải vật lộn với cơn buồn ngủ dữ dội. Đầu cô em gái thứ tư cứ gật gà gật gù như gà mổ thóc, chốc chốc lại giật mình chớp mắt nhìn quanh rồi lại từ từ rũ xuống. Kaiho ngồi cạnh bên cửa sổ, hai tay khoanh trước ngực, ánh mắt bình thản nhìn ra dải bóng tối mịt mùng bên ngoài. Thấy Ryuuhou sắp chúi đầu vào mép cửa xe, không nói một lời, cô nàng chỉ vươn tay chỉnh lại chiếc gối tựa đầu hình chữ U rồi đẩy nhẹ trán em gái ngả vào đó. Đôi mắt Ryuuhou nhắm nghiền lại, miệng lẩm bẩm vài câu vô nghĩa về mấy khuông nhạc rồi thở đều đều chìm sâu vào giấc ngủ.

Ở hàng ghế giữa, Naomi cũng đã thiếp đi từ rất sớm. Bàn tay nhỏ xíu của cô em út vẫn nắm chặt góc áo của Suzuno, đầu ngả hẳn vào vai chị. Suzuno cẩn thận kéo chiếc chăn mỏng đắp ngang ngực Naomi, khẽ đưa tay vuốt lại vài sợi tóc tơ trên trán cô bé với ánh mắt dịu dàng như một người mẹ trẻ.

Thế nhưng, thế cờ thăng bằng ấy bỗng bị phá vỡ bởi kẻ ngồi phía bên kia.

Chloe - người bị kẹp giữa Suzuno và cánh cửa bên phải - ban đầu còn cố giương mắt nhìn màn hình điện thoại. Nhưng vì sợ ánh sáng xanh hắt vào làm Naomi thức giấc, cô nàng đành tiu nghỉu tắt máy. Chiếc áo phông rộng thùng thình và cơn gió điều hòa mát lạnh nhanh chóng đánh gục sức đề kháng của con cá voi sát thủ vốn nổi tiếng lười vận động. Ban đầu Chloe chỉ khẽ ngáp một cái thật dài, mi mắt trĩu nặng dần, rồi chẳng biết từ lúc nào, cả cái đầu tóc bù xù của cô nàng rũ sang một bên, đổ ập vào bờ vai còn lại của Suzuno.

Suzuno giật mình khẽ cứng người. Cô bé nhìn sang bên trái là Naomi đang say giấc nồng, nhìn sang bên phải lại là Chloe đang ngáy nhè nhẹ, cả hai bên vai đều trở thành điểm tựa vững chãi cho hai ``hành khách'' say ngủ.

Tôi ngồi ở ghế phụ phía trước, khẽ liếc qua gương chiếu hậu trong xe rồi xoay nửa người lại, nhịn không được bèn thì thầm trêu: ``Sao rồi Suzuno? Em bị hai bên kẹp chặt như miếng sandwich thịt nguội rồi à?''

Suzuno ngước đôi mắt dị sắc lên nhìn tôi, gương mặt thoáng nét dở khóc dở cười vô cùng đáng yêu. Cô bé giữ nguyên tư thế thẳng đơ, thậm chí không dám thở mạnh: ``Vâng ạ\dots\ Cả hai người đều đang ngủ say quá, em cử động mạnh một chút là sợ họ giật mình thức giấc mất.''

Từ hàng ghế sau, Kaiho vẫn không buồn quay đầu lại, mắt dán chặt vào bóng đêm ngoài cửa sổ nhưng khóe môi khẽ nhếch lên, cất giọng trầm thấp, đều đều: ``Kệ chúng nó đi, Suzu. Cứ thả lỏng lưng dựa vào ghế mà ngủ. Đứa nào trượt xuống tự khắc sẽ biết tìm tư thế khác để ngủ tiếp thôi.''

Suzuno khẽ mỉm cười, giọng dịu hẳn xuống: ``Thôi mà Onee-sama, để hai đứa nhỏ ngủ thêm chút nữa cũng được ạ\dots''

Mẹ tôi ngồi sau vô lăng khẽ bật cười khúc khích, mắt vẫn tập trung cao độ vào vệt đèn pha rọi sáng mặt đường phía trước: ``Mấy đứa cứ tranh thủ chợp mắt đi, để mẹ nhìn bản đồ cho. Còn khoảng hơn một tiếng nữa là tới trạm dừng nghỉ liên tỉnh đấy. Xuống đó tha hồ mà vươn vai vận động!''

\ct

Hai giờ bốn mươi lăm phút sáng.

Ánh đèn neon màu vàng cam rực rỡ của trạm dừng chân quy mô lớn ven đường cao tốc hiện ra như một ốc đảo giữa đêm đen. Hàng chục chiếc xe tải đường dài và xe khách chạy đêm đang đỗ rải rác khắp bãi đỗ rộng lớn. Đoàn xe của công ty mẹ lần lượt tấp vào khoảng sân trống trước sảnh chính.

Cửa xe vừa mở toang, một luồng không khí đêm vùng cao lành lạnh tràn vào khoang lái khiến ai nấy đều rùng mình tỉnh táo hẳn.

Ryuuhou vừa mới nghe tiếng còi xe dừng lại đã lập tức bật dậy như một chiếc lò xo, hai mắt sáng quắc như chưa từng biết buồn ngủ là gì: ``Aaa! Trạm dừng nghỉ rồi! Có quầy đồ ăn đêm không!? Đi ăn mì xíu mại đi mọi người ơi, bụng em bắt đầu biểu tình rồi này!''

Trong khi đó, Naomi được Suzuno khẽ lay vai đánh thức. Cô em út ngồi dậy với mái tóc rối bù xù, hai mắt kèm nhèm chớp liên tục, đôi bàn tay dụi dụi mắt rồi ngước nhìn quanh bằng vẻ ngơ ngác tột độ: ``Suzuno-oneechan\dots\ Mình\dots\ mình tới biển rồi ạ\dots?''

Suzuno bật cười hiền từ, rút chiếc lược nhỏ trong túi xách ra nhẹ nhàng chải lại tóc cho em gái: ``Chưa đâu Naomi-chan. Xe mới đi được một phần ba quãng đường thôi, chúng ta dừng lại nghỉ chân và uống nước ấm nhé.''

Chloe lảo đảo bước xuống xe sau cùng, hai mắt lờ đờ như người mất hồn, chiếc áo phông xộc xệch kéo lệch một bên vai, vừa đi vừa lầm bầm: ``Oaa\dots\ lạnh quá\dots\ Sao ngoài này gió to thế không biết, poe\~{}\dots''

Cách đó một đoạn, chiếc xe bảy chỗ thứ hai cũng vừa tắt máy. Cửa xe bật mở, Kson vươn vai một cái thật dài, xương khớp kêu răng rắc. Theo sau cô nàng là bốn bóng hình nhỏ nhắn của Thế hệ thứ Tư, ai nấy đều quấn khăn hoặc trùm kín áo khoác để giữ ấm.

Kson quay lại nhìn bốn cô em, vung tay ra hiệu đầy hào sảng: ``Mấy đứa xuống xe đi! Tôi đưa mọi người vào quầy kiếm cái gì nóng nóng lót dạ. Ở đây có súp bò với bánh bao hấp ngon đỉnh chóp đấy, cứ ăn thoải mái, tôi bao hết!''

Kanata kéo cao cổ áo len, hai mắt sáng rực nhìn về phía các quầy hàng bốc khói nghi ngút: ``Tuyệt vời! Tôi muốn một tô mì bò hầm cay xè lưỡi! Ngồi xe nãy giờ chân tay tê hết cả rồi!''

Watame xoa xoa hai bên sừng cừu ẩn dưới mũ trùm đầu, miệng nuốt nước bọt: ``Bánh bao hấp nóng hổi nữa nha Kson-chan! Thơm quá đi mất, tôi ngửi thấy mùi thịt nướng từ tít đằng kia kìa!''

Towa thì thở dài, bước đến bên cạnh Luna đang bước đi loạng choạng như chim cánh cụt con: ``Watame, ăn vừa thôi kẻo sáng mai ra biển bụng tròn vo như quả bóng đấy! Này Luna, đi cẩn thận kẻo vấp vào bậc tam cấp bây giờ!''

Luna ngáp ngắn ngáp dài, hai tay níu chặt lấy vạt áo Towa, giọng ngái ngủ nũng nịu: ``Hưm\dots\ Lạnh quá nanora\dots\ Muốn uống sữa đậu nành nóng cơ nanora\dots''

Kson cười lớn, vỗ lưng Towa một cái bộp: ``Yên tâm đi Towa-sama! Tôi canh chừng cho, không để con bé nào bị lạc đâu mà sợ! Đi nào các quý cô!''

Nhìn bóng lưng ríu rít của nhóm Thế hệ thứ Tư kéo nhau vào khu ẩm thực, tôi không khỏi khẽ mỉm cười. Tôi bước sang máy bán nước tự động đặt ở góc hiên ngoài trời, thả vài đồng xu vào khe rồi bấm chọn một lon cà phê đen không đường.

Ngay bên cạnh, một bóng người cao gầy quen thuộc đã đứng tựa lưng vào lan can sắt từ lúc nào.

Kaiho cầm trên tay một lon cà phê sữa đóng hộp, ánh mắt chăm chú nhìn về phía dải đen thăm thẳm của cánh rừng cao su trải dài tít tắp phía bên kia đường cao tốc. Gió đêm lồng lộng thổi tung những lọn tóc ngắn của cô, chiếc nơ cầu vồng cài buộc trên tóc khẽ rung rinh.

Tôi mở nắp lon cà phê, nhấp một ngụm đắng ngắt rồi đứng dựa vào lan can bên cạnh cô em song sinh: ``Không vào trong ăn chút gì cho ấm bụng à, Kaiho?''

Kaiho không quay đầu lại, ngón tay cái miết nhẹ lên thành lon nhôm lạnh ngắt. Giọng cô nàng vang lên trầm thấp, bình thản nhưng lại mang theo một thoáng trầm tư hiếm thấy: ``Không đói. Uống lon này là đủ tỉnh táo rồi.''

Khoảng lặng giữa hai anh em kéo dài vài giây, chỉ có tiếng gió rít qua rặng cây và tiếng cười nói từ khu nhà ăn vọng lại. Đột nhiên, Kaiho lên tiếng, phá vỡ bầu không khí tĩnh mịch: ``Lần trước mình đi xa khỏi thành phố Kawachi như thế này là khi nào nhỉ, Aniki?''

Tôi hơi khựng lại, khẽ nhíu mày lục lọi trong trí nhớ: ``Hồi hai tháng trước\dots\ lúc đưa bà nội xuống bệnh viện dưới ngoại ô để mổ mắt. Anh với em phải thay nhau túc trực suốt ba ngày đêm ở xưởng của bố để chạy xe đưa đón rồi.''

Kaiho khẽ lắc đầu, khóe môi nhếch lên một nụ cười nhàn nhạt: ``Không tính lần đó. Ý em là\dots\ một chuyến đi chơi thực sự ấy.''

Câu hỏi ấy khiến ngực tôi khẽ thắt lại. Một chuyến đi chơi thực sự. Một kỳ nghỉ đúng nghĩa mà không vướng bận công việc, không phải lo toan bệnh tật, cũng chẳng phải chạy trốn những hiểm họa đa vũ trụ ma quái.

``Chắc là\dots\ chưa từng có,'' tôi thở dài, hơi ấm từ hơi thở phả ra tan biến vào màn đêm. ``Từ ngày em với Suzuno và Naomi về với gia đình, anh toàn bị cuốn vào mấy mớ hỗn độn ở công ty với những chuyến du hành bất đắc dĩ. Đây là lần đầu tiên cả nhà mình thực sự bước lên một chiếc xe chỉ để đi ngắm biển.''

Kaiho im lặng một lúc lâu. Đôi mắt màu xanh lam sắc sảo của cô khẽ dịu lại, hướng về bóng tối vô tận phía chân trời xa: ``Ừm. Hồi bố mẹ em còn sống\dots\ trước khi xảy ra vụ tai nạn ấy\dots\ bố cũng từng hứa vào mùa hè sẽ đưa hai chị em ra biển chơi một chuyến. Nhưng rồi chuyện không may xảy đến, cái căn chung cư cũ dột nát kia đã chôn vùi cả tuổi thơ của Suzu. Em đã nghĩ\dots\ lời hứa đi biển đó sẽ mãi mãi là một giấc mơ không bao giờ thành hiện thực.''

Cô nàng dừng lại, bàn tay đang cầm lon cà phê khẽ siết chặt cho đến khi lớp vỏ nhôm phát ra tiếng rắc nhẹ, lõm hẳn một vết: ``Nhưng rốt cuộc, hôm nay em cũng đưa được Suzu đi thật rồi.''

Tôi nhìn sang cô em gái song sinh, trong lòng dâng lên một niềm cảm xúc khó tả. Đằng sau sự gai góc, lạnh lùng và những cú đấm thép có thể bẻ gãy mọi chướng ngại vật, Kaiho vẫn luôn là người gánh trên vai những vết thương lòng sâu hoắm nhất của quá khứ, và cũng là người yêu thương, chở che cho Suzuno hơn bất kỳ ai trên đời.

Tôi đưa lon cà phê của mình sang, khẽ cụng nhẹ vào lon của cô nàng một tiếng \textit{cách}: ``Yên tâm đi. Lần này chúng ta không chỉ đi để ngắm biển cho biết đâu. Chúng ta sẽ cùng nhau tận hưởng trọn vẹn từng con sóng, ăn sạch hải sản ở bãi biển, và tạo ra những kỷ niệm mà không một biến cố nào có thể xóa nhòa được.''

Kaiho liếc nhìn tôi, ánh mắt lóe lên vẻ bất ngờ rồi nhanh chóng trở lại vẻ lạnh lùng thường ngày, nhưng khóe môi cô đã cong lên rõ rệt: ``Biết thế thì liệu mà lo liệu cho chu đáo. Anh mà để Suzu hay Naomi bị say nắng hay đói bụng, coi chừng em đấm thủng tường phòng anh đấy, Aniki.''

Cách đó vài bước chân, Chloe với tay đang cầm nửa chiếc bánh mì kẹp xúc xích vừa mua từ quầy vô tình nghe trọn vẹn cuộc trò chuyện của hai anh em. Bước chân cô nàng khẽ khựng lại giữa sân. Đôi mắt đỏ rực nhìn chăm chú vào bóng lưng của hai người, nơi ánh đèn vàng hắt xuống tạo nên những vệt bóng dài đan xen vào nhau. Một cảm giác ấm áp lẫn chút ghen tị trẻ con khẽ dâng lên trong lòng cô nàng sát thủ đơn độc.

Thế nhưng, khoảnh khắc xúc động ấy lập tức bị dập tắt không thương tiếc bởi tiếng bước chân huỳnh huỵch từ phía sau.

Kson hai tay xách hai túi ni-lông to đùng đựng đầy bánh bao nóng và nước ngọt lon, hùng hổ lao tới đập bốp một phát vào lưng tôi khiến tôi suýt sặc ngụm cà phê trong miệng: ``Này hai anh em nhà kia! Đứng đây tự kỷ tâm sự tuổi hồng thơ ngây đấy à!? Em mua thêm chục cái bánh bao nhân thịt trứng cút với sữa đậu nành nóng cho cả nhà đây! Vào chia nhau ăn nhanh lên kẻo nguội ngắt bây giờ!''

Bốn cô nàng Gen 4 cũng lục tục kéo ra theo sau, miệng nhai nhóp nhép đầy thỏa mãn. Kanata giơ hai ngón tay làm dấu chữ V cười toe toét, trong khi Luna đã cầm trên tay hộp sữa đậu nành nóng hổi, vẻ mặt vô cùng hạnh phúc.

Bầu không khí lắng đọng lập tức vỡ òa trong tiếng cười nói ồn ào. Kaiho thu lại ánh mắt, nhanh như chớp giật lấy một chiếc bánh bao nóng từ tay Kson: ``Cảm ơn. Vừa lúc chị đang đói.''

Nghỉ ngơi và nạp lại năng lượng trong khoảng ba mươi phút, đoàn xe lại một lần nữa nổ máy, tiếp tục hành trình hướng về phương nam.

\ct

Bốn giờ ba mươi phút sáng.

Màn đêm đen đặc quánh dần chuyển sang màu xanh thẫm, rồi từ từ chuyển thành sắc tím xám nhạt nơi đường chân trời xa xôi.

Naomi là người đầu tiên cựa mình tỉnh giấc. Cô em út khẽ chớp mắt, áp sát trán vào mặt kính cửa sổ xe nhìn ra ngoài. Phía trước, dải đồng bằng phẳng lặng bắt đầu hiện rõ dưới làn sương sớm mờ ảo. Và rồi, một vệt sáng màu hồng cam rực rỡ bất ngờ vạch ngang bầu trời phía đông, xé toạc màn sương mờ mịt.

Ánh bình minh đầu tiên của ngày mới bắt đầu tràn xuống mặt đất.

Đôi mắt trong vắt của Naomi mở to hết cỡ, tràn ngập sự ngỡ ngàng. Không dám gọi to vì sợ đánh thức mọi người, cô bé chỉ dám khẽ kéo nhẹ vạt áo của Suzuno bên cạnh, cất giọng thì thầm run run: ``Suzuno-oneechan ơi\dots\ nhìn kìa\dots\ Trời sáng rồi ạ\dots''

Suzuno mắt nhắm mắt mở, khẽ nghiêng đầu nhìn theo hướng ngón tay nhỏ xíu của em gái chỉ ra ngoài cửa sổ, trên môi nở một nụ cười ngái ngủ dịu dàng: ``Ừm\dots\ Đẹp quá phải không Naomi-chan\dots''

Từ hàng ghế sau, Kaiho - vốn dĩ chẳng hề chợp mắt suốt cả chặng đường - vẫn khoanh tay nhìn ra cửa sổ đối diện. Không quay đầu lại, nhưng giọng cô nàng trầm ấm cất lên, mang theo một chỉ dẫn nhẹ nhàng: ``Nhìn chếch sang hướng đông nam một chút nữa đi, Nao.''

Naomi ngoan ngoãn làm theo, áp sát cả khuôn mặt nhỏ nhắn vào tấm kính.

Ở phía chân trời xa tít tắp nơi vệt nắng bình minh đang nhuộm vàng cả một góc trời, dải rừng cây bỗng mở toang ra. Không khí luồn qua khe gió điều hòa bắt đầu mang theo một mùi hương rất lạ: vị mằn mặn, nồng nàn và tươi mát của gió biển.

Và rồi, khi chiếc xe vượt qua con dốc thoai thoải cuối cùng, biển cả bao la chính thức vỡ òa trước mắt họ.

Mặt nước xanh biếc trải dài vô tận dưới ánh bình minh rực rỡ, những con sóng bạc đầu thi nhau xô vào dải cát vàng óng ả, bọt tung trắng xóa như một dải lụa mềm mại.

Đoàn xe rẽ vào con đường rợp bóng dừa dẫn thẳng vào bãi đỗ của khu nghỉ dưỡng ven biển Kaiwa (海和 - Hải Hòa) thuộc địa phận Seika (清化 - Thanh Hóa). Mọi người lục tục bước xuống xe trong trạng thái nửa tỉnh nửa mê sau một chuyến đi dài xuyên đêm.

Thế nhưng, khi vừa đặt chân vào sảnh lớn của khách sạn nơi có bức tường kính khổng lồ nhìn thẳng ra toàn cảnh vịnh biển lộng gió, bước chân của ba chị em nhà Hikawa bỗng đồng loạt khựng lại.

Naomi đứng chết trân ngay giữa sảnh. Chiếc nón cói và bọc đồ chơi trên tay cô bé rơi bịch xuống nền gạch hoa cương lúc nào không hay. Cô bé không chớp mắt, đôi con ngươi màu lục bảo phản chiếu trọn vẹn khoảng không gian bao la, thăm thẳm của biển trời mà trước đây em chỉ được thấy qua những trang sách minh họa chật hẹp hay màn hình tivi bé tí xíu.

Đôi vai nhỏ bé của Naomi khẽ run lên. Rất chậm rãi, cô bé bước tới nép sát vào sau lưng Suzuno, hai tay túm chặt lấy vạt áo chị gái, giọng nói nghẹn ngào vỡ òa trong niềm xúc động thuần khiết: ``Suzuno-oneechan\dots\ Biển\dots\ Biển ngoài đời thật lớn quá\dots\ Lớn hơn trong sách nhiều, nhiều lắm luôn ạ\dots!''

Suzuno đứng chết lặng trước khung kính lớn, đôi bàn tay mảnh mai buông thõng bên hông khẽ run rẩy. Gió biển từ ô cửa thông gió thổi vào làm tung bay mái tóc dài của cô. Đôi mắt hai màu dị sắc của cô gái trẻ ngập tràn những con sóng lung linh. Đã bao nhiêu lần, trong những đêm đông rét mướt ở căn nhà trọ mục nát năm xưa, khi tiếng mưa dột tong tỏng xuống chiếc chậu nhôm thủng đáy, Kaiho đã ôm cô vào lòng và kể cho cô nghe về biển? Về một nơi có bãi cát mịn như nhung, có làn nước trong vắt và những con ốc biển biết hát? Khi ấy, đối với Suzuno, biển là một điều gì đó xa vời như một câu chuyện cổ tích hoang đường.

Và giờ đây, cổ tích ấy đang hiện hữu ngay trước mắt cô, sống động và rực rỡ đến từng giọt nắng.

Suzuno khẽ cúi xuống, nắm chặt lấy bàn tay nhỏ nhắn của Naomi, những giọt nước mắt hạnh phúc khẽ lăn dài trên gò má ửng hồng: ``Ừ\dots\ Lớn lắm Naomi-chan\dots\ Đẹp hơn tất cả những gì chúng ta từng tưởng tượng\dots''

Cách đó vài bước chân, tách biệt khỏi đám đông đang ồn ào nhận phòng, Kaiho đứng lặng lẽ một mình bên cây cột lớn.

Cô nhìn thẳng ra đường chân trời nơi những con sóng đang miệt mài vỗ bờ. Đôi mắt màu lam sâu thẳm của cô gái 26 tuổi không ngấn lệ, nhưng cơ hàm bạnh ra, sống mũi cay cay. Lời hứa của bố mẹ mười năm trước\dots\ giấc mơ dang dở tưởng chừng đã vĩnh viễn bị chôn vùi trong tro bụi của quá khứ\dots\ cuối cùng hôm nay cô đã thay họ thực hiện trọn vẹn cho em gái mình.

Bàn tay Kaiho siết chặt lấy chiếc nơ cầu vồng trước ngực, rồi từ từ buông lỏng ra. Một nụ cười nhẹ nhõm, thanh thản và đẹp đẽ nhất mà tôi từng thấy kể từ ngày gặp lại cô nàng khẽ nở trên môi.

Đúng lúc ấy, Ryuuhou từ phía sau xách ba-lô lao vút ra mép hiên sảnh kính, hai tay dang rộng hít một hơi thật sâu rồi hét toáng lên đầy sảng khoái: ``\textbf{Vậy đây là biển ha~! Gió mát rượi luôn này!}''

Sự phấn khích của cô em thứ năm như một ngòi nổ kéo theo cả đoàn người phía sau ùa tới. Kson khoác vai Luna sải bước tới, ngắm nhìn dải cát trắng mịn màng rồi huýt sáo một tiếng giòn giã: ``Coi bộ trông cũng đẹp lắm nha! Bãi biển Kaiwa này không khí thoáng đãng, lại còn vắng người nữa chứ, quá đã!''

Luna chớp chớp mắt, hai má phúng phính phụng phịu vẻ thích thú: ``Nại còn vắng ngừi nữa-nanora! Tha hồ bơi lội mà không sợ bị chen chúc nanor\~{}!''

Kanata, Watame và Towa cũng lục tục kéo đến sát vách kính. Dù đây chẳng phải lần đầu tiên họ ra biển, nhưng đứng trước khung cảnh vịnh biển Seika hoang sơ, thanh bình lúc bình minh, cả ba cô nàng Thế hệ thứ Tư đều không giấu nổi ánh mắt trầm trồ.

Watame hít hà mùi gió biển mằn mặn, hai mắt long lanh: ``Tuyệt vời thật đấy\dots\ Biển sớm thế này làm em muốn nhảy ùm xuống bơi một vòng ghê!''

Towa khoanh tay, khẽ mỉm cười gật gù: ``Ừm, yên tĩnh hơn hẳn mấy bãi biển đông đúc gần thành phố Điện tử. Chuyến đi này đúng là một lựa chọn sáng suốt.''

Kanata thì vung nắm đấm nhỏ lên không trung, tràn đầy khí thế: ``Tốt lắm! Lát nữa nhận phòng xong, tôi nhất định sẽ kéo cả hội đi dạo một vòng quanh bờ cát!''

Ở phía sau cùng, Chloe đứng lặng trước cửa ra vào, tay kéo chiếc vali nhỏ. Mùi mặn nồng quen thuộc của biển cả phả vào cánh mũi đánh thức toàn bộ bản năng của cô nàng sát thủ cá voi. Đôi mắt đỏ rực của cô sáng bừng lên, khóe môi khẽ cong thành một nụ cười rạng rỡ, tự nhiên nhất kể từ lúc lên xe: ``Hừm, đúng là mùi của biển rồi\dots\ Lâu lắm rồi mới được hít thở cái không khí thân thương này, poe poe!''

Đứng cạnh vách kính, chứng kiến trọn vẹn phản ứng từ sững sờ, xúc động cho đến vỡ òa phấn khích của tất cả mọi người, tôi khẽ nghiêng đầu nhìn sang Kaiho. Cô em gái song sinh của tôi cũng vừa lúc quay sang bắt gặp ánh mắt tôi.

Không cần một lời nói nào, chỉ bằng một cái nhướng mày và nụ cười kín đáo trao nhau, cả tôi và Kaiho đều ngầm hiểu: chuyến đi này, dù có phải trải qua bao nhiêu vất vả hay mệt nhọc trên xe khách đường dài, cũng đã hoàn toàn xứng đáng.

Đại dương rộng lớn đã mở rộng vòng tay đón nhận đại gia đình Hikawa cùng những người bạn như thế đó.

\ct

Khách sạn nơi chúng tôi dừng chân là một quần thể nghỉ dưỡng khá rộng lớn, chia làm hai dãy nhà ba tầng với khoảng năm chục phòng nghỉ, hai khu nhà ăn lớn, vài khu vui chơi giải trí dành cho trẻ em ở tầng trệt và một bể bơi nhỏ ngoài trời—mà thực chất chỉ đóng vai trò là bể tráng nước ngọt sau khi du khách đi tắm biển về. Cơ sở vật chất ở đây mộc mạc và mang đậm phong cách thập niên trước, khác xa với những khách sạn cao tầng hiện đại đầy đủ tiện nghi mà chúng tôi thường thấy ở trung tâm các đô thị lớn.

Sau một chuyến đi dài suốt đêm, cả đoàn ai nấy đều thấm mệt và đói cồn cào. Theo hướng dẫn của ban tổ chức, mọi người kéo nhau vào khu nhà ăn lớn ở tầng một để chuẩn bị dùng bữa điểm tâm sáng kết hợp ăn nhẹ lấy lại sức.

Thế nhưng, ngay lối vào cạnh sảnh lễ tân, tôi và Kson bỗng đồng loạt khựng bước lại.

Đập vào mắt chúng tôi là một khung bảng gỗ mạ vàng to đùng treo ngay ngắn trên tường. Bên trái là bản ``NỘI QUY KHÁCH SẠN'' in bằng tiếng mẹ đẻ chuẩn chỉnh không một tì vết; còn bên phải lại là phiên bản dịch thuật sang tiếng Anh với dòng tiêu đề trang trọng: \textit{``Hotel Regulations''}.

\begin{figure}[H]
    \centering
    \includegraphics[width=12.5cm]{../../../../Image/BookB/hotel_rules.jpg}
\end{figure}

\imgc{hotel_rules}

Sẽ chẳng có gì đáng nói nếu tấm biển nội quy tiếng Anh đó được dịch thuật một cách bình thường. Nhưng khốn nỗi, với một kẻ vốn có chứng chỉ tiếng Anh giao tiếp chuẩn chỉnh như tôi, và một người từng sống nửa đời người ở Mỹ như nàng cựu Tổng trưởng Kson, tấm bảng này chẳng khác nào một bãi mìn ngữ pháp khiến người đọc phải tăng xông máu não.

Kson khoanh hai tay trước ngực, dí sát mặt vào tấm bảng kính, khóe mắt giật giật liên hồi: ``Này Kuon-paisen\dots\ Anh nhìn cái dòng tiêu đề nhỏ dưới tên khách sạn xem? \textit{`In order to be safe and convenient and provide best services for guests\dots\ you are kindly requested to follow the following regulations'}\dots\ Nghe đã lủng củng như google dịch thời kỳ đồ đá rồi!''

Tôi cũng ôm trán thở dài, đưa ngón tay chỉ thẳng vào điều số 2: ``Chưa hết đâu em. Nhìn điều số hai này: \textit{`Before receivinh the room\dots'} Từ \textit{`receivinh'} này là tiếng nước nào vậy!? Chữ `g' biến thành chữ `h' từ bao giờ thế!?''

Kson được đà chỉ tiếp xuống điều số 3, cười sặc sụa: ``Đây này! Điều số ba: \textit{`You are requested toleave valuable assets\dots\ If noy, the hotel will not responsible for the loss.'} Ôi lạy Chúa tôi! \textit{`Toleave'} dính liền không thèm cách, còn \textit{`if not'} thì viết thành \textit{`if noy'}! Đã thế lại còn thiếu luôn động từ to-be ở vế sau nữa chứ!''

Tôi lướt mắt xuống điều số 4 và số 6, da gà da vịt nổi rần rần: ``Chưa dừng lại ở đó đâu! Điều số bốn: \textit{`changing rooms additionnal guests'} - chữ \textit{`additional'} khuyến mãi thêm một chữ `n' to tướng. Còn đỉnh cao nhất là điều số sáu đây này: \textit{`the staying of your guests must be fini shed before 22h00'}! \textbf{Fini shed là cái quái gì!? Chữ \textit{`finished'} bị chém đôi ra thành hai từ như chém chuối thế kia à!?} Rồi phía sau còn có quả \textit{`befone being going out ofthe room'} nữa! \textit{`Befone'} với \textit{`ofthe'}, vừa sai chính tả vừa dính chữ!''

Kaiho đứng khoanh tay bên cạnh, mặt không cảm xúc nhìn hai chúng tôi đang thi nhau săm soi từng chân tơ kẽ tóc của tấm bảng, khẽ cất giọng khô khốc: ``Thôi nào hai người. Đừng có đứng đây bắt bẻ người ta như giám thị chấm thi đại học thế chứ. Dù sao thì người nước ngoài đọc lướt qua vẫn hiểu đại ý mà.''

Kson lập tức giãy nảy lên như đỉa phải vôi, ngón tay chọc lia lịa vào điều số 10: ``Hiểu cái con khỉ mốc ấy! Chị nhìn điều số mười này đi! \textit{`In case you have to bring in weapons for business\dots'} Đi nghỉ dưỡng biển mà lại có vụ mang vũ khí vì mục đích kinh doanh là thế nào!? Rồi cuối cùng còn có quả \textit{`declare them to the recception'} - lại thêm một chữ `c' ở quầy lễ tân nữa! Em còn chưa nói tới việc mang vũ khí mà còn phải đi trình báo trước lễ tân đấy! Khác nào `lạy ông tôi ở bụi này' không!?''

Tôi tiếp lời bồi thêm một cú chót: ``Và hãy chiêm ngưỡng điều số 11 thần thánh: \textit{`You are requested not to in the rooms.'} \textbf{Not to cái gì trong phòng!?} Không được làm gì!? Động từ chính đâu mất rồi!? Đang viết dở thì lăn ra ngủ hay sao mà cụt lủn thế kia!?''

Kaiho nhìn theo ngón tay tôi chỉ vào điều số 11, khóe môi cô nàng khẽ giật giật một hồi rồi đành ngậm ngùi buốt ruột thừa nhận: ``\dots Đúng thật. Quả này thì chịu chết, không đỡ nổi.''

Đúng lúc đó, một ngón tay nhỏ xíu khẽ chạm vào mặt kính của khung tranh.

Bé út Naomi kiễng gót chân, mắt mở to tròn xoe, chiếc mũi nhỏ phập phồng đếm từng dòng một cách vô cùng nghiêm túc: ``Một, hai, ba, bốn\dots\ Ơ? Onii-chan\dots\ Ở bên bảng tiếng mình có bốn mục lớn La Mã, tổng cộng có tới mười sáu điều lận\dots\ Sao sang bên tiếng Anh chỉ có mười một điều thôi ạ? Thiếu mất tận năm điều luôn rồi kìa!''

Tôi và Kson khựng lại, ngơ ngác nhìn theo ngón tay của cô em út. Quả đúng là như thế thật! Bản tiếng mẹ đẻ được phân chia vô cùng khoa học thành bốn phần rõ ràng: từ an ninh khách sạn, tài sản của khách, các hành vi nghiêm cấm cho đến thủ tục thanh toán. Thế nhưng khi phóng tác sang bản tiếng Anh, người biên soạn đã hồn nhiên gộp tuốt tuột lại thành một danh sách độc nhất mười một điều, và cắt phăng đi năm điều quan trọng nhất lúc nào không hay!

Đứng cạnh Naomi, cặp song sinh Suzuno và Ryuuhou cũng nghiêng đầu nhìn tấm bảng. Vốn sở hữu vốn tiếng Anh học thuật và giao tiếp rất cừ khôi - một người tiếp xúc hàng ngày với tài liệu hóa sinh quốc tế, còn một người lăn lộn trong giới underground nghe nhạc Âu Mỹ suốt từ năm 12 tuổi - hai cô gái trẻ chỉ mất vài giây để nắm bắt toàn bộ thảm họa văn bản.

Suzuno đưa ngón tay thon thả lên đỡ cằm, đôi mắt dị sắc ánh lên vẻ bối rối pha lẫn buồn cười. Cô bé khẽ nghiêng đầu nhận xét bằng chất giọng dịu dàng, lễ phép thường ngày: ``Dạ\dots\ Em thấy về mặt cấu trúc văn bản thì bên tiếng Anh bị lỗi định dạng nghiêm trọng ạ. Họ không hề chia ra các phân mục cụ thể như bản gốc, lại còn dịch lược bỏ hoàn toàn phần hướng dẫn nhân viên kiểm tra thiết bị khi nhận phòng nữa\dots\ Nếu khách nước ngoài đọc vào chắc chắn sẽ bị bối rối về quy trình an toàn sinh hoạt cho mà xem.''

Trong khi đó, Ryuuhou thì chẳng thèm nể nang gì, trực tiếp bật cười thành tiếng, khoanh tay trước ngực hất cằm châm chọc: ``Trời đất ơi, tiếng Anh này là tiếng Anh bồi phiên bản miệt vườn rồi chứ dịch thuật nỗi gì! Câu cú thì lủng cà lủng củng, chia thì lung tung xà bèo, đã thế người dịch còn lười đến mức cắt xén bớt điều khoản để đỡ phải gõ phím nữa chứ! Em mà viết bài luận tiếng Anh ở trường kiểu này chắc cô giáo cho âm điểm liệt luôn quá, Onii!''

Cả đám chúng tôi cứ thế vừa đứng vạch lỗi sai, vừa đếm số lượng điều khoản thiếu sót rồi bình luận rôm rả đến mức mặt đỏ tía tai, khiến hai cô nhân viên lễ tân đứng sau quầy chỉ biết cúi gằm mặt xuống đất, mồ hôi hột to như hạt ngô chảy ròng ròng trên trán. Nhìn vẻ mặt của hai cô gái tội nghiệp, có lẽ họ đang khóc thầm trong lòng vì bị một đại gia đình khách bản địa toàn những tay soi mói cự phách kết hợp với một bà chị ngoại quốc chỉ lỗi ngữ pháp đến mức tơi tả không chừa một góc nào.

Trong khi đó, bốn cô nàng Thế hệ thứ Tư đứng phía sau ngơ ngác hoàn toàn. Không đọc hiểu được cả tiếng bản địa lẫn mớ tiếng Anh chắp vá kia, họ chỉ biết tròn mắt nhìn hai chuyên gia thẩm định ngôn ngữ đang thao thao bất tuyệt.

Luna ngơ ngác ngậm kẹo mút, cất giọng thỏ thẻ: ``Họ đang nói zì thế-nanora\dots? Sao trông hai người đó hung hăng zợ\dots?''

Towa khoanh hai tay, thở dài lắc đầu e ngại: ``Trông họ đáng sợ thật đấy\dots\ Cứ như chuẩn bị lật tung cả cái sảnh khách sạn lên không bằng\dots''

Kanata giật giật góc áo Watame, ngập ngừng: ``Chắc là\dots\ chúng ta nên vào nhà ăn trước đi nhỉ? Mọi người trong đoàn đang chờ đông đủ rồi kìa\dots''

Watame chớp chớp mắt, tủm tỉm cười: ``Nhìn kìa, đến cả Kaiho-senpai mặt cũng ngớ người ra vì bất lực kìa, hihi\dots''

Chloe thì đứng xéo góc, tay cầm que kem cắn dở, mắt đảo qua đảo lại giữa Kuon và Kson rồi lẩm bẩm: ``Poe\dots\ Tên Tổng quản này bình thường lười biếng mà đụng tới mấy lỗi vụn vặt này thì soi kinh khủng khiếp thật đấy\dots''

Sau khi mẹ tôi từ trong nhà ăn bước ra lườm cho một phát sắc lẹm, tôi và Kson mới chịu ngậm miệng, ngoan ngoãn kéo vali theo sau mọi người vào dùng bữa sáng.

\ct

Khoảng chín giờ ba mươi phút sáng, sau khi dùng bữa và nghỉ ngơi nhẹ, người cậu ruột của tôi - em trai của mẹ kiêm giám đốc điều hành kho hàng của công ty - bước đến chỉ dẫn đoàn chúng tôi lên nhận phòng nghỉ.

Theo sự sắp xếp của mẹ và ban tổ chức, đại gia đình chúng tôi cùng nhóm khách mời được phân chia vào các phòng ở tầng bốn:

Căn phòng lớn mang số hiệu \textbf{420B} là nơi cư trú của gia đình Hikawa gồm: mẹ, bà ngoại, tôi, Kaiho, Suzuno, Ryuuhou và bé út Naomi.

Căn phòng đối diện mang số hiệu \textbf{404B} được dành riêng cho hội chị em Thế hệ thứ Tư gồm: Kson, Kanata, Watame, Towa và Luna.

Còn một căn phòng nhỏ liền kề mang ký hiệu \textbf{405B} thì được mẹ ưu ái dành riêng cho cô nàng cá voi sát thủ Sakamata Chloe, với lý do cô bé đang thi hành ``nhiệm vụ gián điệp đặc biệt'' cần không gian bí mật để báo cáo cho cấp trên, nhưng thực chất là vì mẹ tôi rất thương con bé lủi thủi một mình nên muốn tạo cho cô nàng một không gian thật thoải mái. Thực ra thì, tại sao mẹ tôi lại gộp cậu ta với phòng của hội Gen 4 chứ?

Tôi đứng trước cửa phòng 420B, nhìn tấm biển số phòng rồi lại nhìn quanh dàn phụ nữ trong nhà, không khỏi thở dài thườn thượt trong lòng: ``\textit{Lại thế nữa rồi\dots}''

Lại một lần nữa, giữa cả một căn phòng đầy ắp phụ nữ từ già đến trẻ, tôi là thằng con trai duy nhất. Gần như năm nào mỗi lần gia đình tổ chức đi nghỉ mát, kịch bản trớ trêu này cũng lặp lại y hệt. Nhưng năm nay, khi trong phòng có thêm cả Kaiho, Suzuno và Naomi, tôi cảm giác gánh nặng trách nhiệm và mức độ phải giữ kẽ của mình còn tăng lên gấp bội phần.

Về cơ bản, các phòng nghỉ ở đây đều có cấu trúc tương tự nhau: mỗi phòng được bố trí hai chiếc giường đôi lớn bằng gỗ ép trải nệm trắng muốt, diện tích vừa vặn đủ cho năm đến sáu người sinh hoạt thoải mái. Trong phòng trang bị đầy đủ những tiện nghi tối thiểu như tủ lạnh mini, điều hòa nhiệt độ hai chiều, quạt treo tường, bàn trà nhỏ và một chiếc tủ quần áo âm tường.

Điểm đặc biệt nhất chính là khu vực nhà tắm liên kết trực tiếp với nhà vệ sinh, chỉ ngăn cách bằng một tấm vách kính mờ và rèm kéo mỏng manh. Với gia đình tôi thì điều này đã quá quen thuộc, nhưng đối với các tiểu thư thần tượng ở phòng đối diện, đây lại là một thử thách tâm lý không hề nhỏ.

Từ hành lang đối diện, cánh cửa phòng 404B đang mở toang, vọng ra những tiếng xì xào bàn tán đầy bối rối của các cô gái Gen 4.

Towa đứng tựa cửa nhà tắm, gương mặt đỏ ửng ngượng ngùng: ``Tôi vẫn thấy nó cứ cấn cấn thế nào ấy\dots\ Cửa kính mỏng tanh thế này thì tắm rửa kiểu gì trời\dots''

Kanata gật gù phụ họa, hai tay ôm má: ``Cứ như thể\dots\ toàn bộ sự riêng tư bị phơi bày hết ấy nhỉ\dots''

Watame khoanh tay đứng cạnh, bỗng nhếch mép nở một nụ cười gian xảo vô cùng khả nghi: ``Có gì đâu mà phải ngại chứ! Toàn con gái với nhau thì sợ cái gì! Với cả như thế này thì mình tha hồ được ngắm nhìn các bồ tắm rửa và đi vệ sinh thoải mái luôn\dots\ Hehe\dots''

Kson đang cởi áo khoác ngoài, nghe thấy thế liền quay ngoắt lại trừng mắt: ``Bồ đang cười cái điệu nham hiểm gì thế hả con Cừu kia!?''

Watame lập tức giật thót, làm mặt ngây thơ vô tội xua tay lia lịa: ``Đ-Đâu có gì đâu! Tôi có nói gì đâu nà\~{}''

Thế nhưng, có vẻ như chuyến xe đưa chúng tôi đến nơi quá sớm so với giờ nhận phòng tiêu chuẩn, nên khi bước vào bên trong, cả hai căn phòng vẫn chưa được dọn dẹp xong. Ga giường cũ vẫn còn vương vãi và rác từ đợt khách trước chưa kịp thu gom. Trên cả tầng ba rộng thênh thang chỉ có vỏn vẹn hai bác nhân viên tạp vụ lớn tuổi đang hì hục lau dọn từng phòng một cách vất vả.

Tôi không hề phàn nàn hay khó chịu. Hai người mà phải cáng đáng cả chục căn phòng rộng lớn thế này thì chuyện chậm trễ là điều khó tránh khỏi.

Tôi liền xắn tay áo lên: ``Thôi được rồi, để con phụ một tay. Cả nhà mình cùng dọn cho nhanh để còn có chỗ ngả lưng.''

Mẹ tôi, Kaiho và Suzuno lập tức hưởng ứng. Mẹ tôi thoăn thoắt gom đống vỏ gối cũ, Suzuno cẩn thận dùng khăn ẩm lau sạch từng mặt bàn ghế và tủ kệ, trong khi Kaiho chỉ bằng một tay đã dễ dàng nhấc bổng cả chiếc giường gỗ lên để quét sạch lớp bụi bám ở góc tường phía dưới. Nhờ sự hiệp lực nhanh gọn của cả nhà, chỉ chưa đầy mười lăm phút sau, căn phòng 420B đã trở nên sạch bong sáng bóng, thơm tho mùi tinh dầu sả chanh dịu nhẹ.

Đấy là tình hình bên phòng chúng tôi. Còn ở căn phòng 404B đối diện thì lại là một câu chuyện hoàn toàn khác.

Đang lúc hai bác nhân viên dọn phòng chuẩn bị bê xô nước bước vào thì cô nàng Luna bỗng ưỡn ngực bước ra giữa phòng. Cô nàng giơ cao chiếc kẹo mút trên tay như một cây đũa phép thần kỳ, xoay một vòng điệu đà rồi dõng dạc hô to bằng chất giọng trẻ con dễ thương: ``\textbf{Úm-ba-la-xì-bùa! Phòng sạch sẽ đi nào-nanora!!}''

Một luồng sáng lấp lánh như bụi tiên bất ngờ lóe lên khắp bốn góc phòng. Chỉ trong chớp mắt, toàn bộ rác rưởi tự động bay vào thùng rác, ga trải giường tự vuốt phẳng phiu không một nếp nhăn, và sàn nhà sáng bóng đến mức có thể soi gương được.

Hai bác nhân viên tạp vụ đứng chết trân ngoài cửa, hai mắt trợn trừng như muốn lồi ra ngoài, chiếc giẻ lau trên tay rơi bẹp xuống sàn nhà: ``Ối\dots\ Ối trời đất quỷ thần ơi! C-Cái gì vừa xảy ra thế này!?''

Đứng ngay cửa phòng đối diện chứng kiến toàn bộ màn trình diễn ma thuật công khai ấy, tôi chỉ biết ôm trán chẹp miệng thở dài ngao ngán: ``Đã dặn là đừng có làm lộ chân tướng rồi cơ mà\dots''

Trong khi đó, ở phòng 405B bên cạnh, Chloe cũng đã tự mình dọn dẹp xong xuôi căn phòng nhỏ với tốc độ đáng kinh ngạc của một sát thủ chuyên nghiệp. Cô nàng đóng sập cửa lại, lập tức giơ điện thoại lên chụp một bức ảnh selfie tự đắc với góc phòng gọn gàng, rồi gửi thẳng vào nhóm chat bí mật của holoX kèm dòng trạng thái:

\begin{quote}
    \textit{`[BÁO CÁO TIẾN ĐỘ]: Căn cứ ẩn náu bí mật của đặc vụ Sakamata Chloe đã được thiết lập thành công ngay cạnh phòng của mục tiêu Hikawa Kuon. Giai đoạn một hoàn tất xuất sắc. Đang tiếp tục giám sát đối tượng!'}
\end{quote}

Nhìn chuỗi những trò hề ngớ ngẩn và những pha xử lý cồng kềnh ngay trong buổi sáng đầu tiên đặt chân đến khách sạn, tôi chỉ biết lắc đầu cười trừ. Chuyến đi này chắc chắn sẽ còn vô số kiếp nạn dở khóc dở cười đang chờ đón phía trước.

\ct

Bốn giờ chiều.

Căn phòng 420B chìm trong không gian lười biếng của buổi chiều hè oi ả. Sau bữa trưa no nê, mẹ và bà ngoại đã đánh một giấc say sưa trên chiếc giường phía trong. Ryuuhou thì nằm lăn lóc dưới thảm, cắm mặt vào chiếc điện thoại của tôi để cày game âm nhạc, chốc chốc lại lẩm bẩm gắt gỏng khi bấm trượt một nốt nhịp điệu.

Suzuno đang ngồi ở mép bàn trà, kiên nhẫn bóc từng múi quýt ngọt đút cho Naomi vừa mới ngủ dậy, hai chị em trò chuyện khe khẽ bằng những thanh âm trong trẻo như tiếng suối.

Tôi và Kaiho thì nằm dài ở chiếc giường còn lại, gối đầu lên hai góc thành giường cùng xem một bộ phim hành động chiếu dở trên màn hình ti-vi treo tường. Thỉnh thoảng, Kson hoặc Luna từ phòng đối diện lại lon ton chạy sang phòng chúng tôi buôn dưa lê vài câu ba điều bốn chuyện rồi lại lượn về. Cứ thế, cả buổi chiều trôi qua trong sự êm đềm đến mức nhàm chán.

Đột nhiên, Kaiho xoay người nằm nghiêng, chống một khuỷu tay xuống nệm nhìn tôi, cất giọng than thở: ``Mà này, anh Kuon\dots''

Tôi dán mắt vào màn hình ti-vi, đáp gọn: ``Hửm?''

``Em thấy chán\~{} quá đi mất!'' cô em gái song sinh phụng phịu, đôi chân dài co duỗi đá nhẹ vào mép đệm.

``Thế em muốn làm gì?'' tôi quay sang hỏi.

Không chần chừ một giây, Kaiho giơ một ngón tay chỉ thẳng qua ô cửa sổ hướng về phía bờ biển rực nắng bên ngoài, đôi mắt xanh lam sáng rực lên đầy háo hức: ``Biển! Em muốn ra ngoài đó để tắm biển!''

Tôi khẽ thở dài. Tất nhiên rồi, với một người lần đầu tiên trong đời được đặt chân đến biển như cô nàng, việc cuồng chân muốn lao ngay xuống làn nước mát lạnh là điều hoàn toàn dễ hiểu. Nhưng khốn nỗi, với một kẻ nội hướng thích giam mình trong phòng máy lạnh như tôi, việc phải phơi mình dưới cái nắng gay gắt ngoài kia chẳng có chút hấp dẫn nào.

Đúng lúc tôi đang định tìm cớ thoái thác thì giọng mẹ tôi từ giường bên bỗng cất lên ngái ngủ: ``Mấy giờ rồi con?''

Tôi liếc nhìn đồng hồ góc màn hình điện thoại: ``Bốn giờ chiều rồi mẹ ạ.''

Bà lập tức vén chăn ngồi dậy, vươn vai dứt khoát: ``Vậy thì chuẩn bị quần áo bơi đi các con.''

Tôi chớp mắt ngơ ngác: ``\dots Dạ?''

``Đi bơi chứ còn sao nữa! Đến biển mà cứ nằm ru rú trong phòng thế này thì đi làm cái gì!'' mẹ tôi vừa xỏ dép vừa cất giọng chỉ đạo đanh thép.

Thấy mẹ ra lệnh, Kaiho sướng rơn người reo lên một tiếng đầy khoái chí, lập tức bật dậy như một chiếc lò xo. Naomi cũng không kém phần phấn khích; cô em út hai mắt sáng long lanh như hai hòn ngọc, nhảy cẫng lên khỏi chiếc ghế bành, hai bàn tay nhỏ xíu vỗ bôm bốp reo hò: ``Đi bơi! Đi tắm biển thật rồi! Em được chạm vào sóng biển rồi!!''

Suzuno ngồi bên cạnh khẽ bật cười hiền từ, nét mặt rạng rỡ hẳn lên. Dù vẫn giữ tác phong điềm đạm, đoan trang thường ngày, đôi bàn tay cô bé thoăn thoắt xếp lại đĩa hoa quả rồi đứng dậy kiểm tra kỹ càng từng tuýp kem chống nắng và chuẩn bị sẵn những chiếc khăn bông lớn: ``Vâng ạ! Để con bôi kem chống nắng cho Naomi và chuẩn bị khăn tắm đem theo nhé mẹ.''

Ryuuhou thì vội vã vứt chiếc điện thoại sang một bên giường, hí hửng nhảy bổ lại chiếc vali to đùng lôi bộ đồ bơi của mình ra: ``Tuyệt vời ông mặt trời! Nằm lì trong phòng cả buổi em sắp mốc meo cả người rồi đây này!''

Nhìn cảnh tượng cả nhà rần rần hưởng ứng, tôi chỉ biết ngậm ngùi thở dài đầu hàng số phận.

Mẹ quay sang tôi, dặn dò: ``Kuon, con sang phòng đối với phòng của bé Chloe gọi mọi người chuẩn bị ra bãi biển luôn đi.''

Tôi gật đầu, bước ra ngoài hành lang rồi gõ cửa phòng 404B.

*CỘC CỘC*

``Ai đấy?'' Giọng Kson vọng ra từ phía trong.

Tôi tựa lưng vào mép cửa: ``Anh đây chứ ai. Chuẩn bị đồ bơi đi, cả nhà chuẩn bị ra biển tắm đây!''

Từ bên trong lập tức vang lên tiếng reo hò ríu rít của nhóm Gen 4. Luna hí hửng hô to: ``\textbf{Kuon-senpai! Đi bơi-nanora!!}'' Watame thì vừa cười khúc khích vừa phụ họa theo. Tôi nghe thấy tiếng cười đầy nham hiểm của cô nàng Cừu mà không khỏi rùng mình; đối tượng này nhìn ngoài mặt thì ngây thơ thảo mai, nhưng độ tinh quái và hám sắc thì có khi chẳng thua kém bất kỳ ai.

Sau khi tạt qua gõ cửa báo cho Chloe ở phòng bên cạnh, tôi xoay người đẩy cửa bước về lại phòng 420B: ``Anh báo xong rồi đấy, mọi người chuẩn bị--''

Thế nhưng, vừa mở toang cánh cửa bước vào, mắt tôi lập tức va phải một khung cảnh chấn động tâm can.

Suzuno và Naomi lúc này đã dẫn nhau vào bên trong phòng vệ sinh đóng kín cửa để thay đồ cho kín đáo. Ở gian ngoài phòng ngủ, chỉ còn lại Kaiho và Ryuuhou đang hồn nhiên thay đồ bơi.

Kaiho đang đứng ngay cạnh mép giường, hai tay vừa kéo chiếc quần lót bơi thể thao màu tím than ngang đùi, để lộ trọn vẹn những đường cong cơ thể săn chắc, khỏe khoắn cùng làn da trắng mịn màng. Trong khi đó, Ryuuhou đứng ngay phía sau lưng chị Hai, hai tay đang vòng ra sau lưng lột chiếc áo phông ngắn tay qua khỏi đầu, để lộ chiếc áo lót bơi màu đen cùng vòng eo thon gọn của một cô nàng dancer underground!

Một cảnh tượng slapstick kinh điển hệt như trong các bộ anime hài hước lãng mạn mà tôi thường hay xem trên mạng xã hội!

Cả hai cô em gái chỉ kịp tròn mắt ``Ơ'' lên một tiếng đồng thanh.

Tôi hoảng hồn như bị điện giật, lập tức quay ngoắt người kéo rầm cánh cửa lại, tim đập thình thịch muốn nhảy ra khỏi lồng ngực. Suýt chút nữa thì toi mạng!

Từ bên trong phòng, giọng Kaiho vọng ra tỉnh bơ: ``Em biết là anh rồi mà! Anh em một nhà cả, không sao đâu!''

Thế nhưng Ryuuhou thì chẳng thèm dễ dãi như chị Hai. Tiếng cô em thứ tư gắt toáng lên vọng qua khe cửa đầy đanh đá: ``Không sao cái đầu chị ấy, chị Kaiho! Anh Kuon kia! Bộ anh bị đứt dây thần kinh gõ cửa rồi đấy à!? \textbf{Lần sau trước khi xông vào phòng con gái thì phải gõ cửa đoàng hoàng nghe chưa, đồ biến thái ngốc nghếch!!}''

Mặt tôi nóng bừng bừng, tựa lưng vào cánh cửa gào lên thất thanh: ``\textbf{Anh vừa mới gõ cửa lúc nãy xong mà!! Ai mà ngờ mấy đứa lại thay đồ ngay giữa phòng khách cơ chứ!? Biết xấu hổ và giữ ý tứ giùm anh cái đi hai cô nương!!}''

Chưa kịp chạm chân xuống nước biển mà tôi đã phải đối mặt với một kiếp nạn thót tim thế này rồi. Thật chẳng dám tưởng tượng mấy ngày tiếp theo sẽ còn hỗn loạn đến mức nào nữa.

\ct

Bãi biển Kaiwa, bốn giờ ba mươi phút chiều.

Khác hẳn với vẻ tĩnh lặng của buổi sáng sớm, bãi biển Seika lúc này đã bắt đầu nhộn nhịp hơn với những nhóm du khách tản bộ. Những con sóng bạc đầu ở vùng biển này cuộn trào mạnh mẽ và dữ dội hơn hẳn so với vùng biển nông quanh thành phố Điện tử, gió biển lồng lộng thổi tung những lọn tóc bồng bềnh.

Cả đoàn chúng tôi vừa đặt chân xuống bãi cát vàng óng ả. Kaiho mặc một bộ đồ bơi hai mảnh màu tím than thể thao làm tôn lên vóc dáng khỏe khoắn, vươn hai tay hít căng lồng ngực, trầm trồ: ``Vậy ra đây là bãi biển Kaiwa ngoài đời thực trông như thế này à\dots''

Luna mặc bộ đồ bơi bèo nhún màu hồng phấn, hai tay ôm chiếc phao tròn hình con vịt, bập bẹ: ``Trông khác hẳn so với bãi bỉn ở thành phố Điện tử chúng ta á-nanora!''

Towa trong bộ bikini màu đen tuyền cá tính khẽ vuốt lại mái tóc dài, gật gù: ``Đông vui ghê! Gió mát rượi luôn này!''

Kaiho đứng cạnh tôi, đưa mắt nhìn quanh dàn idol lộng lẫy trong những bộ áo tắm rực rỡ sắc màu, bất giác bật cười khúc khích: ``Đặc biệt hơn, ngoài bờ biển này còn có cả một dàn hoa hậu \hll\ đang sải bước nữa chứ!''

Tôi lập tức liếc xéo cô nàng, hạ giọng cảnh cáo: ``Này này, tém cái miệng lại giùm anh một chút đi. Ở đây chẳng ai quan tâm đâu, nhưng nói năng hớ hênh lộ ra thì đừng trách.''

Towa cũng hoảng hốt chạy tới huých tay Kaiho: ``Đúng đấy Kaiho-senpai! Bộ chị muốn cả bãi biển nhận ra chúng ta rồi bu kín lại chắc!?''

Kaiho chỉ biết gãi đầu cười trừ. Nhưng trước khi cả đoàn kịp ùa xuống biển, việc quan trọng đầu tiên theo quy định nghiêm ngặt của mẹ Kaien là toàn bộ thành viên phải thực hiện bài khởi động giãn cơ toàn thân trên bờ cát.

Tôi không biết bơi nên chỉ đóng vai trò người giám sát kiêm chân giữ đồ, thảnh thơi đứng lội nước bì bõm ở mép cát nông.

Ryuuhou mặc bộ đồ bơi thể thao màu đen viền cam nổi bật, đôi chân trần tung tăng giẫm lên lớp cát ướt mịn. Cô nàng hăm hở vặn hông, gập người xoay khớp vô cùng thuần thục nhờ những năm tháng luyện tập vũ đạo: ``Anh Kuon! Lát nữa bơi xong ra bãi đá câu cá với em nhé! Bác Toei đã chuẩn bị sẵn cần câu xịn ở cốp xe rồi đấy!''

Trong khi đó, Suzuno mặc bộ đồ bơi liền thân màu tím nhã nhặn, bên ngoài khoác một chiếc áo ren trắng mỏng manh. Em cẩn thận quỳ gối trên tấm bạt trải cát, kiên nhẫn thoa đều từng lớp kem chống nắng lên đôi vai nhỏ và hai cánh tay của Naomi. Bé út Naomi diện bộ đồ bơi một mảnh màu vàng chanh in hình hoa hướng dương, đội chiếc mũ bơi nhỏ xíu, đôi mắt to tròn chăm chú nhìn theo từng đợt bọt sóng trắng xóa liếm vào mép cát dưới chân: ``Mát quá chị Suzuno ơi\dots\ Cát dưới chân em nó cứ nhồn nhột thế nào ấy!''

Suzuno mỉm cười dịu dàng, vuốt lại sợi tóc mai cho em gái: ``Đấy là cát đang mát-xa cho em đấy. Thoa kem xong rồi, hai chị em mình ra mép nước nhặt vỏ ốc trước nhé.''

Thế nhưng, giữa không khí yên bình ấy, cô nàng Thiên sứ Kanata bỗng khựng lại giữa bài tập ép dẻo.

Đôi mắt cô bất giác liếc ngang liếc dọc, quét một vòng qua vòng một đầy đặn, căng tràn sức sống của Kson, rồi trượt qua sự đẫy đà của Watame và cả thân hình cân đối của Towa. Chưa dừng lại ở đó, ánh mắt cô nàng còn liếc sang hai cô gái lớn nhà Hikawa: Kaiho sở hữu vóc dáng thể thao nở nang, quyến rũ chết người, còn Suzuno dù dáng người thanh mảnh nhưng vòng một lại đẫy đà, tròn trịa chẳng kém cạnh bất kỳ ai. Để rồi khi cúi đầu nhìn xuống phần ngực áo bơi của chính mình, Kanata bàng hoàng nhận ra một sự thật cay đắng: trước ngực cô nàng phẳng lì như một dải đồng bằng không một gợn sóng!

Kson đang xoay hông, nhận thấy ánh mắt thất thần của cô bạn liền cất tiếng hỏi: ``Ủa, bồ đang nhìn cái gì chăm chú thế, Kanatan?''

Kanata giật mình đánh thót, hai má đỏ bừng ấp úng: ``K-Không có gì đâu mà\dots''

Thế nhưng, cô nàng Cừu đứng bên cạnh lập tức nở một nụ cười gian xảo như nhìn thấu tim gan đối phương, chớp chớp mắt trêu chọc: ``A\~{}! Có phải bà đang nhìn những trái dưa căng mọng của bọn tôi không nè\~{}?''

Câu nói trúng tim đen ấy chẳng khác nào một tiếng sét đánh ngang tai. Toàn bộ thế giới xung quanh Kanata như sụp đổ tan tành. Cô nàng đứng chết trân, gương mặt tái mét.

Thế nhưng, kẻ chịu tổn thương sâu sắc đâu chỉ có mình Kanata!

Đứng ngay cạnh đó, Ryuuhou—người vốn sở hữu thân hình tomboy mảnh khảnh và vòng một phẳng lặng hệt như một chiếc sân bay quân sự—vừa nghe thấy Watame mở mồm khoe ``trái dưa'' thì cả người lập tức giật bắn lên. Câu nói ấy chẳng khác nào một nhát dao xỉa xói trực diện vào nỗi tự ái ngút trời của cô em thứ năm!

Ryuuhou lập tức chống nanh, lông mày dựng ngược, trừng mắt gắt toáng lên chát chúa: ``Này bà Cừu kia! Phẳng thì sao hả!? Phẳng thì đắc tội với dòng họ nhà bà chắc!? Người ta tập nhảy breakdance với hoạt động underground, ngực phẳng nó mới gọn gàng, thanh thoát và cơ động nghe chưa!!''

Tôi và Kaiho đứng gần đó không khỏi toát mồ hôi hột, đồng thanh cất tiếng: ``N-Này bình tĩnh nào hai đứa\dots''

Luna thấy vậy liền phụng phịu mắng Watame: ``Này Watame! Sao bà có thể chêu bà í ác như thế đực chứ-nanora!?''

Nhưng nàng Cừu tinh quái chỉ nghiêng đầu giả ngơ vô tội: ``Đâu có, tôi có nói gì Ryuuhou-chan đâu nà\~{}!''

Không chịu nổi cú sốc kép, Kanata lập tức ngồi thụp hẳn xuống bờ cát, hai tay ôm chặt lấy mặt, chỉ quay tấm lưng trần cô độc về phía mọi người rồi nức nở giả khóc: ``Tại sao bà lại nói những lời tàn nhẫn như thế với tôi chứ!? Thật là quá đáng mà!!''

Tôi thở dài thườn thượt, liếc nhìn nàng Cừu: ``Em làm con bé khóc thật rồi kìa, Wata-chan.''

Towa cũng khoanh tay, nghiêm giọng: ``Mau nói lời xin lỗi Kanata đi!''

Watame thấy tình hình hơi quá đà liền lật đật ngồi xổm xuống bên cạnh, đưa tay vỗ nhẹ lưng cô bạn: ``Rồi rồi, cho tôi xin lỗi mà! Tôi chỉ đùa một chút thôi, không có ý đó đâu\dots''

Thế nhưng, nàng Thiên sứ cứng đầu vẫn không thèm quay mặt lại, hậm hực cất giọng nghẹn ngào: ``Như thế vẫn chưa đủ!''

Kson nhướn mày: ``Chứ bồ muốn sao nữa?''

Nàng Thiên sứ giậm chân xuống cát: ``Bà ấy phải dõng dạc nói rằng ngực em có `boing-boing' cơ!!''

Cả đám chúng tôi đồng thanh há hốc mồm: ``\dots Hả?''

Kanata tiếp tục mè nheo: ``Nói đi! Nếu không em sẽ quay lưng về phía mọi người suốt cả chuyến đi biển này luôn cho xem! Mọi người muốn nhìn thấy một Kanata nhí nhảnh đáng yêu hay là một Kanata đau khổ cả ngày hả!?''

Watame còn đang bối rối chưa biết phải khen trái với lương tâm thế nào thì cô nàng Luna bỗng hồn nhiên bồi thêm một cú chí mạng: ``Nhưng mà bà chước xau như một mà-nanora!''

Ngòi nổ lập tức bốc cháy ngùn ngụt. Kanata bật dậy như lò xo, trừng trừng mắt nhìn nàng công chúa kẹo ngọt: ``Bà vừa nói cái gì đấy!?''

Luna vênh mặt khiêu khích: ``Tôi bít bà ăn vạ rồi-nanora! Khum cần phải ziễn nữa đâu-nanora!''

Kanata bẻ khớp ngón tay nghe \textit{rắc rắc}, sát khí đằng đằng tỏa ra chuẩn bị lao vào bóp nát chiếc phao vịt của Luna. Cả tôi, Kson và Towa hoảng hồn phải lao vào can ngăn hết sức bình sinh, giữ chặt lấy hai cánh tay sở hữu lực nắm năm mươi ký kinh hoàng của nàng Thiên sứ.

Trong lúc hỗn loạn, Ryuuhou hậm hực sải bước tới cạnh cô nàng Thiên sứ, vỗ mạnh vào vai người bạn cùng cảnh ngộ: ``Đừng thèm nghe mấy lời kích bác của bọn họ, Kanata-san! Chúng ta phẳng nhưng chúng ta có phong cách! Sau này ai dám chê ngực chị, cứ bảo em, em đập nát đàn guitar vào đầu đứa đó luôn!''

Kanata ngẩng phắt đầu lên, nhìn sang cô nàng tomboy cá tính nhà Hikawa rồi hai mắt bỗng sáng bừng lên như tìm thấy tri kỷ ngàn năm: ``E-Em nói thật sao!? Cuối cùng mình cũng tìm được đồng minh cùng chung chiến tuyến rồi!!''

Nàng Thiên sứ lập tức nhào tới ôm chầm lấy Ryuuhou khóc ròng trong niềm xúc động tột cùng. Ryuuhou cũng ôm lại cô bạn, vỗ lưng an ủi: ``Chuẩn luôn! Chiến hữu phẳng muôn năm!''

Suzuno đứng cạnh Naomi, khẽ mỉm cười đưa khăn lau mặt cho Kanata: ``Kanata-san và Ryuuhou đều rất dễ thương và năng động mà, đừng buồn nữa nhé.''

Chloe đứng ngoài rìa, ngậm que kem cắn dở, chứng kiến toàn bộ màn hài kịch từ đầu đến cuối chỉ biết lắc đầu thở dài: ``Poe\dots\ Ở đây còn hỗn loạn và kỳ quặc hơn cả văn phòng holoX của tụi tôi nữa cơ đấy\dots''

Thế nhưng, màn tranh cãi ồn ào ấy mới chỉ là khúc dạo đầu êm ả. Cơn ác mộng thực sự chỉ bắt đầu khi cả đoàn cùng nhau nhảy xuống làn nước biển cuộn trào.

\ct

``\textbf{Lại có một đợt xóng kìa-nanora!!}'' tiếng hét thất thanh của Luna vang lên xé toạc tiếng gió rít.

``\textbf{Bám chắc vào phao đi mọi người ơi!!}'' Kson gào lớn.

``\textbf{Nhưng mà sóng ở đây mạnh quá!!}'' Towa bám chặt lấy sợi dây viền của chiếc phao tròn khổng lồ mà tôi vừa thuê ở bãi trượt nước.

Một con sóng bạc đầu cao ngút ngàn đổ ập xuống như một bức tường nước khổng lồ.

*RẦM!*

Bọt nước trắng xóa bắn tung tóe mù mịt. Bảy con người gồm Ryuuhou, Kson, Kanata, Towa, Watame, Luna và cả Kaiho đang bám trên chiếc phao lớn lập tức bị con sóng cuốn phăng đi.

Nàng Công chúa xứ Kệo bị sóng tạt thẳng vào mặt, lảo đảo tuột tay: ``Xóng kìa-- Ụp!''

``Luna! Đừng có buông tay!'' Towa vội vã với tay chộp lấy cổ áo cô bạn.

``Bám chặt vào lưng tôi này!'' Kson rống lên giữa làn nước ngập ngụa.

Ryuuhou thì cười toe toét, đập nước ầm ĩ: ``Sóng tới nữa kìa! Vui quá đi mất!!''

Tôi đứng ở vùng nước nông ngang thắt lưng, hai tay chống hông nhìn đám người đang vật lộn la hét đầy thích thú mà không khỏi vừa buồn cười vừa bất lực. Trông họ lúc này chẳng khác nào một bầy vịt con nhốn nháo hơn là những thần tượng nổi tiếng của một tập đoàn giải trí hàng đầu.

Cách đó không xa ở vùng nước êm ả sát bờ, Suzuno đang cẩn thận dắt tay bé út Naomi đi dạo dọc theo dải bọt sóng trắng. Naomi đeo chiếc phao bơi mini hình chú cá heo màu xanh, hai chân nhỏ xíu bì bõm giẫm lên làn nước trong vắt, chốc chốc lại cúi xuống nhặt một chiếc vỏ sò hoa ngũ sắc rồi giơ lên khoe rối rít: ``Suzuno-oneechan nhìn này! Vỏ sò lấp lánh như ngọc trai luôn này chị!''

Suzuno mỉm cười hiền dịu, cúi người đỡ lấy vỏ sò bỏ vào chiếc xô nhựa nhỏ, vuốt lại mái tóc ướt nước muối của em gái: ``Đẹp lắm Naomi-chan. Lát nữa về phòng chị em mình rửa sạch rồi xâu thành chuỗi vòng lưu niệm nhé.''

Thế rồi, hai chị em cùng ngước nhìn ra ngoài xa. Ở đó, Chloe đang trình diễn bản lĩnh vượt trội của loài sinh vật biển sâu. Rời xa mặt đất, con cá voi sát thủ như được trở về với vương quốc thực sự của mình. Cô nàng bơi lội thoăn thoắt, lặn ngụp dưới những ngọn sóng sâu hoắm rồi bất thình lình phóng vút lên khỏi mặt nước xoay một vòng điệu nghệ, bọt nước bắn tung tóe.

Ryuuhou đang ôm chặt mép phao tròn cùng nhóm Gen 4, thấy cô nàng làm xiếc dưới nước liền nổi máu ganh đua, buông một tay khỏi phao đập nước bắn tung tóe về phía cô nàng cá voi: ``Khá đấy đồ cá voi sát thủ! Nhưng đừng hòng qua mặt được dân bơi lội chuyên nghiệp như ta nhé!''

Chloe trồi đầu lên khỏi mặt nước, lè lưỡi làm mặt quỷ: ``Poe poe! Có giỏi thì bơi đuổi theo tôi này, đồ nhóc con tomboy!''

Naomi đứng trên bờ thấy cảnh hai người đấu khẩu té nước chí tử liền thích thú vỗ tay nhảy cẫng lên: ``Ryuuhou-oneechan cố lên! Chloe-oneechan cố lên!!''

Thế nhưng, trò đùa của số phận chưa bao giờ dừng lại ở đó.

Khi tôi vừa quay đầu lại định gọi mọi người cẩn thận dạt vào bờ thì một con sóng quái vật - cao hơn gấp đôi những đợt sóng trước—từ ngoài khơi xa bất ngờ trườn tới với tốc độ kinh hoàng.

``\textbf{Cẩn thận kìa Onii!!}'' tiếng Ryuuhou hét lên trong gió.

Tôi còn chưa kịp định thần xem chuyện gì đang xảy ra thì cả một khối nước khổng lồ nặng hàng tấn đã giáng thẳng xuống đầu tôi.

*RẦM!!*

Cả thế giới xung quanh tôi chao đảo dữ dội. Tôi bị dòng nước cuốn nhào lộn tùng phèo dưới đáy cát, mắt cay xè, miệng ngập ngụa vị mặn chát của muối biển. Tôi cố dùng hết sức bình sinh đạp chân vào nền cát để ngoi lên khỏi mặt nước.

``\textit{Phù\dots! Khục khục\dots!}''

Tôi ho sặc sụa, vừa trồi lên được thì trước mắt bỗng tối om, một vật thể mềm mại và ướt sũng nào đó trùm kín mít lấy toàn bộ khuôn mặt tôi.

Trong cơn hoảng loạn vì không thở được, tôi vội vã đưa hai tay lên giật phăng thứ đang chắn ngang tầm nhìn của mình ra. Thế nhưng trong lúc quờ quạng, bàn tay phải của tôi vô tình quẹt trúng một mảng da thịt mềm mại ở phía dưới mặt nước, và ngay lập tức, một tiếng thét chói tai quen thuộc vang lên xé toạc cả góc biển:

``\textbf{KYAAAHH!!}''

Kèm theo đó là hàng loạt tiếng la hét thất thanh của các cô gái xung quanh:

``Chị ổn chứ, Kaiho-senpai!?''

``K-Không sao, chị không--''

``Kaiho-senpai, đừng có đứng dậy vội!!''

``Ể!? \textbf{Á á á! Kyaaah!!}''

``Có chuyện gì thế Kaiho\dots\ Ối chà!! Cả tôi nữa này!!''

Tôi vội vàng dụi sạch nước biển cay xè bám trên khóe mắt để nhìn rõ xung quanh. Và khi cảnh tượng trước mắt hiện ra rõ mồn một, toàn bộ não bộ của tôi chính thức bị đình trệ hoàn toàn.

Trên tay trái của tôi lúc này\dots\ là một chiếc áo ngực bơi màu trắng viền cam rực rỡ của Kson.

Còn trên các ngón tay bên phải của tôi\dots\ là một dải dây quần bơi màu tím than sẫm, chính là chiếc quần bơi thể thao của Kaiho!

Ngay trước mặt tôi giữa làn nước trong vắt, cô nàng Tổng trưởng đang dùng cả hai tay ôm chặt lấy vòng một trần trụi của mình, mặt đỏ bừng như quả ớt chín. Còn cô em song sinh thì đang ngồi thụp hẳn xuống đáy nước, hai tay che kín phần dưới cơ thể, cả gương mặt lẫn vành tai bốc khói nghi ngút vì ngượng ngùng cực độ!

Sóng biển dữ dội đã đánh tụt sạch đồ bơi của hai người họ, và bằng một sự sắp đặt trớ trêu đến quỷ khốc thần sầu, toàn bộ những mảnh vải ấy lại nằm gọn lỏn trong hai bàn tay tôi!

Chưa kịp để tôi mở miệng thanh minh lấy nửa lời, Watame nổi lên ngay bên cạnh, nở một nụ cười nham hiểm tột cùng rồi cất giọng ngân nga đầy châm chọc: ``Kuon-senpai\dots\ Anh có vẻ khá là dâm đãng nghen\~{}''

Kaiho và Kson xấu hổ đến mức đồng thanh gào lên trong làn nước: ``\textbf{Tại con sóng lớn quá đánh tụt chứ bộ!!}''

Watame lại tiếp tục chêm dầu vào lửa: ``Uầy\~{} Sóng đánh tụt cả quần áo bơi là có thật kìa\~{} Chuyến này Kuon-senpai tha hồ bổ mắt nhé\~{}''

Đứng cách đó một đoạn, Chloe vừa trồi lên khỏi mặt nước, nhìn thấy cảnh tượng chấn động địa cầu liền trợn tròn đôi mắt đỏ rực, luống cuống thò tay vào túi áo phao: ``P-Poe?! Khoan đã! Chờ một chút! Cảnh tượng phạm tội đắt giá thế này phải chụp ảnh lại làm bằng chứng tống tiền--''

Chưa kịp dứt câu, Ryuuhou đứng ngay sau lưng đã thẳng tay ấn dúi đầu con cá voi sát thủ chìm nghỉm xuống đáy nước: ``Im mồm đi đồ cá voi nhiều chuyện!''

Trước những ánh mắt bốc hỏa như muốn thiêu rụi mình ra tro của Kaiho và Kson, tôi chỉ biết lắp bắp run rẩy: ``C-Cho anh xin phép!!''

Nói rồi, tôi vội vàng lặn ngụp một mạch rồi co giò phóng thục mạng lên bờ cát. Nhưng khốn nỗi, trong lúc hoảng loạn tột cùng, hai bàn tay tôi vẫn nắm chặt chiếc áo ngực của Kson và chiếc quần bơi của Kaiho mà quên béng mất chưa trả lại cho chủ nhân của chúng!

Thấy ``tên trộm'' cầm theo y phục bỏ chạy lên bờ, Kson và Kaiho vừa xấu hổ vừa tức điên người.

Kson che ngực lội nước đuổi theo: ``Này tên Kuon-paisen kia! Đứng lại! Trả áo cho tôi mau!!''

Kanata và Towa hoảng hồn bơi chắn phía trước cô nàng: ``Kson! Ngực kìa! Coi chừng lộ hết bây giờ! Để bọn tôi che cho!''

Luna cũng quẫy nước bơi theo: ``Che cho bà í với-nanora!''

Trong khi đó, Kaiho đang lúp xúp ngâm mình dưới nước, vừa ngượng chín mặt vừa gào thét trong tuyệt vọng: ``Ê\dots! Không có đứa nào thèm che chắn cho chị à!?''

Watame bơi lững lờ bên cạnh, cười hí hí: ``Hehe\dots''

Kaiho đỏ mặt gắt lên: ``Đừng có đứng đấy mà cười nữa! Mau che chắn cho chị rồi cùng đuổi theo giật lại đồ từ tay ông anh ngốc kia mau lên!!''

Ở mép nước nông gần bờ cát, Suzuno và Naomi đang nhặt vỏ sò thì nghe tiếng thét thất thanh liền ngẩng đầu lên.

Vừa nhìn thấy cảnh tượng tôi cả hai tay đang vung vẩy chiếc áo ngực bơi cùng chiếc quần lót ướt sũng phóng thục mạng lên bờ, hai má Suzuno lập tức đỏ ửng như gấc chín. Em giật thót mình, không một giây chần chừ, vội vã đưa cả hai bàn tay che kín mít đôi mắt trong veo của bé út Naomi lại: ``Naomi-chan, không được nhìn! Cảnh này trẻ con tuyệt đối không được nhìn thấy đâu em!''

Bản thân em cũng ngượng ngùng quay ngoắt mặt sang hướng khác, lắp bắp không nên lời: ``Onii-sama\dots\ Sao anh lại\dots\ cầm mấy thứ đó trên tay chạy lông nhông thế kia chứ\dots!''

Naomi bị bịt mắt bất ngờ, chỉ biết quờ quạng hai tay nhỏ xíu lên mu bàn tay của chị, giọng ngây thơ ngơ ngác: ``Ơ\dots\ Suzuno-oneechan ơi, có chuyện gì thế ạ? Em nghe tiếng Onii-chan hét to lắm mà\dots\ Sao tự nhiên trời tối thui vậy chị?''

Thế là trên bờ cát vàng rực rỡ của bãi biển Kaiwa chiều hôm ấy, một cuộc rượt đuổi hỗn loạn chưa từng có trong lịch sử đã diễn ra giữa chàng Tổng quản lý bất đắc dĩ và một bầy thần tượng đang la ó vang trời.

Một ngày đầu tiên tại biển khơi ngập tràn những tiếng cười, những pha tai nạn thót tim và sự gắn kết không thể nào quên.

\ct

Tám giờ tối.

Sau một buổi chiều vận động tơi bời dưới làn sóng biển và bữa tiệc tối no nê với các món hải sản đậm đà ở nhà ăn trung tâm, cả gia đình chúng tôi kéo nhau về lại tầng ba để nghỉ ngơi.

Vừa bước đến hành lang trước cửa phòng 420B, tôi đang định mở cửa vào trong đánh một giấc thì tiếng gọi oang oang quen thuộc của Kson từ phòng đối diện đã cất lên: ``Ê Kuon! Kaiho! Hai anh chị có muốn sang đây làm vài ván cờ tỷ phú Mon*p*ly không!?''

Kaiho vừa lau tóc bằng khăn bông, nghe thấy hai chữ cờ tỷ phú thì đôi mắt màu lam lập tức sáng bừng lên như mèo thấy mỡ. Vốn mang bản tính hiếu thắng và ưa thích những trò chơi tính toán cơ học, cô nàng gật đầu cái rụp không chút đắn đo: ``Chơi chứ! Em đang muốn phục thù ván trước đây!''

Đứng ngay cạnh bên, Ryuuhou cũng nhanh nhảu huých nhẹ khuỷu tay vào hông tôi, nháy mắt tinh quái: ``Đi thôi Onii! Có em tham gia cùng phe thì đảm bảo cái khoản thương lượng, mặc cả đất đai với đám Gen 4 là anh em mình chấp cả bản đồ luôn! Giới underground chưa bao giờ ngán chuyện trả giá nhé!''

Tôi khẽ mỉm cười. Lâu lắm rồi tôi mới lại có dịp ngồi vào bàn cờ tỷ phú kể từ chuyến đi dã ngoại của lớp thời cấp ba. Với bộ ba gồm một kẻ có tư duy chiến lược logic như tôi, một chuyên gia tính toán vật lý không gian như Kaiho và một tay ngoại giao lém lỉnh, am hiểu tâm lý xã hội như Ryuuhou, liên minh nhà Hikawa quả thực là một đối thủ vô cùng đáng gờm.

Căn phòng 404B nhanh chóng trở thành một sới cờ rực lửa.

Bàn cờ tỷ phú được trải ngay ngắn trên chiếc bàn trà tròn giữa phòng. Sáu người tham gia tranh tài gồm: phe nhà Hikawa (tôi, Kaiho và Ryuuhou), phe Thiên sứ - Ác ma (Kanata và Towa) và phe Đùi Cừu - Công chúa (Watame và Luna). Nàng cựu Tổng trưởng Kson đóng vai trò quản trò kiêm thủ quỹ ngân hàng, ngồi bắt chéo chân phát tiền và thẻ cơ hội.

Ở góc giường phía trong, không gian lại yên bình đến lạ kỳ. Bé út Naomi sau một ngày bơi lội mệt nhoài đã ngủ say như một chú cún con dưới lớp chăn ấm. Suzuno ngồi bên cạnh khẽ lật từng trang sách hóa sinh, chốc chốc lại mỉm cười ngước nhìn về phía bàn cờ nhộn nhịp. Chloe cũng kéo ghế ngồi cạnh Suzuno, tay cầm cốc trà đào chăm chú quan sát trận đấu. Dù không hiểu hết luật chơi, cô nàng sát thủ vẫn kiên trì dán chặt mắt vào từng đường đi nước bước của tôi để thực hiện trọn vẹn ``nhiệm vụ theo dõi mục tiêu'' được giao.

``Mười ô! Một, hai, ba\dots'' Watame vừa đếm bước chân của quân cờ vừa lẩm bẩm.

Kson đập bàn hô to: ``\dots Chín, mười! Rơi trúng vào ô đất của Kanatan rồi nhé! Nộp phạt ba nghìn đồng mau lên!''

Nàng Cừu ôm đầu rên rỉ thảm thiết: ``Trời đất ơi\dots\ Tận ba nghìn lận! Mới đầu trận mà đã bay đứt nửa gia tài rồi này\~{}!''

Luna ngồi vắt vẻo bên cạnh, huơ huơ hai con xúc xắc trên tay: ``Thui, để tui gỡ gạc lại cho-nanora!''

Ván cờ trôi qua được khoảng ba mươi phút và thế trận bắt đầu phân hóa rõ rệt. Nhờ vào sự phối hợp ăn ý không tưởng của bộ ba anh em chúng tôi, liên minh nhà Hikawa đã càn quét và thâu tóm được gần hai phần ba diện tích bàn cờ, đặc biệt là những khu đất vàng đắt đỏ bậc nhất nằm sát vạch xuất phát.

Ở giai đoạn đầu, khi mọi người còn đang ngập ngừng thăm dò, Kaiho đã nhanh chóng phát huy thế mạnh độc nhất vô nhị của mình. Dù cô nàng vốn cực kỳ dị ứng với những công thức đại số trừu tượng khô khan trên sách vở, nhưng khi chuyển hóa thành bài toán không gian và quỹ đạo chuyển động thực tế, trực giác hình học và cơ học của Kaiho lại nhạy bén đến mức đáng sợ. Thay vì mua đất theo cảm tính, cô nàng nhìn toàn bộ bàn cờ 40 ô như một hệ tọa độ khép kín. Cô nhẩm tính khoảng cách bước nhảy phổ biến nhất từ ô Ra Tù dựa trên sự kết hợp va chạm của hai khối lập phương, từ đó vạch ra chính xác cung đoạn từ 6 đến 8 bước—nơi có mật độ va chạm cao nhất. Bằng sự quyết đoán lạnh lùng của một tay thực chiến, Kaiho không ngần ngại vung tiền thế chấp những mảnh đất nhỏ lẻ phân tán để dồn vốn ôm trọn hai dải đất màu cam và đỏ, biến cung đường ngay sau cổng nhà tù thành một bãi mìn đón lõng hoàn hảo.

Khi ván cờ bước vào giai đoạn thương lượng phức tạp giữa các bên, sàn đấu chính thức thuộc về Ryuuhou. Nhận thấy Kanata đang nắm giữ một mảnh đất màu cam then chốt nhưng lại thiếu tiền mặt để nâng cấp nhà, Ryuuhou liền phát huy toàn bộ bản lĩnh ngoại giao dẻo miệng của một trưởng ban nhạc underground. Cô em thứ năm vừa cười toe toét vừa nắm tay Kanata rót mật vào tai: ``Kanata-san, chị giữ mảnh đất cam này một mình thì khác nào ôm cục nợ không sinh lời đâu nè! Chi bằng nhượng lại cho bọn em, bọn em sẽ bù cho chị một khoản tiền tươi rói kèm theo quyền miễn phí thuê nhà hai lần khi dẫm vào khu đất màu hồng của bọn em nhé? Hợp đồng thế này là chị vừa có tiền xây nhà bên kia, lại vừa có `bảo hiểm nhân thọ' dọc đường đi, quá hời luôn còn gì!''

Sự khéo léo, nắm bắt tâm lý cực nhạy của Ryuuhou khiến Kanata gật đầu lia lịa ký giấy nhượng đất trong sự bất lực của đồng đội Towa: ``Kanata ngốc ơi! Bà bị con bé lừa bán sang biên giới rồi kìa!!''

Về phần tôi, với cái đầu lạnh của một Tổng quản lý chuyên xử lý các bài toán tài chính hóc búa, tôi đóng vai trò nhạc trưởng điều phối dòng tiền. Tôi tính toán điểm hòa vốn và kiên quyết chỉ nâng cấp tối đa lên ba căn nhà gỗ ở mỗi ô đất thay vì vội vã xây khách sạn lớn ngay từ đầu, vừa giữ lại lượng tiền mặt an toàn để đối phó với thẻ cơ hội, vừa gây sức ép triệt để lên nguồn vốn lưu động của hai đội còn lại.

Luna nhẩm tính số bước rồi run rẩy thốt lên: ``Nể tình xíu đi-nanora\dots\ Rơi vào ô của Kuon-senpai với hai bà chị nhà Hikawa là xạt nghịp lun á-nanora\dots''

Kson gõ nhẹ cây bút xuống bàn: ``Đến lượt phe nhà Hikawa! Ai đổ xúc xắc đây?''

Kaiho xắn tay áo lên, cười khì: ``Để em đổ cho! Tay em đang đỏ lắm!''

Hai con xúc xắc lăn lông lốc trên bàn rồi dừng lại ở mặt sáu và sáu: tổng cộng mười hai nút!

Kson hô: ``Mười hai ô! Đổ đôi thì được quyền đổ thêm một lượt nữa!''

Kaiho lại vung tay. Hai con xúc xắc lại nhảy ra hai mặt sáu! Lại mười hai nút! Cả căn phòng đồng loạt ồ lên kinh ngạc: ``Gì vậy trời!? Tay con bé này gắn nam châm đấy à!?''

Kaiho đắc chí nháy mắt với tôi: ``Thấy chưa Aniki! Em đã bảo là em may mắn mà!''

Tôi nhìn thế cờ, bình thản nhấp ngụm nước rồi nhắc khéo: ``Cái tiếp theo mà ra mười hai nữa là coi chừng vào tù đấy em.''

Kaiho nhăn mũi: ``Làm sao mà trùng hợp thế được chứ!''

Vừa dứt lời, cô nàng lại tung hai con xúc xắc lên không trung. Và như một trò đùa quái ác của định mệnh, hai khối vuông lại dừng lại ở mặt sáu và sáu! Ba lần liên tiếp đổ ra hai con sáu!

Cả đám con gái đang trầm trồ bỗng khựng lại khi Kson nở một nụ cười rạng rỡ như ánh mặt trời, chỉ tay thẳng vào mặt Kaiho: ``Chúc mừng quý cô may mắn! Ba lần đổ đôi liên tiếp, theo luật thì tống thẳng vào tù ngay và luôn!!''

Cô em song sinh chết sững người, mắt tròn xoe như hai quả trứng gà: ``Ể!? T-Tại sao chứ!? Đổ điểm cao mà lại bị phạt là thế nào!?''

Ryuuhou ôm trán thở dài: ``Luật cờ tỷ phú nó thế đấy Kaiho-nee\dots\ Ham hố cho lắm vào rồi giờ bóc lịch!''

Kanata ngồi đối diện khoái chí vỗ đùi bôm bốp: ``Haha! Ai bảo chị không thèm đọc kỹ luật làm gì! Cho chừa cái tội thâu tóm đất đai nhé!''

Kson rút chiếc còng đồ chơi huơ huơ: ``Nào, nếu muốn ra tù ngay thì phải nộp phạt năm mươi đô lệ phí bảo lãnh, hoặc dùng thẻ Ra Tù Miễn Phí! Bằng không thì phải ngồi im trong đó đổ đôi tiếp trong tối đa ba lượt!''

Kaiho bực dọc toan rút xấp tiền giấy trong tay ra trả: ``Năm mươi đô chứ gì? Nộp phạt luôn cho xong, em phải ra ngoài tiếp tục gom đất--''

Tôi nhanh như chớp đưa tay giữ chặt cổ tay cô nàng lại: ``Khoan đã. Cứ ngồi yên trong tù đi, không việc gì phải trả tiền để ra vội cả.''

Em trố mắt nhìn tôi, giãy nảy: ``Tại sao chứ anh!? Ở lại trong tù thì mình bị giam chân, không đi lại gom thêm đất được! Anh tính để hai đứa mình chịu trói trong chốn lao tù à!?''

Tôi điềm tĩnh lắc đầu, khóe môi khẽ nhếch lên một nụ cười bí hiểm: ``Cứ nghe lời anh. Theo luật chính thống, người ở trong tù vẫn có đầy đủ quyền thu tiền thuê đất, mua bán nhà và tham gia đấu giá như người bình thường. Giai đoạn này bàn cờ đầy rẫy bẫy rập, ở trong tù mới là nơi an toàn nhất.''

Ryuuhou nhìn biểu cảm tự tin của tôi liền lập tức hiểu ý, vỗ vai bà chị song sinh cười khúc khích: ``Chuẩn bài luôn! Anh ấy nói đúng đấy chị, cứ tin ở anh ấy đi!''

Kaiho tuy chưa hiểu hết nước cờ sâu xa nhưng đành ngậm ngùi thu tiền lại, phồng má ngồi im chịu trận.

Và rồi, sự thật diễn ra sau đó đã khiến toàn bộ người chơi trong phòng phải ngả mũ thán phục.

Giai đoạn cuối trận, gần như toàn bộ các ô đất trên bàn cờ đều đã mọc đầy nhà cửa với giá thuê cắt cổ. Việc quân cờ của chúng tôi an tọa trong tù đồng nghĩa với việc chúng tôi hoàn toàn miễn nhiễm với mọi rủi ro dẫm phải ô đất của đối phương! Suốt ba lượt đổ xúc xắc tiếp theo, Kaiho cố tình đổ không ra đôi để được hợp pháp ở lại trong buồng giam.

Trong khi đó, ở thế giới bên ngoài, Kanata, Towa, Watame và Luna cứ liên tục dẫm đạp lên chuỗi nhà phố màu cam và đỏ của nhà Hikawa.

Dù không phải tốn một xu tiền thuê đất nào cho đối phương, nhưng tiền phạt thuê nhà từ các đội khác lại chảy vào túi chúng tôi liên hồi như nước lũ!

Towa ôm mặt gào thét: ``\textbf{Áaaa! Lại rơi trúng ô đất của Kuon-senpai rồi!! Sao ô nào cũng là đất nhà anh ấy thế này!?}''

Watame nước mắt ngắn nước mắt dài móc hầu bao: ``Tận sáu nghìn đồng lận\dots\ Bán cả đàn cừu đi cũng không đủ trả nợ nữa rồi\dots''

Kson đếm tiền xoành xoạch, cười ha hả: ``Towa và Watame, mỗi người nộp phạt sáu nghìn đồng cho tù nhân nhà Hikawa nhé! Khà khà khà!''

Sau ba lượt, khi buộc phải nộp năm mươi đô tiền phạt để rời khỏi buồng giam theo luật định, khối tài sản trong tay phe nhà Hikawa đã phình to gấp mấy lần. Kaiho kinh ngạc đến mức há hốc mồm, quay sang nhìn tôi bằng ánh mắt cảm phục hiếm thấy, khẽ đập tay với Ryuuhou: ``Ghê thật đấy\dots\ Hóa ra ngồi tù lại là chiến thuật phòng thủ kiếm tiền đỉnh cao đến mức này à!?''

Tôi cười trừ: ``Anh đã bảo rồi mà. Trò này phải biết dùng cái đầu lạnh để tính toán thời điểm.''

Trận chiến kéo dài đến tận gần mười một giờ đêm. Dù về cuối trận, Kanata và Watame cũng có vài pha lật kèo ngoạn mục nhờ rút được thẻ cơ hội miễn trả tiền thuê, nhưng chung cuộc, liên minh ba anh em nhà Hikawa vẫn giành chiến thắng áp đảo với số tài sản khổng lồ.

Towa buông hai con xúc xắc xuống nệm, thở dài chịu thua: ``Thật luôn\dots\ Không thể tin được là ba anh em nhà này lại phối hợp ăn ý đến mức quái vật như vậy\dots''

Luna ngáp một cái thật dài, hai mắt díp lại: ``Khum thể ngờ đực anh chị chơi giỏi thịt đó-nanora\dots\ Buồn ngủ quá rùi nanora\dots''

Kaiho phấn khích quay sang đập vào vai tôi: ``Thắng đậm luôn rồi Aniki ơi! Anh thấy em đổ xúc xắc đỉnh-- Ơ? Aniki?''

Thế nhưng, đáp lại lời cô nàng chỉ là tiếng thở đều đặn. Tôi đã gối đầu lên cánh tay ngủ thiếp đi ngay trên mép bàn cờ từ lúc nào không hay. Cơn mệt mỏi sau chuyến xe đêm dài và một ngày lăn lộn ngoài bãi biển cuối cùng cũng đánh gục sức chịu đựng của tôi.

Watame mỉm cười dịu dàng: ``Kuon-senpai ngủ mất rồi. Chắc anh ấy mệt lắm rồi đấy.''

Kson vươn vai đứng dậy, bắt đầu thu dọn bàn cờ: ``Cũng gần nửa đêm rồi, hôm nay dừng ở đây thôi nhé các cô gái. Mai còn cả một ngày dài quậy phá ngoài biển nữa!''

Kaiho nhìn tôi ngủ say như chết, khóe môi khẽ cong lên thành một nụ cười ấm áp. Không nói một lời, cô nàng khẽ cúi người xuống, dễ dàng luồn hai cánh tay cơ bắp nhấc bổng cả người tôi lên vai nhẹ tênh như vác một bao tải bông gòn, khiến các cô nàng Gen 4 được một phen tròn mắt trầm trồ trước sức mạnh phi thường của cô em song sinh.

Ryuuhou cười tít mắt, mở toang cửa phòng: ``Kaiho-nee vô địch! Vác đại ca về phòng ngủ thôi nào!''

Đêm đầu tiên tại vùng biển Kaiwa khép lại trong sự bình yên, ấm áp và những nụ cười rạng rỡ của tất cả mọi người. Một ngày mới với những con sóng bạc đầu và những trò quậy phá bùng nổ đang chờ đón họ phía trước\dots

\ct

\stpt{-Jour 2-}
\stcap{Bình minh ký ức, bữa sáng muộn màng và đại chiến sóng thần}

Năm giờ ba mươi phút sáng.

Khi những tia nắng đầu tiên của ngày mới vừa ló rạng nơi chân trời, rọi những vệt sáng lấp lánh như dát vàng lên mặt biển Seika phẳng lặng, tôi đã khẽ cựa mình thức giấc.

Sau một giấc ngủ đẫy đà không mộng mị, đầu óc tôi sảng khoái lạ thường. Nhìn sang chiếc giường bên cạnh, Kaiho vẫn đang quấn tròn trong chăn như một chiếc kén tằm khổng lồ, tiếng thở đều đặn và sâu hoắm cho thấy cô em song sinh còn lâu mới chịu rời giường. Mẹ và bà ngoại thì đã thức dậy từ sớm để ra khuôn viên khách sạn tập dưỡng sinh cùng các cô bác trong đoàn. Suzuno và Naomi vẫn đang say giấc nồng, đôi má hai đứa ửng hồng dưới ánh mai dịu nhẹ.

Tôi rón rén xỏ đôi dép quai hậu, nhẹ nhàng khép cửa phòng rồi bước xuống sảnh.

\ct

Gió sớm thổi lồng lộng, mang theo vị mằn mặn và tươi mát của đại dương xua tan đi chút ngái ngủ cuối cùng. Tôi bước thong thả trên bãi cát ướt mịn, lắng nghe tiếng sóng vỗ rì rào êm ả dưới chân. Thế nhưng, khi vừa đi được một đoạn ngắn dọc theo mép nước, tôi bất chợt nhìn thấy một bóng hình nhỏ nhắn quen thuộc đang ngồi bó gối một mình trên tảng đá nhô ra biển.

Sakamata Chloe.

Cô nàng mặc một chiếc áo phông xám dài tay trùm qua đầu gối, mái tóc xám để xõa bay lòa xòa trong gió sớm, hai bàn chân trần ngâm dưới làn nước trong vắt, ánh mắt đỏ thẫm lơ đãng nhìn về phía đường chân trời xa xôi.

Tôi bước lại gần, tiếng cát lạo xạo dưới chân làm Chloe khẽ giật mình quay đầu lại.

``Ủa? Kuon-kun\dots?'' đôi mắt cô nàng mở to vẻ ngạc nhiên: ``Sao ông dậy sớm thế? Tưởng tối qua bị vác về phòng ngủ say như chết rồi cơ mà?''

Tôi phì cười, bước lên tảng đá rồi ngồi xuống cạnh cô bạn: ``Ngủ từ chín giờ tối thì giờ này không dậy mới là lạ đấy. Còn bà? Không ngủ nướng như mọi khi à?''

Chloe khẽ bĩu môi, dùng ngón chân khuấy nhẹ mặt nước tạo thành những vòng sóng lăn tăn: ``Người ta là sinh vật biển sâu mà! Hít thấy mùi biển sớm là tự nhiên đồng hồ sinh học nó bật dậy thôi chứ bộ\dots''

Khoảng lặng giữa hai chúng tôi kéo dài vài giây, chỉ còn lại tiếng gió sớm và tiếng hải âu chao lượn trên cao.

Đột nhiên, Chloe nhìn xuống hai bàn chân đang ngâm nước của mình, cất giọng thì thầm: ``Này Kuon-kun\dots\ Đã mười sáu năm rồi nhỉ.''

Tôi khẽ nghiêng đầu nhìn cô nàng: ``Ừ. Mười sáu năm kể từ mùa hè ở làng chài ấy.''

Chloe khẽ nhếch mép, một nụ cười hiền lành và hoài niệm hiếm thấy nở trên môi cô: ``Hồi đó hai đứa mình ngố thật đấy. Ngày nào cũng rủ nhau chạy dọc theo cái bờ đê chắn sóng, cào ngao đến trầy cả ngón tay, rồi lại thi nhau bắt mấy con còng gió chui vào lỗ cát. Lúc nào áo quần cũng dính đầy bùn đất với mùi tanh của cá, thế mà chẳng đứa nào thấy phiền.''

``Thế mà bây giờ bà lại biến thành con cá voi lười tắm trứ danh của công ty,'' tôi châm chọc.

Chloe đỏ bừng mặt, huých mạnh cùi chỏ vào sườn tôi: ``Poe! Đã bảo chuyện tắm rửa là do công việc bận rộn chứ bộ! Đừng có nhắc lại nữa chứ!''

Cả hai chúng tôi cùng bật cười thành tiếng. Tiếng cười tan vào làn gió sớm bao la. Chloe ngước nhìn lên vòm trời xanh ngắt, ánh mắt ánh lên một niềm xúc động chân thành: ``Nhưng mà\dots\ được gặp lại ông, rồi được đi biển cùng gia đình ông như thế này, tôi thấy vui lắm. Hồi đó khi ông đột ngột chuyển nhà đi theo gia đình, tôi đã nghĩ chắc cả đời này sẽ chẳng bao giờ được gặp lại người bạn thuở nhỏ nữa\dots''

Tôi đưa tay vỗ nhẹ lên vai cô nàng sát thủ: ``Duyên phận cả đấy. Dù là mười sáu năm hay bao nhiêu năm đi nữa, ký ức tuổi thơ ở làng chài ấy sẽ không bao giờ mất đi đâu. Cứ tận hưởng kỳ nghỉ này đi, đồ cá voi lười biếng.''

Chloe mỉm cười, đôi mắt đỏ rực long lanh dưới ánh nắng sớm. Một khởi đầu ngày mới thanh bình và ấm áp đến lạ kỳ.

\ct

Tám giờ bốn mươi phút sáng.

Sau khi đi dạo biển về, tôi cùng Suzuno, Ryuuhou và Naomi kéo nhau xuống nhà ăn lớn ở tầng một để dùng điểm tâm.

Bữa sáng tự chọn của khách sạn diễn ra từ sáu giờ ba mươi đến tám giờ ba mươi sáng, với đầy đủ các món nóng hổi như phở bò, cháo hải sản, bánh mì ốp la và hoa quả tráng miệng. Cả đoàn công ty mẹ lẫn hội chị em Gen 4 đều đã dùng bữa no nê, vừa cười nói rôm rả vừa chuẩn bị đồ đạc cho buổi dã ngoại chiều.

Thế nhưng, có một nhân vật đã hoàn toàn vắng mặt trong bức tranh tươi đẹp ấy.

Mãi đến gần chín giờ kém mười lăm, khi nhân viên nhà bếp đã thu dọn sạch bách các quầy đồ ăn và tắt đèn khu buffet, cánh cửa nhà ăn mới bật mở.

Kaiho uể oải bước vào với mái tóc ngắn bù xù như tổ quạ, một tay vừa dụi mắt vừa ngáp một cái thật dài, người mặc bộ quần áo ngủ rộng thùng thình: ``Oaaa\dots\ Đói bụng quá đi mất\dots\ Bữa sáng còn gì không, Aniki--''

Lời nói của cô nàng nghẹn ứ lại nơi cuống họng khi nhìn thấy toàn bộ các khay đồ ăn trên bàn đều trống trơn, sạch bóng không còn lấy một cọng hành!

Cô nhân viên phục vụ đang lau bàn thấy vậy liền ái ngại nhìn Kaiho: ``Xin lỗi chị nhé, nhà ăn của khách sạn đóng cửa phục vụ buffet sáng lúc tám giờ ba mươi rồi ạ. Bây giờ bếp bắt đầu chuẩn bị thực đơn cho bữa trưa rồi.''

Gương mặt ngái ngủ của Kaiho bỗng chốc cứng đờ. Đôi mắt màu lam mở to hết cỡ, sụp đổ hoàn toàn trước thực tại tàn khốc: ``H-Hết giờ rồi á!? Một bát cháo cũng không còn sao ạ!?''

Ryuuhou ngồi ở bàn trà gần đó đang nhâm nhi ly nước cam, thấy cảnh tượng bà chị song sinh bị bỏ đói liền đập bàn cười sặc sụa: ``Hahaha! Đáng đời chị chưa! Ai bảo tối qua hiếu thắng cho lắm vào, vác Onii về phòng xong rồi thức khuya bấm điện thoại, giờ ngủ nướng tới tận chín giờ sáng! Khách sạn chứ có phải quán ăn phục vụ hai tư trên bảy đâu mà đòi ăn sáng giờ này!''

Kaiho nghiến răng kèn kẹt, mặt đỏ bừng vì tức tối: ``Ryuu! Em có tin chị ném em ra ngoài biển làm mồi cho cá mập không hả!?''

Rốt cuộc, người phụ nữ sở hữu sức mạnh ngàn cân có thể nhấc bổng cả xe tải đành phải lủi thủi đi ra tiệm tạp hóa nhỏ cạnh sảnh lễ tân, ngậm ngùi móc ví mua tạm một hộp mì tôm cốc cùng một cây xúc xích ăn liền.

Ngồi ở góc bàn ăn, Kaiho vừa cắm mặt xì xụp húp nước mì tôm nóng bỏng lưỡi vừa phồng má càu nhàu trong uất ức.

Suzuno ngồi bên cạnh khẽ mỉm cười dịu dàng, ân cần đẩy đĩa dưa hấu và hai chiếc bánh sừng bò mà mình đã cẩn thận gói ghém cất riêng từ trước sang trước mặt chị gái: ``Onee-sama ăn thêm bánh ngọt với dưa hấu đi ạ. Em biết chị ngủ dậy muộn sẽ đói nên đã xin phép các cô nhà bếp giữ lại một phần cho chị rồi đây.''

Đôi mắt Kaiho lập tức sáng bừng lên như bắt được vàng ròng. Cô nàng nhào tới ôm chầm lấy Suzuno, cảm động suýt rơi nước mắt: ``Suzu\dots!! Đúng là chỉ có Suzu là thương chị nhất trên đời thôi!! Không như cái đồ em gái tomboy mất nết nào đó chỉ biết đứng cười trên nỗi đau của người khác!''

Ryuuhou ngồi đối diện bĩu môi lè lưỡi: ``Xì! Suzu-nee lúc nào cũng nuông chiều chị quá đà!''

Nhìn cảnh ba chị em chí chóe rồi lại yêu thương đùm bọc nhau quanh chiếc bàn ăn nhỏ, tôi không khỏi bật cười lắc đầu. Một buổi sáng ồn ào và tràn ngập sinh khí hệt như ở nhà.

\ct
Sau khi giải quyết xong bữa sáng muộn của Kaiho, Suzuno rủ Chloe cùng dạo bước ra khu phố chợ nhỏ ven biển gần khách sạn để mua chút đặc sản địa phương và đồ ăn vặt.

Suzuno mở cuốn sổ tay nhỏ bìa da quen thuộc, vừa đi vừa tỉ mỉ ghi chép từng tên gọi đặc sản vùng Seika: từ bánh nhãn giòn tan, mực khô một nắng cho đến các loại kem dừa ép thủ công. Trong khi đó, cô nàng sát thủ cá voi Chloe thì mắt sáng rực như đèn pha ô tô khi đứng trước quầy kem bảy màu. Suzuno dịu dàng mua hai que kem ốc quế mát lạnh, đưa một que cho Chloe còn một que mình cầm trên tay định thưởng thức. Thế nhưng, trong lúc Suzuno còn đang chăm chú ghi chú công thức tỷ lệ đường trong món bánh mè vào sổ tay, Chloe đã ngoạm sạch que kem của mình trong vỏn vẹn ba nốt nhạc. Và rồi, bằng một phản xạ vô thức của loài sinh vật biển sâu phàm ăn, cô nàng tiện tay với luôn que kem trên tay Suzuno ăn sạch bách lúc nào không hay!

Đến khi Suzuno gấp sổ lại, nhìn bàn tay trống trơn của mình rồi ngơ ngác chớp mắt, Chloe mới giật mình ngậm chiếc ốc quế cuối cùng trong miệng, hai má phúng phính dính đầy vệt kem sô-cô-la, ngơ ngác nghiêng đầu hỏi một câu tỉnh bơ: ``Ủa\dots? Phần kem của Suzuno-san đâu mất tiêu rồi, poe?''

Suzuno chỉ biết mở to đôi mắt hai màu nhìn cô bạn nhỏ, rồi khẽ bật cười khúc khích, lấy chiếc khăn tay thêu hoa trong túi áo ra nhẹ nhàng lau khóe miệng cho con cá voi ngốc nghếch: ``Chloe-chan ăn hết luôn phần của chị rồi đấy\dots\ Thôi được rồi, để chị mua cho em thêm một que nữa nhé.''

Chloe đỏ bừng mặt vì xấu hổ, hai vành tai nóng ran, lí nhí cúi đầu bước theo sau cô nàng gia sư dịu dàng.

\ct

Bốn giờ chiều.

Kỳ nghỉ biển thực sự bùng nổ khi cái nắng gay gắt của buổi trưa đã dịu hẳn, nhường chỗ cho làn gió nồm mát rượi lồng lộng thổi qua vịnh Seika.

Toàn bộ lực lượng gồm gia đình Hikawa, hội chị em Gen 4 và cô nàng sát thủ Chloe lại một lần nữa tề tựu đông đủ trên bãi cát vàng. Sau bài khởi động nhanh gọn và nghiêm túc - lần này không có bất kỳ màn so sánh vòng một nào xảy ra nữa - cả đám lập tức ùa xuống làn nước biển mát lạnh.

Thế nhưng, khác với vẻ êm ả của buổi sáng, buổi chiều hôm ấy sóng biển Seika cuộn trào mạnh mẽ khác thường. Những con sóng bạc đầu cao ngút ngàn dồn dập đổ vào bờ, tạo nên một đấu trường nước đầy phấn khích.

Vừa chạm chân xuống nước, Luna đã ôm chiếc phao vịt, hai tay vỗ bôm bốp xuống mặt biển bắn nước tung tóe về phía Watame: ``Nhận lấy nè-nanora!!''

Watame bị bọt nước bắn vào mặt, lập tức vuốt tóc cười tít mắt: ``Ái chà chà! Dám khiêu chiến với tôi sao công chúa kẹo ngọt!? Đỡ lấy đòn phản công sừng cừu này!!''

Nàng Cừu dùng cả hai tay múc nước tạt mạnh trở lại. Hai cô nàng bắt đầu màn đấu té nước rộn rã tiếng cười, mở màn cho một khung cảnh nghỉ dưỡng rực rỡ hệt như những thước phim anime mùa hè kinh điển.

Kaiho đứng cạnh tôi, thấy không khí náo nhiệt liền nổi máu nghịch ngợm. Cô em song sinh bỗng nở một nụ cười tinh quái, hai bàn tay vục sâu xuống nước rồi hất một đợt sóng mini thẳng vào mặt tôi: ``Nhận lấy này Aniki!''

Tôi bị nước biển táp ướt sũng cả mặt mũi, lập tức ho khùng khục: ``Này nhé Kaiho! Em bao nhiêu tuổi rồi mà còn chơi trò trẻ con th!?''

``Thì đã sao nào!'' Kaiho cười khì, hai tay lại tiếp tục té nước xối xả: ``Đây là lần đầu em được đi biển thỏa thích mà! Lại còn được đi cùng anh nữa chứ! Không quậy phá cho đã thì phí hoài chuyến đi!''

Tôi thở dài lắc đầu, nhưng khóe môi đã cong lên từ lúc nào. Tôi cúi người vốc nước tạt trả đũa cô nàng: ``Được lắm! Xem ai té nước khỏe hơn nào!''

Hai anh em chúng tôi lập tức lao vào cuộc chiến té nước bì bõm, tiếng cười nói vang vọng khắp một góc bờ cát.

Cách đó một đoạn, bộ ba Kson, Kanata và Towa cũng bắt đầu nhập cuộc. Nhưng với bản tính hiếu thắng và thể lực phi phàm của các thành viên này, trò té nước thông thường nhanh chóng biến chất thành một cuộc đọ sức kinh thiên động địa.

Nàng Tổng trưởng đứng dạng chân vững như bàn thạch giữa làn sóng, hai cánh tay cơ bắp quạt nước rầm rầm như mái chèo tàu thủy: ``Nhận lấy này mấy đứa!! Đại tỷ cho các bồ tắm nước biển mát lạnh nhé!!''

Nàng Thiên sứ và nàng Ác ma bị từng đợt sóng nước giáng thẳng vào người đau điếng, lảo đảo lùi lại. Towa gạt nước khỏi trán, càu nhàu: ``Aiya! Đau đấy Kson! Bà té nước hay là đang ném bom nước thế!?''

Kanata siết chặt nắm đấm, đôi mắt rực lửa: ``Không thể để thua được! Towa, hợp lực lại!!''

``Được!! Nhận lấy này Kson!!''

Hai cô nàng Thiên sứ và Ác ma đồng loạt vung tay quạt nước hết tốc lực. Sự cộng hưởng giữa sức nắm năm mươi ký kinh hoàng của Kanata và sự linh hoạt của Towa đã tạo nên một luồng áp lực nước khủng khiếp. Kson cũng chẳng chịu kém cạnh, dốc toàn lực hất nước đáp trả.

Và rồi, một cảnh tượng phi lý đến mức bẻ gãy mọi định luật vật lý bắt đầu xuất hiện trước mắt toàn thể bãi biển.

Từ những cú quạt tay điên cuồng của ba người họ, làn nước biển xung quanh bỗng bắt đầu cuộn tròn lại, xoay tít mù rồi từ từ hình thành nên một vòng xoáy nước khổng lồ cao hơn đầu người ngay giữa vùng nước nông.

Tôi, Kaiho và các du khách xung quanh đồng loạt trợn tròn mắt, mồ hôi hột tuôn ra như tắm:

``Cái quái gì thế kia!? Xoáy nước thật kìa!!''

``Chạy mau!! Coi chừng nó hút chúng ta vào trong đấy!!''

Tiếng la hét thất thanh của đám đông bắt đầu vang lên hỗn loạn. Tôi ôm trán rên rỉ trong tuyệt vọng: ``Lạy Chúa tôi\dots\ Họ đang biến bãi biển thành bãi thử vũ khí ma thuật đấy à?''

Kaiho khoanh tay, trầm trồ: ``Mấy cô nàng này nhiệt tình ghê thật đấy\dots''

``Nhiệt tình cái con khỉ mốc ấy!'' tôi gắt lên: ``Đấy là điều anh đang lo sốt vó đây này!! Lộ hết danh tính bây giờ!!''

Watame và Luna đứng dạt sang một bên, vừa gặm kem vừa gật gù bình phẩm. Watame cười toe toét: ``Cảm giác hệt như mấy lần họ làm loạn ở văn phòng \hll\ nhỉ!''

Luna phụ họa: ``Đúng là zậy đó-nanora! Náo nhiệt xỉu lun á-nanora!''

Thế nhưng, thảm họa vẫn chưa chịu dừng lại ở đó.

Bị kích động bởi sự cổ vũ của đám đông và khí thế ngùn ngụt của Kson, cô nàng Thiên sứ Kanata bỗng bộc phát toàn bộ bản năng của một ``chú khỉ đột'' sở hữu thể lực vô địch. Kanata gầm lên một tiếng vang dội cả mặt biển:

``\textbf{Nhận lấy chiêu hiểm của ta đây--!!}''

Hai bàn tay hộ pháp của Kanata đập thẳng xuống mặt biển với một lực đạo ngàn cân.

*ẦM ĐÙNGGG!!*

Mặt nước biển nổ tung! Một bức tường sóng thần nhân tạo cao gần năm mét bất ngờ trồi lên khỏi mặt nước, cuộn trào sủi bọt trắng xóa như một con quái vật đại dương rồi hung hãn quét thẳng về phía bờ.

``\textbf{Lặn xuống!}'' tôi chỉ kịp gào lên một tiếng thất thanh trước khi con sóng khổng lồ ập xuống đầu.

*RẦM!!*

Cả tôi, Kaiho, Kson, Towa và toàn bộ những người đứng gần đó đều bị ngọn sóng thần nhân tạo cuốn phăng đi, lăn lông lốc cả chục vòng trên mặt cát rồi dạt thẳng lên bờ như những chú cá mắc cạn.

Tôi sặc sụa ho ra ngụm nước biển mặn chát, lồm cồm bò dậy thì thấy cảnh tượng xung quanh hoang tàn chẳng khác nào vừa trải qua một trận bão quét. Kson nằm ngửa thở dốc, Towa thì dính đầy rong biển trên đầu, còn Kaiho thì đang lắc lắc mái tóc ướt sũng, miệng cười khúc khích như vừa trải qua một trò chơi cảm giác mạnh.

Ở ngoài xa, Kanata đứng trơ trọi giữa làn nước phẳng lặng, chớp chớp mắt nhìn quanh với vẻ mặt ngơ ngác tột độ: ``Ơ\dots? Mọi người đâu cả rồi!? Sao tự nhiên bãi biển vắng tanh thế này!?''

Tôi nằm bẹp trên bãi cát, chỉ biết thều thào thở dài trong bất lực: ``Em dọa người ta chạy mất dép hết cả rồi chứ còn đâu nữa\dots''

Chưa bị đội cứu hộ bờ biển bế lên đồn cảnh sát vì tội gây rối trật tự công cộng đã là một kỳ tích vĩ đại lắm rồi. Đúng là một kỳ nghỉ hè long trời lở đất không một phút giây nào được bình yên!

\ct

Sau khi dẹp yên bãi chiến trường do cú đập nước long trời lở đất của Kanata gây ra, cả đám quyết định chuyển sang một hoạt động thể thao mang tính hòa bình hơn: thi bơi ba mươi mét tự do ngoài biển.

Tất nhiên, bơi ở một vùng biển có sóng bạc đầu dồn dập như Seika là một thử thách không hề đơn giản. Từng đợt sóng vỗ liên hồi có thể hất văng người bơi trở lại bờ, hoặc nguy hiểm hơn là cuốn phăng ra ngoài vùng nước sâu nguy hiểm nếu không có thể lực vững vàng.

Và dĩ nhiên, người không biết bơi như tôi thì xin phép đứng ngoài cuộc. Các bạn có thể sẽ chê bai rằng: ``Trời ạ, đàn ông con trai hai mươi sáu tuổi đầu mà không biết bơi!'' Nhưng xin thưa, người bình thường tối ngày cắm mặt vào giấy tờ sổ sách và xử lý khủng hoảng ở công ty như tôi thì lấy đâu ra thời gian mà rèn luyện bơi lội. May thay, tôi không phải là người duy nhất bơi như khúc gỗ. Luna cũng từ chối xuống bơi cự ly, cô nàng cứ ôm khư khư chiếc phao vịt màu hồng dạt dào ở mép nước nông.

Vì thế, tôi và cô nàng Công chúa tự động trở thành ban trọng tài quyền lực kiêm đội cổ vũ. Cuộc thi bơi ba mươi mét từ mép nước ra chiếc cờ phao màu cam - vốn dĩ là cái cờ để đánh mốc giới hạn có thể bơi an toàn - rồi quay về gồm có: Kaiho, Kson, Towa, Kanata, Watame và Chloe.

Tôi cầm chiếc còi nhựa giơ lên cao: ``Đường bơi ba mươi mét! Bơi từ vạch xuất phát ra chiếc phao cam đằng kia rồi quay ngược về mép cát! Ai chạm chân vào bờ trước là thắng!''

Luna ngậm kẹo mút, cất giọng đếm ngược: ``Ba, hai, một\dots\ GO-nanora!!''

*TÚUUT!!*

Tiếng còi vừa dứt, sáu bóng người lập tức lao vút xuống mặt nước.

Người dẫn đầu ngay từ những sải tay đầu tiên không ai khác chính là Chloe. Con cá voi sát thủ trườn đi như một mũi tên xé toạc mặt nước, uyển chuyển và nhanh đến mức mắt thường khó lòng bắt kịp. Ngay sát phía sau là Kaiho; với sải tay dài và cơ bắp săn chắc, cô em song sinh quạt nước mạnh mẽ như một chiếc cano cao tốc, rẽ sóng lao đi phăng phăng.

Towa bơi ở vị trí thứ ba. Dáng bơi của cô nàng Ác ma trông khá kỳ dị: hai tay quạt nước lạch bạch, hai chân đạp liên hồi hệt như một con gọng vó khổng lồ đang nổi lềnh phềnh trên mặt sóng, nhưng bằng một cách thần kỳ nào đó tốc độ vẫn rất đáng nể.

Tôi gật gù: ``Xem ra ba người này thi đấu rất nghiêm túc và bình thường đấy chứ.''

Thế nhưng, vừa dứt lời thì ba thí sinh còn lại lập tức chứng minh cho tôi thấy khái niệm ``bình thường'' chưa bao giờ tồn tại trong từ điển của họ.

Kanata bơi ở vị trí thứ tư, nhưng cô nàng hoàn toàn\dots\ không thèm huơ tay quạt nước! Chiếc vầng hào quang hình ngôi sao trên đỉnh đầu cô nàng bỗng quay tít mù như một chiếc cánh quạt trực thăng mini, tạo thành lực đẩy phản lực đưa nàng Thiên sứ lướt đi vù vù trên mặt biển.

Luna tròn xoe mắt: ``Nư zậy chẳng phải là ăn dan sao nanora!?''

Tôi ôm trán: ``Gian lận rành rành ra đấy rồi còn gì!''

Nhưng đó vẫn chưa là gì so với Watame. Nàng Cừu lặn ngụp hẳn xuống dưới đáy nước, chỉ nhô lên mặt biển một chiếc ống thở bằng nhựa cùng chỏm tóc vàng hoe bồng bềnh. Và ở ngay phía sau thắt lưng cô nàng\dots\ là một chiếc động cơ đẩy nước chạy bằng pin mini gắn chân vịt đang quay sùng sục, đẩy cô phóng đi vun vút như một quả ngư lôi.

Chưa dừng lại ở đó, đỉnh cao của sự phi lý chính là nàng cựu Tổng trưởng Kson.

Không thèm bơi lội tốn sức, Kson chẳng biết bằng cách nào đã nhảy tót lên một chiếc ca-nô cứu hộ gắn động cơ của đội tuần tra bãi biển. Cô nàng đứng hiên ngang ở mũi tàu, hai tay chống nạnh, áo hoa bay phần phật trong gió, gầm lên như một hải tặc: ``\textbf{Tiến lên!! Tiến thẳng về phía trước cho ta!!}''

Luna rớt cả chiếc kẹo mút xuống nước: ``Cái béo gì thế-nanora!?''

Tôi há hốc mồm hét toáng lên: ``Oắt-đờ-phấc!? Em kiếm đâu ra cái ca-nô đó thế hả Kson!?''

Câu trả lời lập tức xuất hiện ngay phía sau chiếc ca-nô: một toán nhân viên cứu hộ bãi biển mặc áo phao đỏ đang điên cuồng chèo hai chiếc thuyền nan đuổi theo sau, vừa thở không ra hơi vừa gào thét khản cả cổ bằng loa cầm tay: ``\textbf{Cướp!! Cô gái tóc tím kia mau tắt máy trả ca-nô lại cho chúng tôi!!}''

Cả bãi biển Seika được một phen nhốn nháo tột độ trước màn rượt đuổi bằng ca-nô có một không hai. Tiếng động cơ gầm rú ầm ĩ khiến Kaiho, Chloe và Towa đang bơi dẫn đầu cũng phải kinh hãi quay đầu tháo chạy thục mạng vào bờ: ``Tránh ra mau!! Nó tông vào người bây giờ!!''

Sau gần mười phút hỗn loạn, đội cứu hộ rốt cuộc cũng vây bắt được chiếc ca-nô của Kson. Nàng Tổng trưởng bị mẹ tôi và tôi lôi lên bờ mắng cho một trận vuốt mặt không kịp. Thế nhưng với bản tính bạt mạng trời không sợ đất không sợ, cô nfang chỉ gãi đầu cười hề hề rồi huýt sáo bài hát ``Tôi xin lỗi, tôi xin lỗi được chưa'' khiến ai nấy chỉ biết dở khóc dở cười.

\ct

Năm giờ chiều.

Sau màn thi bơi náo loạn, mọi người chia nhau ra các góc bãi biển để tận hưởng khoảng thời gian thư giãn cuối ngày.

Mẹ và bà ngoại ngồi nghỉ mát dưới tán ô cọ lớn, vừa uống nước dừa tươi vừa ngắm nhìn cảnh biển. Cách đó không xa ở mép nước êm dịu, Suzuno và bé út Naomi đang cùng nhau kiên nhẫn xây một tòa lâu đài cát khổng lồ.

Naomi hai tay xúc từng xẻng cát ướt đổ vào xô nhựa, mắt tròn xoe ngước nhìn chị gái thắc mắc: ``Chị Suzu ơi, sao cát khô ở trên bờ cứ bị tơi tả trôi tuột đi, mà cát ướt ở mép nước lại đắp thành tháp cao được thế ạ?''

Suzuno mỉm cười dịu dàng, vừa khéo léo dùng bay gọt phẳng bờ tường cát vừa cất giọng giải thích như một cô giáo hóa học nhỏ tuổi: ``Đó là nhờ hiện tượng sức căng bề mặt của nước đấy em. Giữa các hạt cát nhỏ li ti có những khoảng trống li ti. Khi có một lượng nước vừa đủ len lỏi vào, các phân tử nước sẽ tạo thành những cây cầu lỏng siêu nhỏ liên kết các hạt cát lại với nhau bằng lực mao dẫn, giúp khối cát kết dính bền vững hơn nhiều.''

Naomi chăm chú lắng nghe, đôi mắt màu lục bảo long lanh tràn đầy sự ngưỡng mộ tuyệt đối: ``Oa\dots\ Chị Suzu giỏi quá đi mất! Cái gì chị cũng biết hết trơn!''

Thế nhưng, vừa đắp xong chiếc tháp canh hình xoắn ốc thì một con sóng lớn bất ngờ trườn vào bờ, cuốn phăng đi một nửa chân thành. Không hề nản lòng hay buồn bã, hai chị em nhìn nhau bật cười khúc khích, rồi lại cùng nhau xắn tay áo múc cát đắp lại từ đầu với niềm vui thích say mê.

Cách đó một đoạn, tôi đang cầm chiếc máy ảnh du lịch compact nhỏ gọn, đứng ở mép cát canh góc chụp lại khoảnh khắc bình yên của gia đình và bãi biển.

Nhìn thấy tôi đang mải mê ngắm ống kính, cô nàng sát thủ cá voi Sakamata Chloe bỗng nảy ra một mưu đồ tình báo táo bạo. Chloe liếc nhìn điện thoại, rón rén bước từng bước nhẹ như mèo mò lại phía sau lưng tôi, lẩm bẩm trong đầu: \textit{`Thời cơ ngàn năm có một đây rồi! Mục tiêu đang sơ hở không phòng bị! Phải chụp trộm một bức ảnh dìm hàng lúc hắn đang ngơ ngác để nộp báo cáo cho Tổng soái La+ mới được!'}

Chloe giơ điện thoại lên, nheo một mắt ngắm thẳng vào gáy tôi, từ từ bước lùi lại để lấy góc rộng. Thế nhưng, cô nàng lại không hề để ý rằng ngay phía sau gót chân mình, Ryuuhou đang ngồi chồm hổm hì hục đắp một đụn cát lớn!

\textit{Bộp!}

Gót chân Chloe va thẳng vào đùi Ryuuhou. Cô nàng sát thủ mất đà, hai tay quờ quạng trong không trung: ``Oái! Poe poe--''

*XOÀNG!!*

Chloe ngã dập mặt xuống bãi cát ướt, cả người lấm lem bùn cát từ đầu đến chân! Chiếc điện thoại trên tay cô nàng rơi bẹp xuống đất, vô tình kích hoạt nút chụp ảnh cảm ứng đúng lúc chiếc mũi dính đầy cát của Ryuuhou đang thò sát vào màn hình với vẻ mặt hầm hố:

*TÁCH!*

Ryuuhou phủi cát trên đầu gối, trừng mắt nhìn con cá voi đang nằm đo ván dưới chân: ``Bà đi đứng kiểu gì mà thụt lùi như cua bò thế hả đồ cá voi ngốc!? Suýt nữa đạp nát đụn cát của tôi rồi đấy!''

Chloe lồm cồm bò dậy, nhổ ra một ngụm cát biển, run rẩy nhìn vào màn hình điện thoại: bức ảnh chụp trộm mục tiêu đâu chẳng thấy, chỉ thấy một bức ảnh macro chụp cận cảnh lỗ mũi và sống mũi dính đầy cát của Ryuuhou to tổ bố chiếm trọn khung hình! Cô nàng sát thủ chỉ biết ôm đầu khóc không ra tiếng vì chiến dịch tình báo bị phá sản hoàn toàn.

Tôi và Kaiho nằm dài trên hai chiếc ghế xếp gần đó chứng kiến toàn bộ màn kịch ngớ ngẩn chỉ biết ôm bụng cười nghiêng ngả.

Trong khi đó, ở ngoài khơi xa, Watame đang phấn khích cầm lái chiếc mô-tô nước cho thuê, kéo theo chiếc dù lượn khổng lồ chở Towa lơ lửng trên tầng không.

Towa ngắm nhìn toàn cảnh vịnh biển Seika từ trên cao, gió biển luồn qua tóc mát rượi. Thế nhưng khi tầm mắt quét qua bãi cát trắng trải dài gần trạm phát thanh, cô nàng bỗng nhìn thấy một cảnh tượng bất ngờ.

Một bé trai chừng năm tuổi mặc quần bơi xanh đang đứng ôm mặt khóc thút thít vì lạc bố mẹ. Nhưng trước khi Towa kịp ra hiệu cho Watame hạ dù lượn xuống, một bóng hình nhỏ nhắn đã tiến lại gần đứa trẻ trước.

Đó là bé út Naomi.

Đang ngồi nhặt vỏ ốc gần đó nghe tiếng khóc, Naomi đã lẳng lặng buông chiếc xô nhựa xuống, bước từng bước êm ả đến trước mặt cậu bé. Naomi khẽ ngồi xổm xuống ngang tầm mắt đứa trẻ, đôi mắt màu lục bảo tĩnh lặng như mặt hồ soi bóng vào đôi mắt đẫm lệ của cậu bé. Naomi đưa bàn tay nhỏ nhắn ra, nhẹ nhàng xoa lấy bờ vai đang run rẩy của em, cất giọng nói êm đềm, trong trẻo và thanh tịnh như tiếng chuông gió ngân nga giữa buổi chiều:

``Em đừng khóc nữa nhé\dots\ Có chị ở đây rồi. Mình cùng nhau đi tìm bố mẹ thôi nào.''

Kỳ diệu thay, luồng khí tức thanh khiết toát ra từ linh hồn Naomi dường như đã xoa dịu toàn bộ nỗi sợ hãi trong lòng đứa trẻ. Cậu bé bỗng nín bặt tiếng khóc, chớp chớp mắt nhìn cô bé trước mặt như nhìn thấy một thiên sứ nhỏ, rồi ngoan ngoãn gật đầu nắm lấy tay Naomi.

Đúng lúc ấy, chiếc dù lượn của Towa vừa vặn đáp xuống bãi cát gần đó. Towa và Watame hùng hổ lao tới định giúp đỡ, nhưng khi tận mắt chứng kiến khoảnh khắc Naomi dịu dàng dỗ dành đứa trẻ nín khóc chỉ bằng một câu nói êm dịu, cả nàng Ác ma lẫn nàng Cừu đều khựng lại, đứng ngẩn người ra vì kinh ngạc và xúc động.

Sau khi nắm rõ tình hình, cả nhóm quyết định dắt cậu bé vào thẳng buồng phát thanh bãi biển để phát loa tìm phụ huynh.

\textit{Rẹt\dots\ rẹt\dots}

Tiếng loa phóng thanh của bãi biển Seika bỗng vang lên khắp vịnh: ``\textbf{Thông báo tìm trẻ lạc! Thông báo tìm trẻ lạc!}''

Tôi và Kaiho đang ngồi uống nước dừa liền ngẩng đầu lên nghe ngóng. Kaiho tò mò: ``Trẻ lạc à anh?''

Tôi gật đầu: ``Mùa du lịch đông đúc thế này thì trẻ con mải chơi bị lạc là chuyện thường ngày ở huyện ấy mà.''

Chiếc loa lại tiếp tục phát ra tiếng phát thanh viên bối rối: ``Cháu bé\dots\ ờm, cháu tên gì nhỉ? Mau nói tên để bác đọc loa nào\dots''

Đáp lại chỉ là tiếng thút thít lí nhí của đứa trẻ. Và rồi, một giọng nữ quen thuộc bất ngờ giật lấy micro quát toáng lên:

Tôi giật bắn mình suýt sặc ngụm nước dừa: ``Ủa!? Giọng Towa-chan đấy à!? Sao em ấy lại ở trong buồng phát thanh thế kia!?''

Bị giọng quát đầy uy lực của nàng Ác ma làm cho giật mình, cậu bé thút thít đáp qua loa: ``D-Dạ\dots\ Cháu tên là Kuro ạ\dots''

Towa hài lòng: ``Tốt! Bố mẹ cháu tên gì?''

``Mẹ cháu là Hasaka\dots\ Bố cháu là Rojou ạ\dots''

Ngay lập tức, chất giọng đanh thép của Towa vang vọng khắp bờ biển qua hệ thống loa công suất lớn: ``\textbf{Gia đình anh Rojou và chị Hasaka ở đâu, xin mời đến ngay trạm cứu hộ số ba nhận cháu Kuro đang bị lạc! Nhắc lại, cháu Kuro đang ở trạm cứu hộ số ba, gia đình mau đến nhận diện!}''

Nghe tiếng thông báo dõng dạc và đầy uy quyền của nàng Ác ma, tôi không khỏi ôm trán lắc đầu cười trừ: ``Cái cô nàng này\dots\ bình thường thì bẽn lẽn, lúc cần kíp thì ra dáng phát thanh viên hình sự ghê thật.''

Chỉ chưa đầy mười phút sau, một cặp vợ chồng trẻ hớt hải chạy đến trạm cứu hộ nhận lại đứa con trai bé bỏng. Nhìn cảnh hai mẹ con ôm chầm lấy nhau khóc trong mừng rỡ, Towa đứng bên cạnh tuy khoanh tay quay mặt đi làm vẻ lạnh lùng nhưng khóe môi lại khẽ cong lên một nụ cười ấm áp. Đúng là một Ác ma nhưng trái tim lúc nào cũng nhân hậu như một thiên thần.

Thế nhưng, vừa giải quyết xong vụ trẻ lạc thì một biến cố mới lại ập đến.

Tôi ngó quanh bãi biển để điểm danh lại các thành viên thì bỗng khựng lại: ``Ủa? Kaiho này\dots\ Nãy giờ em có thấy Luna-chan đâu không?''

Kaiho chớp mắt, lập tức trèo tót lên vai tôi đứng kiễng chân ngó ra xa: ``Để em nhìn thử\dots\ Oái! Aniki, nhìn đằng kia kìa!!''

Tôi nhìn theo hướng tay cô nàng chỉ, tim muốn rớt ra ngoài lồng ngực.

Cách bờ cát gần trăm mét, vượt qua cả phao cờ cảnh báo nguy hiểm, chiếc phao tròn màu hồng phấn của Luna đang trôi dạt vô định giữa những con sóng lớn ngoài khơi xa. Cô nàng công chúa kẹo ngọt không hề hay biết hiểm nguy, vẫn đang nằm ngửa ngậm kẹo mút thảnh thơi tắm nắng!

``Chết tiệt! Con bé không biết bơi!'' Tôi hoảng hốt gào lên.

Kaiho không chần chừ nửa giây, đôi mắt màu lam lóe lên tia quyết liệt. Cô nàng lập tức lao mình xuống mép nước định bơi thẳng ra khơi cứu người: ``Để em ra lôi con bé về!!''

Hai cô nàng đang đắp lâu đài cát nghe thấy vậy liền bật dậy như lò xo. Kson nhìn ra ngoài khơi, nghiến răng: ``Để em! Em bơi ra túm con bé về ngay!''

Thế nhưng ngay khi chân Kaiho vừa chạm làn nước sâu ngang gối, một bàn tay mảnh mai đã nhanh như cắt túm chặt lấy cổ tay cô kéo giật ngược trở lại!

Suzuno thở hổn hển, gương mặt tái nhợt nhưng ánh mắt vô cùng cương quyết: ``Onee-sama, dừng lại ngay! Khu vực ngoài cờ phao đó có dòng chảy xa bờ ngầm cực kỳ nguy hiểm! Dù thể lực của chị có khỏe đến đâu thì bơi ngược dòng cũng sẽ bị kiệt sức ngay lập tức đấy ạ!!''

Bị Suzuno dùng kiến thức khoa học giữ chặt lại, Kaiho buộc phải khựng bước. Cô nàng nắm chặt lấy cột phao chắn sóng bên cạnh đến mức thanh thép rên lên ken két, ánh mắt bồn chồn căng thẳng tột độ nhìn về phía tôi: ``Aniki\dots! Làm sao bây giờ!?''

Đúng lúc đó, Chloe đang đứng cạnh vội vã xắn tay áo phao lên: ``Để tôi bơi ra! Tôi là sinh vật biển sâu, tôi có thể lặn dưới dòng chảy ngược--''

Chưa kịp dứt lời, Ryuuhou từ phía sau đã thẳng tay đè dúi đầu cô nàng sát thủ xuống cát: ``Chị là cá voi sát thủ chứ có phải tàu ngầm hạt nhân đâu mà đòi vượt dòng chảy xiết! Điên à mà lao ra đó để bị cuốn đi mất xác hả đồ ngốc!!''

Chloe bị đè đầu ú ớ trong cát: ``Poe poe\dots\ Người ta có lòng tốt muốn cứu người mà\dots''

Tình thế ngàn cân treo sợi tóc khiến cả bãi biển như nín thở. Đúng lúc ấy, Kson—người đang khởi động bên cạnh—nhìn ra vị trí chiếc phao của Luna, tính toán hướng gió và dòng chảy ven bờ rồi dõng dạc tuyên bố: ``Để đó cho tôi! Tôi bơi vòng theo rìa dòng xoáy từ mạn trái sang, bảo đảm an toàn trăm phần trăm!''

Với thể lực phi thường và kinh nghiệm lướt sóng dạn dày vùng biển California, Kson lao mình xuống biển như một mũi tên rẽ nước. Đôi cánh tay vạm vỡ quạt nước theo đường vòng cung chiến thuật xé toạc những đợt sóng bạc đầu, chỉ mất chưa đầy bốn mươi giây cô nàng đã áp sát được chiếc phao của Luna.

Kson một tay tóm chặt lấy sợi dây viền phao, một tay quạt nước kéo chiếc phao nương theo con sóng thuận chiều quay ngược về phía bờ: ``Bám chắc vào con nhóc này! Trôi ra tít ngoài này mà còn nằm ngửa ngậm kẹo mút được à!?''

Luna chớp chớp mắt, phụng phịu: ``Tui có bít zì đâu-nanora\dots\ Đang ngắm mây tự nhiên thấy biển sâu hoắm à-nanora\dots''

Khi Kson kéo được Luna an toàn vào bờ cát, cả đoàn chúng tôi mới thở phào nhẹ nhõm một hơi dài như vừa trút được tảng đá ngàn cân. Mẹ Kaien lập tức bước tới cốc nhẹ lên đầu cô công chúa kẹo ngọt một cái rõ đau để răn đe, khiến nàng Công chúa xứ Kẹo phải ôm đầu mếu máo xin lỗi trong tiếng cười rộn rã của mọi người.

Hoàng hôn dần buông xuống trên vịnh biển Seika, nhuộm cả vùng trời thành một màu đỏ tía rực rỡ và lãng mạn. Một ngày thứ hai ngập tràn những pha rượt đuổi, cứu hộ và quậy phá tưng bừng đã khép lại trong sự an toàn trọn vẹn.

\ct

\begin{center}
    \uline{Chín giờ tối.\\Phòng 420B.}
\end{center}

Sau bữa tối ăn vội ở nhà ăn khách sạn (thực đơn tối nay có phần thanh đạm hơn buổi trưa vì ngư dân ven biển bắt đầu neo tàu tránh đợt áp thấp ngoài khơi) mọi người lại lục tục kéo về phòng nghỉ ngơi.

Naomi sau cả ngày vui đùa mệt nhoài đã được mẹ và bà ngoại dỗ dành cho ngủ say tít trên chiếc giường lớn phía trong. Trong khi đó, Ryuuhou, Suzuno vàLuna đã rủ nhau xin phép ra bãi cát trước sảnh đi dạo hóng gió biển đêm.

Thế là gian phòng khách rộng rãi của phòng 420B bỗng trở thành một đấu trường mới toanh khi Kson, Kanata, Watame và Towa ùa sang, tụ tập cùng tôi và Kaiho để làm một chầu bài UNO xuyên màn đêm.

Chloe cũng kéo một chiếc đôn nhỏ ngồi nép ở góc bàn, hai tay ôm gói snack mực nướng vừa nhai nhồm nhoàm vừa chớp chớp mắt theo dõi, kiên trì không rời mục tiêu nửa bước.

Chúng tôi đã chiến đấu bền bỉ suốt hơn một tiếng rưỡi đồng hồ qua tám ván đấu nảy lửa. Tôi và Kaiho chia nhau mỗi người hai ván thắng, Watame tinh quái thắng ba ván, Kson ẵm một ván. Đây đã là ván thứ chín, ván chốt hạ cuối cùng của đêm nay. Cặp đôi Kanata và Towa đang ở thế đường cùng, mắt bốc hỏa hừng hực quyết tâm giành lấy một chiến thắng danh dự trước giờ đi ngủ.

Vòng tròn lượt chơi được ấn định: Kson, tôi, Kanata, Kaiho, Watame, và cuối cùng là Towa, trước khi quay về cô nàng Tổng trưởng.

Tôi đặt một lá bài đổi màu xuống giữa bàn, dõng dạc hô to: ``Đổi sang màu đỏ! Bốc thêm một lá!''

Nàng Thiên sứ nhìn bộ bài trên tay, mặt méo xệch: ``Ể!? Em làm gì có lá màu đỏ nào đâu trời!!''

Nàng Ác ma khoanh tay đắc chí: ``Không có thì bốc thêm bài đi chứ còn chờ gì nữa!''

Kanata càu nhàu bốc thêm một lá từ xấp bài chung. Thế nhưng ngay khi nhìn thấy quân bài vừa rút được, sắc mặt cô nàng bỗng biến đổi một trăm tám mươi độ, hai mắt sáng quắc lên như đèn pha: ``Rồi! Mấy người tới số với tôi rồi!''

Kaiho ngồi ngay cạnh liền nhếch mép cười nhàn nhạt: ``Chị thì\dots\ xin chặn họng em nhé!''

Nói rồi, cô nàng dằn mạnh một lá bài Cấm Lượt (Skip) màu đỏ xuống bàn. Watame ngồi kế bên xuýt xoa: ``Hể\~{} Ghê thật đấy nha\~{}!''

Đến lượt Towa, cô nàng Ác ma đập ngay một lá số 7 màu đỏ xuống mặt bàn: ``Em đánh lá số bảy! Theo luật gia đình thì em được quyền đổi toàn bộ bài trên tay với người khác! Và em chọn đổi với\dots\ Kson!''

Kson trố mắt: ``Ối chà! Khôn thế! Bài tôi xấu mù thế này mà cũng đòi đổi à!?''

Hai người tráo bài cho nhau. Towa hí hửng lật xấp bài mới lên, miệng cười toe toét: ``Haha! Có lá cộng bốn này!!''

Và thế là, một trong những ngón đòn tàn nhẫn và hủy diệt tình bạn nhất của bài UNO chính thức khai hỏa: luật cộng dồn lá phạt!

Towa vung tay ném xuống: ``\textbf{Cộng bốn!}''

Tôi mỉm cười nham hiểm, thả ngay một lá bài xuống chồng: ``Hay đấy. \textbf{Cộng hai!}''

Kanata gầm gừ tiếp chiêu: ``Em chờ giây phút này lâu lắm rồi! \textbf{Cộng hai!}''

Kaiho không hề nao núng, bồi thêm: ``\textbf{Cộng bốn!}''

Watame cười khúc khích: ``Cho em góp vui với nha\~{}! \textbf{Cộng hai!}''

Kson đập bàn hô lớn: ``Một vòng tròn hoàn hảo rồi nhé! \textbf{Cộng hai!}''

Tôi lại tiếp tục ném: ``\textbf{Cộng hai!}''

Kanata nghiến răng: ``\textbf{Cộng bốn!!}''

Cả căn phòng bỗng nín thở khi tổng số lá bài phạt tích tụ đã lên tới con số khủng khiếp. Mặt Kaiho thoáng tái mét, toan đưa tay ra bốc bài thì đột nhiên cô nàng bật cười khanh khách: ``Đùa chút thôi! \textbf{Cộng hai!}''

Watame giật thót nhưng vẫn kịp lật bài: ``\textbf{Cộng hai!}''

Towa không một giây do dự, dằn tiếp lá bài tẩy cuối cùng: ``\textbf{Cộng bốn!}''

Kson trợn tròn mắt: ``Vãi chưởng thật chứ! Tôi còn một lá cuối đây: \textbf{Cộng hai!}''

Cả bàn cờ lúc này chỉ còn lại duy nhất một người có khả năng gánh chịu quả bom nổ chậm ấy. Tôi nhìn xấp bài trên tay, ung dung rút ra quân bài đen tuyền cuối cùng nhẹ nhàng đặt xuống, dõng dạc tuyên bố: ``\textbf{Cộng bốn! UNO!!}''

Cả căn phòng nổ tung trong sự kinh hoàng. Tổng cộng chuỗi cộng dồn liên hoàn từ nãy đến giờ đã đạt tới kỷ lục vô tiền khoáng hậu: \emph{ba mươi sáu} lá bài phạt!

Tôi chỉ còn đúng một quân bài duy nhất trên tay. Và người ngồi ngay cạnh tôi gánh trọn quả bom ba mươi sáu lá ấy không ai khác chính là\dots\ Amane Kanata!

Mặt nàng Thiên sứ trắng bệch không còn một giọt máu. Cô nàng run rẩy nhìn xấp bài trên tay mình, rồi ngước nhìn xấp bài chung dày cộp giữa bàn. Cơm toi thật rồi! Em ấy không còn bất kỳ một lá cộng dồn nào để đỡ đòn nữa!

Kanata nghiến răng ken két, hai tay ôm đầu gào thét trong cơn uất ức tột cùng: ``\textbf{Tại sao mọi chuyện lại trở nên như thế này!!}''

Kson cười trừ: ``Thì chính bồ là người hào hứng đồng ý chơi luật cộng dồn cơ mà, sao bây gờ còn ý kiến?''

Towa châm chọc: ``Có chơi thì có chịu đi bạn hiền ơi!''

Watame xoa tay cười gian: ``Bốc bài đi bạn ơi\~{} Tròn ba mươi sáu lá đấy, tha hồ mà xếp quạt nhé\~{}!''

Kanata tức giận đến mức lăn lộn ra sàn nhà ăn vạ, hai chân đạp lung tung, hét toáng lên bằng chất giọng thánh thót làm rung chuyển cả căn phòng: ``\textbf{KHÔNG CHỊU ĐÂU! KHÔNG CHỊU! KHÔNG CHỊU ĐÂU MÀ HU HU HU--!!}''

Tiếng hét long trời lở đất của nàng Thiên sứ vừa dứt thì\dots

*ĐOÀNGGG!!*

Một tia sét xé toạc màn đêm đen đặc ngoài cửa sổ, kèm theo tiếng sấm sét nổ vang rền như bom dội ngay trên nóc khách sạn! Gió rít gầm rú điên cuồng làm toàn bộ cửa kính ban công rung lên bần bật.

*LẠCH CẠCH* *LẠCH CẠCH*

Mẹ và bà ngoại ở giường bên giật bắn mình tỉnh giấc: ``Cái gì đấy!? Tiếng gì nổ om sòm thế!?''

Kanata luống cuống xua tay chối tội: ``K-Không phải cháu hét làm vỡ kính đâu ạ! Ở bên ngoài! Bão ở bên ngoài đánh sấm đấy ạ!!''

Mẹ tôi bỗng giật thót, nhìn ra ban công rồi thảng thốt la lên: ``Chết dở rồi! Đống quần áo bơi hồi chiều mẹ đang phơi ngoài dây thép ban công! Kuon! Kaiho! Hai đứa mau chạy ra thu dọn quần áo vào đây cho mẹ nhanh lên!!''

Tôi lập tức lao ra cửa ban công, nắm lấy tay nắm cửa giật mạnh. Thế nhưng ngay khi cánh cửa vừa hé mở, một luồng cuồng phong kèm theo mưa quất rát mặt như đạn bắn đã thổi thốc vào trong phòng.

Sức gió khủng khiếp đẩy ngược cả người tôi văng dội trở lại, đập mạnh vào cánh cửa phòng ngủ kêu một tiếng *RẦM!*.

Tôi đau điếng ôm lưng lồm cồm bò dậy: ``Vãi chưởng\dots\ Gió gì mà như cuồng phong cấp mười thế này!?''

Thấy ông anh cả bị gió thổi bay ngược, Kaiho không chần chừ nửa giây. Bản năng bảo vệ gia đình của cô nàng lập tức trỗi dậy. Kaiho siết chặt nắm đấm, hạ thấp trọng tâm, dùng sức mạnh thể chất ngàn cân bước từng bước vững như bàn thạch lao ra mép cửa: ``Để đó cho em!! Aniki, bám sát sau lưng em!!''

Tôi lập tức chồm dậy, ghì chặt lấy thắt lưng của Kaiho làm điểm tựa. Hai anh em phối hợp nhịp nhàng như một cỗ xe tăng bọc thép: Kaiho dùng bờ vai và cơ bắp vững chãi chắn trọn luồng gió rít điên cuồng phía trước, hai tay thoăn thoắt giật từng chiếc áo quần trên dây thép ném ngược ra sau; còn tôi ở phía sau chịu trách nhiệm ôm trọn lấy đống quần áo ướt sũng ném thẳng vào trong phòng cho Kson và Watame đỡ lấy.

Chỉ trong chưa đầy một phút chiến đấu nghẹt thở giữa tâm gió, hai anh em chúng tôi đã thu dọn sạch bách dây phơi rồi cùng nhau dồn sức đóng sập cánh cửa ban công lại.

*RẦM!!*

Hai anh em thở hồng hộc ngồi bệt xuống sàn nhà, quần áo ướt nhẹp nước mưa. Kson vội đưa khăn bông cho hai chúng tôi, mắt nhìn ra ngoài trời mưa như thác đổ: ``Đệch mợ, gió bão ở vùng biển này kinh khủng thật đấy!''

Kanata lúc này đã mở vội chiếc máy tính xách tay của tôi lên để tra cứu radar thời tiết, giọng hốt hoảng: ``Ủa! Đúng luôn này mọi người ơi! Trung tâm khí tượng vừa phát thông báo khẩn: một dải áp thấp nhiệt đới vừa bất ngờ mạnh lên thành bão và quét thẳng vào bờ biển Seika đây này!!''

Towa biến sắc: ``Bão thật á!?''

Nghe đến hai chữ ``bão biển'', tim tôi bỗng khựng lại một nhịp. Một luồng điện lạnh toát chạy dọc sống lưng.

Tôi quay ngoắt sang nhìn quanh phòng: ``Khoan đã\dots\ Suzuno, Ryuuhou với Luna đâu rồi!?''

Cả căn phòng bỗng chốc im phăng phắc.

Mẹ Kaien tái mặt: ``Lúc nãy\dots\ lúc nãy ba đứa nó xin phép mẹ ra bãi biển đi dạo hóng mát mà!!''

Bão biển đổ bộ bất thình lình, gió giật cấp mười quất sập cả cây cối ngoài đường, mà ba đứa con gái - một đứa út ít lơ ngơ, một đứa ngây thơ không biết bơi và một cô bé mảnh mai - lại đang kẹt lại giữa bãi cát mênh mông không một bóng người che chắn!

``Chết tiệt!!''

Tôi không kịp nghĩ ngợi thêm một giây nào nữa, lập tức vùng dậy lao thẳng ra cửa phòng.

Kaiho cũng nhanh như chớp giật lấy chiếc áo khoác gió bên cạnh, đôi mắt màu lam lóe lên tia sát khí đằng đằng, lao vút theo sau tôi: ``Chờ em với Aniki! Em đi cùng anh!!''

Đêm giông bão kinh hoàng tại bờ biển Seika chính thức bắt đầu.

\ct

Hai anh em chúng tôi lao như bay xuống cầu thang bộ. Tiếng gió rít qua khe cửa sổ thông gió gầm rú như tiếng gầm của dã thú, từng đợt chớp lóe lên xanh lè soi rõ vẻ mặt căng thẳng tột độ của Kaiho.

Cô em song sinh của tôi bình thường điềm tĩnh và sắt đá là thế, nhưng lúc này hai bàn tay cô siết chặt đến mức các khớp ngón tay trắng bệch, cả người run lên bần bật vì lo lắng: ``Suzu\dots\ Thể lực Suzu yếu nhất nhà, lại hay bị nhiễm lạnh khi gặp mưa gió\dots\ Nếu con bé có mệnh hệ gì, em sẽ không tha thứ cho bản thân mình mất!!''

Tôi vừa sải bước dài vừa đặt tay lên vai cô nàng trấn an: ``Bình tĩnh lại Kaiho! Suzuno thông minh lắm, con bé sẽ biết tìm chỗ trú an toàn cùng Ryuuhou và Luna! Chúng ta ra đón ba đứa nó về ngay bây giờ!''

Xuống đến sảnh chính tầng một, toàn bộ các cửa kính lớn hướng ra ngoài biển đều đã được khóa chặt và chốt then cài bằng những thanh gỗ lớn. Hai bác bảo vệ khách sạn đang túc trực bên cánh cửa ngách, gương mặt đầy vẻ căng thẳng.

Nhìn thấy hai anh em tôi xông tới toan mở chốt cửa, bác bảo vệ lập tức dang tay cản lại: ``Hai cháu định làm gì đấy!? Bão quét thẳng vào bờ rồi, gió giật gãy cả cột đèn ngoài kia kìa, tuyệt đối không được ra ngoài!''

Kaiho bước lên một bước, đôi mắt màu lam sắc lẹm lóe lên luồng uy áp đáng sợ khiến người bảo vệ bất giác lạnh gáy: ``Em gái cháu và bạn của chúng cháu vẫn đang kẹt ngoài bãi cát! Cháu phải ra đón chúng nó về!!''

Tôi tiếp lời dứt khoát: ``Xin bác mở cửa cho chúng cháu! Chúng cháu sẽ đi men theo bờ tường tìm một chút rồi quay về ngay ạ!''

Nhận thấy sự kiên định không gì lay chuyển nổi trong ánh mắt của hai anh em, người bảo vệ già khẽ thở dài, đành nhượng bộ mở hé cánh cửa ngách rồi dúi vào tay tôi một chiếc ô sắt cỡ lớn của khách sạn: ``Cầm lấy cái này! Đi men sát chân tường nhà hàng, thấy người là phải kéo về ngay đấy nhé!''

``Vâng ạ! Cháu cảm ơn bác!''

\textit{ẦM!!}

Vừa bước chân ra khỏi mái hiên sảnh, một luồng cuồng phong lạnh buốt kèm theo màn mưa dày đặc như những mũi kim châm thẳng vào da thịt quất tới tấp vào người chúng tôi. Cả đất trời tối sầm lại, chỉ có ánh chớp thi thoảng rạch ngang bầu trời rọi rõ những cành cây phi lao bị quật gãy đổ rạp sang một bên.

Kaiho lập tức hạ thấp trọng tâm, đứng chắn phía trước đón đầu ngọn gió, hai chân ghim chặt xuống nền gạch ướt như mọc rễ: ``Aniki! Bám sát em!! Đi theo hướng nhà hàng ven biển hồi trưa!!''

Tôi bung chiếc ô sắt lớn lên để che bớt những mảnh tôn và cành cây vụn đang bay loạn xạ. Thế nhưng, tôi đã đánh giá quá thấp sức mạnh của bão biển.

Ngay khi chiếc ô vừa mở toang, một luồng gió lốc cực mạnh từ phía biển bất ngờ thốc ngược từ dưới lên. Toàn bộ vòm ô phồng to như một cánh buồm căng gió, nhấc bổng cả thân hình nặng gần tám mươi ký của tôi lên khỏi mặt đất.

``\textbf{Á Á Á!! Cái gì thế nầy!??}'' tôi hoảng loạn gào thét giữa không trung.

Kaiho quay phắt lại, mắt trợn tròn kinh hãi: ``\textbf{Aniki!!}''

Cô nàng nhanh như chớp vươn tay chộp lấy cổ chân tôi, nhưng thay vì kéo tôi xuống được, sức nâng khủng khiếp của cơn bão kết hợp với trọng lượng nhẹ hơn của Kaiho lại kéo luôn cả hai anh em chúng tôi bay là là trên tầng không, hướng thẳng về phía dải nhà hàng ven biển.

Tôi vừa dùng hai tay ghì chặt cán ô vừa la toáng lên: ``\textbf{Kaiho! Đừng buông tay! Nhìn xuống dưới và tìm ánh đèn pin đi!}''

Kaiho một tay bám chặt lấy mắt cá chân tôi, cả người lơ lửng trong gió bão, đôi mắt màu lam căng ra quét xuyên qua màn mưa mịt mù: ``Đằng kia!! Cạnh mái hiên nhà hàng có ánh sáng đang đung đưa kìa anh!!''

\ct

Tại mái hiên của một căn nhà hàng đóng kín cửa sát mép biển.

Suzuno, Ryuuhou và Luna đang nép sát vào góc tường đá thụt sâu vào trong. Quần áo cả ba ướt sũng, dán chặt vào người lạnh buốt. Nàng Kẹo run rẩy ôm chặt lấy cánh tay Suzuno, hai hàm răng đánh vào nhau lập cập: ``Hức\dots\ Lạnh quá nanora\dots\ Gió thổi bay cả dép rùi nanora\dots''

Suzuno cởi chiếc áo khoác mỏng của mình ra, cẩn thận quấn quanh bờ vai nhỏ bé của Luna rồi ôm chặt cô em vào lòng, khẽ xoa hai bàn tay lạnh ngắt của em: ``Cố lên một chút nữa nhé Luna-san. Chúng ta nấp ở đây rất an toàn, tường bê tông kiên cố sẽ chắn được cây đổ. Ở yên đây chờ gió dịu nhé.''

Thế nhưng, giọng của Suzuno cũng không giấu nổi chút run rẩy vì lạnh.

*RẮC!! RẦM!!*

Một thân cây dừa lớn cách đó vài mét bất ngờ bị gió bão vặn gãy đôi, đổ sập xuống bờ cát tạo nên một tiếng động kinh hoàng. Luna sợ hãi hét lên một tiếng thất thanh, rúc sâu vào ngực Suzuno.

Ryuuhou đứng chắn phía ngoài, trên tay cầm chiếc điện thoại bật sáng đèn flash pin, đôi mắt cá tính không giấu nổi vẻ sốt ruột. Cô em thứ năm vừa lia đèn pin vào màn mưa vừa càu nhàu gắt gỏng để át đi nỗi sợ: ``Mẹ kiếp! Cái thời tiết khốn nạn này ở đâu ra vậy không biết! Đang dạo biển ngon lành tự nhiên nổi giông! Không biết Onii với Kaiho-nee ở khách sạn có nhận ra bọn mình chưa nữa!''

Suzuno khẽ ngước lên, ánh mắt tràn ngập niềm tin tưởng tuyệt đối: ``Onii-sama và Onee-sama chắc chắn sẽ nhận ra và đến tìm chúng mình ngay thôi. Ryuuhou-chan, em cứ tiếp tục giơ đèn pin báo hiệu nhé.''

Đúng lúc đó, từ trên vòm trời tối đen mịt mùng, một âm thanh gào thét quen thuộc hòa cùng tiếng gió rít vọng xuống: ``\textbf{Ryuu! Suzu! Luna-tan! Mấy đứa trốn ở đâu rồi!!}''

Ryuuhou giật thót mình, vểnh tai lên nghe ngóng rồi nhảy cẫng lên mừng rỡ: ``Tiếng hét gãy giọng của Kaiho-nee kìa!! Luna-senpai! Suzu-nee! Anh ấy đến rồi!!''

Luna ngước lên chớp chớp mắt: ``Ở đâu zậy nanora!?''

Ryuuhou ngẩng đầu nhìn lên không trung, rồi cả hàm dưới của cô nàng suýt chút nữa rớt xuống đất: ``C-Cái quái gì thế kia!? Sao Onii với Kaiho-nee lại\dots\ bay lơ lửng trên trời như Mary Poppins thế kia trời!?''

Từ trên cao, tôi và Kaiho nhìn thấy vệt đèn pin đang đung đưa dữ dội dưới mái hiên. Nhận thấy mục tiêu, tôi lập tức nghiêng cán ô sang một bên để giảm sức gió: ``\textbf{Họ ở đó! Kaiho, chuẩn bị tiếp đất!}''

Cơn gió quái ác kéo chiếc ô lao xuống như một quả đại bác không phanh.

``\textbf{Tránh ra!!!}''

*ẦM ĐÙNGGG!!*

Cả tôi và chiếc ô sắt lao thẳng vào bờ tường cạnh mái hiên. Tôi đau điếng khuỵu xuống nền gạch, sống mũi va mạnh vào gờ tường rớm máu, chiếc ô gãy gập văng ra một bên. Trong khi đó, Kaiho nhờ phản xạ vận động đỉnh cao đã kịp buông chân tôi ra trước một giây, tiếp đất bằng một cú lộn nhào hoàn hảo rồi lao vút tới ôm chầm lấy Suzuno.

Kaiho thở hổn hển, hai tay ôm chặt lấy bả vai và mặt Suzuno, kiểm tra khắp người cô bé với giọng run rẩy nghẹn ngào: ``Suzu!! Em có bị thương ở đâu không!? Có bị đá dăm hay cành cây quẹt trúng không em!?''

Suzuno rưng rưng nước mắt, đưa hai bàn tay nhỏ bé áp lên gò má ướt đẫm nước mưa của chị gái: ``Em không sao ạ\dots\ Onee-sama đừng lo, em không bị thương tí nào cả\dots''

Trong khi đó, Ryuuhou và Luna chạy vội lại đỡ tôi dậy. Con bé nhìn thấy một vệt máu đỏ tươi đang chảy ra từ khóe mũi tôi liền tròn mắt kinh hãi: ``Trời đất ơi Onii!! Anh bay kiểu gì mà lấy mặt phanh vào tường thế hả!? Có sao không đấy!?''

Tôi lồm cồm bò dậy, quệt ngang mu bàn tay lau vết máu, thở dốc phì phò: ``Anh\dots\ anh ổn\dots\ Mấy đứa bình an vô sự là được rồi\dots''

Ryuuhou vội vàng móc trong túi ra một mẩu khăn giấy vo tròn nhét tạm vào lỗ mũi tôi để cầm máu, miệng vừa càu nhàu vừa lộ rõ vẻ xót xa: ``Đầu anh đúng là cứng như đá hoa cương ấy! Đến lúc này rồi mà còn mạnh mồm!''

Luna đứng bên cạnh hai mắt đỏ hoe, nắm lấy vạt áo tôi: ``Hức\dots\ Kuon-senpai chảy máu mũi rùi nanora\dots\ Cảm ơn anh đã đến cứu bọn em nanora\dots''

Gió bão bên ngoài càng lúc càng gầm rú dữ dội hơn, từng mảng tôn của các ki-ốt xung quanh bắt đầu bay loảng xoảng. Nơi trú ẩn này chắc chắn sẽ không trụ được lâu nếu bão đổ bộ trực diện.

Kaiho quay lại nhìn tôi, gương mặt lấy lại vẻ đanh thép quyết đoán: ``Aniki! Chỗ này nguy hiểm lắm, mái tôn sắp sập rồi! Mình phải đưa ba đứa nó về khách sạn ngay lập tức!''

Tôi nhìn con đường ngập ngụa nước phía trước, gió quật ngược chiều cực mạnh: ``Gió ngược thế này chạy bộ bình thường sẽ bị thổi bay ngay. Nhưng anh em mình cùng mở đường thì được! Kaiho, em đi trước dẫn đầu cản gió cho ba đứa nhỏ! Anh bọc hậu phía sau giữ chặt Luna và Suzu! Cả nhóm nắm chặt tay nhau thành một hàng dọc, tuyệt đối không được buông tay!!''

``Rõ!!''

Kaiho khoác chiếc áo khoác gió lên người Suzuno, rồi hiên ngang bước ra đón đầu ngọn cuồng phong. Với thể lực phi phàm của mình, Kaiho như một bức tường thành thép sừng sững xé toạc làn gió bão, che chắn trọn vẹn cho Suzuno và Luna ở giữa. Ryuuhou nắm chặt tay Luna, còn tôi dùng cả tấm lưng rộng lớn của mình che chắn phía sau cùng, vừa đẩy vừa dìu cả nhóm tiến từng bước vững chắc qua màn mưa xối xả.

Mười lăm phút vượt bão gian nan dường như dài đằng đẵng như cả thế kỷ.

Nhưng cuối cùng, ánh đèn vàng ấm áp của sảnh khách sạn cũng hiện ra ngay trước mắt.

*RẦM!!*

Cánh cửa ngách khách sạn bật mở. Bác bảo vệ già cùng mẹ Kaien, Kson và Kanata đang đứng chờ trong ruột gan như lửa đốt liền ùa ra kéo cả năm anh em chúng tôi vào bên trong sảnh an toàn.

Cánh cửa gỗ dày cộp lập tức được đóng sập lại, then cài then chốt đóng chặt, ngăn cách hoàn toàn tiếng gầm rú cuồng loạn của bão biển ở phía bên ngoài.

Mẹ Kaien lao tới ôm chầm lấy cả năm đứa con, thở phào nhẹ nhõm: ``Ơn trời lạy Phật\dots'' còn Naomi thì rưng rức nước mắt: ``May quá\dots\ Mọi người về nhà an toàn rồi\dots''

Kson và Towa vội mang khăn bông khô và nước gừng ấm tới quấn quanh người Suzuno, Ryuuhou và Luna. Toàn bộ sảnh khách sạn như vỡ òa trong niềm vui mừng và sự nhẹ nhõm vô bờ bến.

Ryuuhou vừa lau tóc vừa cười toe toét nhìn sang chiếc mũi đang nhét giấy ăn của tôi: ``Công nhận Onii bay trên trời trông ngầu bá cháy bọ chét luôn ấy! Cứ như siêu nhân phiên bản lỗi vậy!''

Tôi nhăn mặt huých nhẹ đầu cô em gái: ``Lỗi cái đầu em ấy! Lần sau xem dự báo thời tiết trước nghe chưa!''

Nhìn Kaiho đang ân cần lau khô tóc cho Suzuno, nhìn Naomi đang lăng xăng chạy lại ôm chầm lấy chân các chị, trong lòng tôi dâng lên một cảm giác ấm áp và bình yên đến lạ kỳ.

Một đêm giông bão kinh hoàng đã thực sự khép lại trước ngưỡng cửa gia đình.

\stpt{-Jour 3-}
\stcap{Dư âm bão tố, hương vị đồng quê và nẻo về Kawachi}

Bảy giờ ba mươi phút sáng ngày thứ ba.

Cơn bão biển dữ dội đêm qua đã tan biến nhanh chóng như cách nó ập đến, để lại bầu trời sớm vùng vịnh Seika trong vắt không một gợn mây. Mặt biển phẳng lặng như gương soi, không khí sau cơn mưa rào trở nên trong lành và mát rượi đến nao lòng.

Thế nhưng, tàn tích mà trận cuồng phong đêm qua để lại cho dải bờ biển Seika thì lại hiện rõ mồn một khắp mọi nẻo đường.

Dọc theo con đường ven biển trước khách sạn, hàng chục cây phi lao và dừa chắn sóng bị gió quật bật gốc, thân cây nằm chắn ngang lối đi bộ. Những chiếc ô cọ che nắng bên bãi cát bị gió bão cuốn phăng rách tả tơi, rác biển và rong rêu dạt trắng xóa cả một dải bờ. Một vài ki-ốt bán nước giải khát lợp mái tôn ven bờ bị gió bốc bay mất phân nửa mái che, các chủ quán đang tất bật thu dọn lại bàn ghế gãy đổ. Nhìn cảnh tượng ngổn ngang ấy, ai nấy đều thầm rùng mình khi nghĩ lại khoảnh khắc kinh hoàng đêm qua.

Trước cổng chính của khách sạn—ngay tại vị trí tôi tiếp đất bằng mặt—toàn bộ ba mươi sáu thành viên của đoàn du lịch công ty mẹ cùng đại gia đình chúng tôi đã tập trung đông đủ để chuẩn bị lên xe quay về Kawachi.

Tôi đứng ở cốp xe bảy chỗ, trán dán một miếng băng gạc y tế to tướng, hai tay xách hai chiếc ba-lô nặng trịch chứa đầy những vỏ ốc, đồ lưu niệm và quần áo của mấy đứa em gái. Dù mẹ tôi đã dặn đi dặn lại là đi nghỉ mát không cần mang vác lỉnh kỉnh, nhưng Ryuuhou vẫn cố nhét thêm cả một bọc đá cuội biển và mấy bộ dây đàn sơ-cua vào hành lý.

Kson khoác ba-lô, tiến lại gần vỗ nhẹ vai tôi, nhếch mép cười ranh mãnh: ``Ê Kuon-paisen, cái đầu với cái mũi của anh vẫn ổn đấy chứ?''

Tôi nhăn mặt sờ nhẹ vào miếng băng gạc: ``Ngoài trừ việc hơi ê ẩm với nhức ra thì anh vẫn sống tốt chán. Ăn nhằm gì so với mấy cái vụ ẩu đả ở văn phòng chứ.''

Kanata đứng cạnh chắp hai tay sau lưng, nghiêm túc gật gù: ``Công nhận là đầu của Kuon-senpai cứng thật đấy\dots''

Ryuuhou đang xách vali kéo ngang qua liền đệm thêm một câu chốt hạ sắc lẹm: ``Cả về nghĩa đen lẫn nghĩa bóng luôn ấy chứ!''

Tôi trừng mắt: ``Này nhé, hai cô nương có ý gì đấy hả!?''

Đứng cạnh đó, mẹ vừa kiểm tra danh sách hành khách vừa bồi thêm một cú: ``Mấy đứa nó nói có sai câu nào đâu con. Đã không biết bay lại còn thích cầm ô làm siêu nhân, không va vào tường thì mới là lạ đấy!''

Tôi thở dài thườn thượt nhìn lên trời. Đến cả mẹ ruột còn quay sang phe đàn em thì tôi còn biết thanh minh với ai được nữa đây?

Đúng lúc đó, Suzuno nhẹ nhàng bước lại gần, hai bàn tay đan vào nhau trước ngực, đôi mắt hai màu dị sắc ngập tràn sự ân cần và xót xa: ``Onii-sama đừng chạm tay vào vết thương nhiều quá kẻo nhiễm trùng ạ\dots\ Sáng nay em đã sát trùng lại và bôi thuốc mỡ thảo dược cho anh rồi, nhưng về đến nhà anh vẫn phải nghỉ ngơi tĩnh dưỡng đấy nhé. Tối qua nếu không có anh và Onee-sama kịp thời lao ra đón, em thật sự không biết ba đứa em sẽ xoay xở thế nào giữa tâm bão nữa\dots''

Bé út Naomi đứng nép cạnh vạt áo Suzuno, hai tay ôm chặt con gấu bông nhỏ, ngước đôi mắt tròn xoe trong veo lên nhìn miếng băng gạc trên trán tôi, cất giọng thỏ thẻ đầy thương cảm: ``Onii-chan có đau lắm không ạ\dots? Tối qua lúc anh bay về trán anh chảy nhiều máu lắm, Naomi sợ muốn khóc luôn\dots\ Lần sau anh đừng bay bằng ô nữa nha anh\dots''

Chloe xách chiếc vali nhỏ lon ton chạy lại, nghe thấy thế liền gãi đầu, đôi mắt đỏ rực thoáng nét hối lỗi: ``Poe\dots\ Tối qua tôi cũng định lao ra hỗ trợ bơi ngược dòng cứu người, ai dè lại bị Ryuuhou đè đầu dúi xuống cát\dots\ Dù sao thì, nhìn ông bay liệng trên trời trông cũng ra dáng người hùng gia đình lắm đấy, đồ Tổng quản lì lợm!''

Towa đứng tựa cửa xe, khẽ mỉm cười dịu dàng: ``Nhưng mà công nhận chuyến đi biển lần này vui thật đấy, senpai.''

Tôi nhìn sang dàn idol Thế hệ thứ Tư, nhìn Chloe đang lon ton xách túi quà lưu niệm cho sếp La+, rồi nhìn Kaiho, Suzuno và Naomi đang ríu rít bên nhau, khóe môi bất giác khẽ cong lên: ``Ừ. Vui đến mức anh suýt nữa thì tăng xông mấy lần, nhưng\dots\ thực sự rất đáng giá.''

Chuyến xe khởi hành lăn bánh rời khỏi bãi biển Seika lúc tám giờ sáng.

\ct

Mười một giờ trưa.

Đoàn xe rẽ vào một thị trấn nhỏ thuộc địa phận Neihei (寧平 - Ninh Bình) - vùng đất non nước hữu tình với những dãy núi đá vôi sừng sững nhấp nhô giữa cánh đồng lúa chín vàng, và cũng chính là quê ngoại của tôi. Cả đoàn dừng chân tại một nhà hàng ẩm thực dân dã nằm nép mình dưới chân núi đá để dùng bữa trưa.

Đặc sản nổi tiếng trứ danh của vùng đất Neihei này không gì khác ngoài các món thịt dê núi.

Vừa bước chân vào gian nhà ăn ngát hương thơm của sả, gừng và tương bần, cô nàng Cừu Watame bỗng khựng lại. Cô nàng hai tay ôm chặt lấy cặp sừng xoắn ốc trên đầu, ánh mắt đầy vẻ cảnh giác nhìn quanh những chiếc đĩa thịt nướng xèo xèo trên bàn rồi cất giọng thanh minh trước: ``Này nhá! Em nói trước cho cả nhà biết, em là linh vật cừu, tuyệt đối không phải là thức ăn đâu đấy nhá!!''

Lời tuyên bố đầy cảnh giác của nàng Cừu làm cả bàn ăn ngơ ngác một giây rồi nổ ra một tràng cười nghiêng ngả.

Luna đang cầm đũa gắp một miếng cơm cháy giòn rụm, tròn xoe mắt bập bẹ: ``Ủa? Đã ai nói zì đâu-nanora?''

Towa khoanh hai tay, thở dài bất lực nhìn cô bạn thân: ``Bà bị chứng hoang tưởng đấy à, Watame? Chúng ta vào đây ăn thịt dê núi, chứ có ai đòi làm thịt cừu đâu mà bà giãy nảy lên thế!?''

Watame xua tay lia lịa, mặt mếu máo: ``Nhưng mà trong phân loại sinh học thì dê với cừu cùng thuộc phân họ Caprinae mà!! Họ hàng gần gũi như thế, lỡ người ta nhìn nhầm rồi bắt em mang xuống bếp thì sao!!''

Mẹ và Ryuuhou ngồi cạnh nghe vậy liền đồng loạt gật gù: ``\dots Ừ thì, xét về mặt sinh học thì con bé nói cũng có lý đấy chứ.''

Watame ôm đầu hoảng hốt: ``Thấy chưa!! Em đã bảo mà!!''

Kanata phì cười vỗ vai cô bạn: ``Tôi nghĩ bà nhạy cảm quá mức rồi đấy con Cừu ngốc ạ!''

Đúng lúc đó, Suzuno nghiêng đầu nhìn Watame, khẽ mỉm cười giải thích bằng giọng điệu học thuật dịu dàng: ``Watame-san đừng lo lắng quá ạ. Thịt dê núi Neihei sở hữu hàm lượng myoglobin cao hơn thịt cừu thông thường, thớ thịt săn chắc và ít mỡ dưới da hơn hẳn do đặc tính leo trèo vách đá vôi hiểm trở. Hơn nữa, mùi hương đặc trưng của acid béo chuỗi nhánh 4-methyloctanoic trong thịt dê nướng ở đây rất khác biệt so với cừu, các đầu bếp chuyên nghiệp ngửi mùi là phân biệt được ngay, không thể nào nhầm lẫn được đâu ạ.''

Nghe bài phân tích hóa sinh cặn kẽ và chuẩn xác của Suzuno, cả bàn ăn đều tròn xoe mắt thán phục, còn Watame thì thở phào nhẹ nhõm như vừa trút được bản án tử hình.

Bé út Naomi ngồi cạnh Suzuno, vừa bốc một miếng cơm cháy chấm nước sốt dê thơm lừng vừa ngây thơ chớp mắt nhìn Watame: ``Thịt dê ngon lắm chị Cừu ơi, không có vị giống chị đâu ạ! Chị ăn thử một miếng đi!''

Chloe thì đã ngồi phồm xoàm ở góc bàn đối diện từ bao giờ. Mặc kệ mọi tranh luận sinh học xung quanh, cô nàng sát thủ cá voi cầm cả một chiếc đùi dê nướng mật ong to tướng cắn ngập răng, vừa nhai ngấu nghiến vừa tấm tắc khen ngợi: ``Ngon quá xá là ngon luôn!! Thịt dê nướng chấm tương bần này đỉnh chóp hơn hẳn mấy món ăn nhanh ở Tokyo, poe!!''

Tôi kéo chiếc ghế bên cạnh Watame ra rồi ngồi xuống, gắp cho cô nàng một miếng thịt dê hấp cuốn lá mơ chấm sốt tương ngọt lịm: ``Yên tâm đi, chẳng ai rảnh mà vác một cô nàng idol trăm ký xuống bếp làm thịt đâu. Nhìn em tròn ủng thế kia người ta lại tưởng thú nhồi bông thì có. Mau ăn đi kẻo nguội ngắt bây giờ.''

Bữa trưa diễn ra trong không khí vô cùng ấm cúng và rôm rả. Thịt dê núi Neihei săn chắc, ngọt mềm, từ món dê tái chanh thanh mát, dê nướng tảng thơm lừng cho đến nồi lẩu dê nghi ngút khói chấm cùng tương bần đậm đà đều khiến mọi người ăn đến no căng bụng.

\ct

Hai giờ chiều.

Trước khi đoàn xe chính thức lăn bánh chặng cuối cùng trở về Kawachi, mẹ tôi dẫn cả đoàn ghé vào một trạm dừng chân mua sắm đặc sản và quà lưu niệm quy mô lớn ven quốc lộ.

Khu mua sắm sáng rực ánh đèn, bày la liệt hàng trăm gian hàng từ đồ thủ công mỹ nghệ, đá phong thủy, chuỗi hạt gỗ cho đến các loại bánh kẹo đóng hộp gắn mác ``Đặc sản truyền thống vùng cao''.

Tôi và Kaiho vừa bước vào cửa hàng, nhìn lướt qua bảng giá niêm yết trên mấy quầy quà tặng lưu niệm thì đồng loạt khựng lại, khóe miệng giật giật không ngừng.

Một chiếc chuông gió bằng vỏ sò nhỏ xíu chỉ to bằng nắm tay mà dán nhãn giá tương đương hai bát phở bò ở thủ đô! Hay một hộp bánh nhãn đóng gói bao bì màu mè sặc sỡ lại có giá đắt gấp ba lần hộp bánh nhãn y hệt đang bày bán ở siêu thị ngay gần ký túc xá Hikawa!

Kaiho khoanh hai tay trước ngực, cầm một chiếc móc khóa hình thuyền buồm bằng gỗ ép lên xem rồi tặc lưỡi lắc đầu ngao ngán: ``Trời đất ơi\dots\ Một miếng gỗ vụn thế này mà chém tận năm chục nghìn!? Cái này ra chợ đầu mối ở Kawachi mua sỉ chắc chưa tới mười nghìn một chiếc! Đúng là chém đẹp khách du lịch mà!''

Tôi cũng buốt ruột gật đầu đồng tình: ``Chuẩn luôn em ạ. Mấy cái danh xưng `đặc sản gia truyền' này thực chất chỉ là chiêu trò marketing đánh vào tâm lý người đi xa mua quà thôi. Toàn hàng nhập từ xưởng gia công ngoại ô chuyển về chứ gia truyền cái nỗi gì.''

Cả hai anh em tôi - những kẻ luôn tôn thờ chủ nghĩa thực dụng và ám ảnh bởi giá trị đồng tiền - chỉ biết đứng nhìn nhau cười trừ, kiên quyết không móc ví cho những món đồ giá trên trời ấy.

Thế nhưng, ở phía bên kia gian hàng, hội chị em lại đang chìm đắm trong cơn sốt mua sắm không lối thoát.

Ryuuhou thì mắt sáng rực trước quầy bán nhạc cụ dân tộc và đồ lưu niệm bằng tre nứa. Cô nàng nâng niu một chiếc sáo trúc khắc hoa văn sóng biển, vừa ướm thử lên môi thổi vài nốt \textit{vi vu} vừa hí hửng: ``Chiếc sáo này âm vang thanh thoát ghê! Lát nữa về phải đem khoe với ban nhạc underground mới được!''

Suzuno thì đứng ở quầy trà thảo mộc và thảo dược khô, cẩn thận dùng tay nâng từng gói trà hoa cúc và nụ tam thất lên ngửi mùi để kiểm tra độ tươi mới. Cô bé tỉ mỉ chọn ra hai hộp trà hảo hạng nhất: ``Hộp trà hoa cúc này tính mát, đem về hãm nước cho bà nội và bác Hijaki uống dưỡng tâm an thần thì tốt lắm đây ạ.''

Naomi thì say mê đứng trước dãy móc khóa hình sinh vật biển đan bằng len thủ công. Cô bé kiễng chân, cẩn thận chọn ra một chú rùa biển tí hon màu xanh ngọc bích, ôm vào lòng với vẻ mặt nâng niu vô cùng hạnh phúc.

Trong khi đó, hội Gen 4 và Chloe cũng nhộn nhịp chẳng kém. Kson hào phóng vung tay mua luôn hai thùng bánh cáy và cơm lam đóng gói để mang về chia cho ban quản lý COVER; Kanata và Towa thì chụm đầu chọn mấy chiếc móc chìa khóa đôi hình thiên sứ và ác ma; Watame ôm trọn năm hộp bánh nhãn mật ong ngọt lịm; còn Chloe thì hớn hở xách hai túi bánh mực khô to đùng, miệng lẩm bẩm tính toán: ``Một túi nộp cho sếp La+ làm bằng chứng công tác, còn một túi giấu dưới gầm giường để ăn đêm dần!''

Sau gần một tiếng đồng hồ nhộn nhịp chọn lựa và trả giá, mọi người ai nấy đều tay xách nách mang những túi quà lưu niệm rực rỡ sắc màu, vui vẻ bước lên xe tiếp tục hành trình.

\ct

Sáu giờ ba mươi phút chiều.

Bóng hoàng hôn vàng cam buông dài trên những con phố quen thuộc của thành phố Kawachi. Chiếc xe bảy chỗ của gia đình chúng tôi chầm chậm rẽ vào con ngõ nhỏ và đỗ xịch trước cửa ký túc xá Hikawa.

Chuyến hành trình ba ngày hai đêm đầy ắp những con sóng, những pha tai nạn thót tim và vô vàn tiếng cười giòn giã đã chính thức khép lại.

Bố tooi mặc bộ quần áo thợ cơ khí lấm lem dầu mỡ cùng Delutaya đang đứng trước hiên nhà đón chúng tôi. Nàng Tam giác nở nụ cười dịu dàng như hoa nở mùa xuân, đỡ lấy túi hành lý từ tay Suzuno: ``Mọi người về rồi! Chuyến đi biển thế nào ạ? Mike-chan với bà nội ở nhà nhớ mọi người lắm đấy!''

Bé út Naomi lập tức chạy ùa tới ôm chầm lấy chân bố Ryujin, giơ chiếc vòng tay xâu bằng vỏ sò khoe rối rít: ``Bố ơi! Biển lớn lắm bố ạ! Chị Suzuno dạy con làm vòng tay đẹp chưa này!''

Tôi tựa lưng vào mép cốp xe, nhìn sang hội chị em Gen 4 đang chuẩn bị bắt taxi về lại ký túc xá của họ, mỉm cười hỏi: ``Thế nào các cô gái? Kỳ nghỉ ở bãi biển Kaiwa có đủ để xả stress chưa?''

Luna giơ hai tay lên cao, reo to: ``Vui quá xá là vui-nanora!! Cảm ơn Kuon-senpai và cả nhà nhìu nhìu nanora!''

Kson cười khà khà, vỗ mạnh vào vai Kanata: ``Tuyệt vời ông mặt trời luôn! Lâu lắm rồi bọn tôi mới có dịp tụ tập đông đủ và quậy tưng bừng như thế này!''

Kanata hai mắt sáng ngời: ``Được đi chơi xa cùng Kson và cả nhà Hikawa là kỷ niệm tuyệt vời nhất mùa hè này của em đấy ạ!''

Towa khoanh tay làm bộ phụng phịu: ``Này, còn tôi với Watame thì bị coi như không khí đấy à!?''

Watame cười khúc khích, ôm lấy cánh tay Towa: ``Đâu có đâu nà\~{} Cứ như thể tình bạn của Thế hệ thứ Tư tụi mình lại càng bền chặt hơn sau bao nhiêu sóng gió ấy nhỉ!''

Chloe đứng nép cạnh cửa xe, ôm chặt bọc quà lưu niệm bánh mực khô, đôi mắt đỏ rực nhìn đại gia đình chúng tôi bằng một ánh nhìn ngập tràn sự lưu luyến. Cô nàng khẽ cúi đầu chào mẹ Kaien rồi quay sang tôi, cất giọng lí nhí: ``Này Kuon-kun\dots\ Cảm ơn ông nhé. Chuyến đi này\dots\ thực sự rất có ý nghĩa với tôi.''

Tôi khẽ xoa đầu cô bạn thuở nhỏ: ``Về văn phòng holoX nhớ báo cáo tốt với sếp La+ đấy nhé, đồ sát thủ tập sự.''

Chloe đỏ mặt phồng má: ``Biết rồi mà, poe poe!''

Chuyến đi biển lần này thực sự là một trải nghiệm vừa quen thuộc lại vừa mới mẻ. Quen thuộc, bởi nó mang lại cảm giác ấm áp, bình dị như những ngày đầu tôi chập chững bước chân vào thế giới của các cô gái ấy; nhưng mới mẻ, bởi đây là lần đầu tiên đại gia đình Hikawa hoàn chỉnh - với sự hiện diện của Kaiho, Suzuno và Naomi - cùng nhau chia sẻ từng con sóng, từng bát cơm và từng khoảnh khắc sinh tử giữa giông bão cuộc đời.

Tựa lưng vào cánh cổng sắt của ngôi nhà thân thương, nhìn ánh đèn ấm áp hắt ra từ khung cửa sổ tầng sáu, tôi khẽ mỉm cười thầm tự nhủ: lần tới, nhất định tôi sẽ kéo bằng được cả bố Ryujin, Delutaya và tất cả mọi người cùng xách ba-lô lên đường.

Bởi vì chỉ khi có đầy đủ tất cả các thành viên bên cạnh, gia đình mới thực sự là một mái ấm trọn vẹn nhất.

\end{document}
